% Supplementary Materials for:
% Fixed-Width vs. Multiprecision Digit Generation for JSON Canonicalization:
% A Four-Algorithm Empirical Study
%
% Mark Lenhardt, Lattice Substrate
\documentclass[10pt]{article}

\usepackage[margin=0.75in]{geometry}
\usepackage{booktabs}
\usepackage{longtable}
\usepackage{pdflscape}
\usepackage{xspace}
\usepackage{amsmath}
\usepackage{hyperref}
\usepackage{caption}
\usepackage{fancyhdr}

\newcommand{\schubfach}{Schubfach\xspace}
\newcommand{\burgerdybvig}{Burger--Dybvig\xspace}
\newcommand{\ryualg}{Ry\=u\xspace}
\newcommand{\dragonbox}{Dragonbox\xspace}
\newcommand{\ecma}{ECMA-262\xspace}
\newcommand{\rfcjcs}{RFC~8785\xspace}

% Provide labels referenced by auto-generated tables that are defined in the main paper
\newcommand{\fakesec}[1]{\label{#1}}
\AtBeginDocument{%
  \fakesec{sec:discussion}%
  \fakesec{sec:appendix-tables}%
  \fakesec{sec:methodology}%
}

\pagestyle{fancy}
\fancyhf{}
\fancyhead[L]{\small Supplementary Materials --- Lenhardt (2026)}
\fancyhead[R]{\small \thepage}
\renewcommand{\headrulewidth}{0.4pt}

\begin{document}

\begin{center}
{\Large\bfseries Supplementary Materials}\\[6pt]
{\large Fixed-Width vs.\ Multiprecision Digit Generation\\
for JSON Canonicalization: A Four-Algorithm Empirical Study}\\[6pt]
Mark Lenhardt\\
Lattice Substrate\\[6pt]
\today
\end{center}

\vspace{1em}
\noindent
This document contains the complete benchmark tables and statistical
comparisons referenced as ``supplementary materials'' in the main
paper.  All data are auto-generated from raw benchmark output by
\texttt{scripts/generate-paper-tables.go} and are reproducible from
the archived artifact (DOI:
\href{https://doi.org/10.5281/zenodo.18890835}{10.5281/zenodo.18890835}).

\tableofcontents

\clearpage

% ─── Section 1: Full API Benchmarks ───────────────────────────────────────────

\section{Full API Benchmark Results}
\label{sec:supp-api}

Table~\ref{tab:api-benchmarks} reports the complete API benchmark
results for all 16~comparison pairs across all workloads and modes
(canonicalize and verify).  Each row reports the mean of $n{=}10$
iterations.  Allocation columns ($\Delta$allocs/op, $\Delta$B/op)
report the difference between the two implementations in each pair;
negative values indicate the first implementation allocates less.

For Rust implementations, allocation data is measured via a counting
global allocator (\texttt{std::alloc::GlobalAlloc}) that intercepts
all heap allocations.

\begin{landscape}
% Auto-generated by scripts/generate-paper-tables.go
\scriptsize
\begin{longtable}{lllrrrrr}
\caption{API benchmark results (mean of $n=10$ iterations).}
\label{tab:api-benchmarks}\\
\toprule
Pair & Workload & Mode & \multicolumn{2}{c}{ns/op} & Speedup & $\Delta$allocs/op & $\Delta$B/op \\
\midrule
\endfirsthead
\toprule
Pair & Workload & Mode & \multicolumn{2}{c}{ns/op} & Speedup & $\Delta$allocs/op & $\Delta$B/op \\
\midrule
\endhead
\midrule
\multicolumn{8}{r}{\textit{Continued on next page}} \\
\endfoot
\bottomrule
\multicolumn{8}{l}{\footnotesize B/op = bytes allocated per operation. allocs/op = heap allocations per operation.} \\
\multicolumn{8}{l}{\footnotesize $\Delta$ columns show Impl B $-$ Impl A for each pair.} \\
\endlastfoot
go:schubfach-vs-bd & array-2048 & canonicalize & 2538189 & 4058559 & 1.60$\times$ & +8210 & +186691 \\
go:schubfach-vs-bd & array-256 & canonicalize & 304201 & 472012 & 1.55$\times$ & +1020 & +24976 \\
go:schubfach-vs-bd & canonical-minimal & canonicalize & 351 & 530 & 1.51$\times$ & +0 & +0 \\
go:schubfach-vs-bd & control-escapes & canonicalize & 1333 & 1252 & 1.06$\times$ & +0 & +0 \\
go:schubfach-vs-bd & deep & canonicalize & 36663 & 36648 & 1.00$\times$ & +0 & +0 \\
go:schubfach-vs-bd & deep-64 & canonicalize & 16218 & 16221 & 1.00$\times$ & +0 & +0 \\
go:schubfach-vs-bd & escaped-key-order & canonicalize & 706 & 1125 & 1.59$\times$ & +2 & +41 \\
go:schubfach-vs-bd & large & canonicalize & 11952645 & 20608504 & 1.72$\times$ & +32792 & +663257 \\
go:schubfach-vs-bd & long-string & canonicalize & 75536 & 75430 & 1.00$\times$ & +0 & +0 \\
go:schubfach-vs-bd & medium & canonicalize & 341381 & 579627 & 1.70$\times$ & +1014 & +21130 \\
go:schubfach-vs-bd & mixed-prod & canonicalize & 135688 & 179064 & 1.32$\times$ & +255 & +6237 \\
go:schubfach-vs-bd & nested-mixed & canonicalize & 2177 & 3065 & 1.41$\times$ & +4 & +85 \\
go:schubfach-vs-bd & number-heavy & canonicalize & 13812 & 21307 & 1.54$\times$ & +20 & +830 \\
go:schubfach-vs-bd & numeric-boundary & canonicalize & 12603 & 18700 & 1.48$\times$ & +14 & +706 \\
go:schubfach-vs-bd & rfc-key-sorting & canonicalize & 2999 & 2963 & 1.01$\times$ & +0 & +0 \\
go:schubfach-vs-bd & small & canonicalize & 1235 & 1685 & 1.36$\times$ & +2 & +50 \\
go:schubfach-vs-bd & surrogate-pair & canonicalize & 646 & 638 & 1.01$\times$ & +0 & +0 \\
go:schubfach-vs-bd & unicode & canonicalize & 1546 & 1541 & 1.00$\times$ & +0 & +0 \\
go:schubfach-vs-bd & verify-whitespace & canonicalize & 906 & 1559 & 1.72$\times$ & +4 & +82 \\
go:schubfach-vs-bd & canonical/array-2048 & verify & 2475365 & 4019281 & 1.62$\times$ & +8208 & +186522 \\
go:schubfach-vs-bd & canonical/array-256 & verify & 314132 & 483772 & 1.54$\times$ & +1021 & +25208 \\
go:schubfach-vs-bd & canonical/canonical-minimal & verify & 359 & 548 & 1.53$\times$ & +0 & +0 \\
go:schubfach-vs-bd & canonical/control-escapes & verify & 1319 & 1277 & 1.03$\times$ & +0 & +0 \\
go:schubfach-vs-bd & canonical/deep & verify & 37929 & 38260 & 1.01$\times$ & +0 & +0 \\
go:schubfach-vs-bd & canonical/deep-64 & verify & 16664 & 16930 & 1.02$\times$ & +0 & +0 \\
go:schubfach-vs-bd & canonical/escaped-key-order & verify & 685 & 1134 & 1.65$\times$ & +2 & +41 \\
go:schubfach-vs-bd & canonical/large & verify & 11465484 & 19941996 & 1.74$\times$ & +32790 & +662980 \\
go:schubfach-vs-bd & canonical/long-string & verify & 76827 & 77335 & 1.01$\times$ & +0 & +0 \\
go:schubfach-vs-bd & canonical/medium & verify & 338837 & 576583 & 1.70$\times$ & +1015 & +21370 \\
go:schubfach-vs-bd & canonical/mixed-prod & verify & 132472 & 175149 & 1.32$\times$ & +255 & +6323 \\
go:schubfach-vs-bd & canonical/nested-mixed & verify & 2203 & 3098 & 1.41$\times$ & +4 & +86 \\
go:schubfach-vs-bd & canonical/number-heavy & verify & 13838 & 21562 & 1.56$\times$ & +20 & +834 \\
go:schubfach-vs-bd & canonical/numeric-boundary & verify & 12820 & 18752 & 1.46$\times$ & +14 & +709 \\
go:schubfach-vs-bd & canonical/rfc-key-sorting & verify & 2840 & 2831 & 1.00$\times$ & +0 & +0 \\
go:schubfach-vs-bd & canonical/small & verify & 1214 & 1651 & 1.36$\times$ & +2 & +51 \\
go:schubfach-vs-bd & canonical/surrogate-pair & verify & 651 & 653 & 1.00$\times$ & +0 & +0 \\
go:schubfach-vs-bd & canonical/unicode & verify & 1540 & 1544 & 1.00$\times$ & +0 & +0 \\
go:schubfach-vs-bd & canonical/verify-whitespace & verify & 929 & 1588 & 1.71$\times$ & +4 & +82 \\
go:schubfach-vs-bd & noncanonical/control-escapes & verify & 1318 & 1257 & 1.05$\times$ & +0 & +0 \\
go:schubfach-vs-bd & noncanonical/escaped-key-order & verify & 701 & 1128 & 1.61$\times$ & +2 & +41 \\
go:schubfach-vs-bd & noncanonical/large & verify & 11919728 & 20550017 & 1.72$\times$ & +32792 & +663497 \\
go:schubfach-vs-bd & noncanonical/medium & verify & 343792 & 574323 & 1.67$\times$ & +1014 & +21131 \\
go:schubfach-vs-bd & noncanonical/mixed-prod & verify & 137067 & 177569 & 1.30$\times$ & +255 & +6226 \\
go:schubfach-vs-bd & noncanonical/nested-mixed & verify & 2191 & 3070 & 1.40$\times$ & +4 & +85 \\
go:schubfach-vs-bd & noncanonical/number-heavy & verify & 13881 & 21352 & 1.54$\times$ & +20 & +830 \\
go:schubfach-vs-bd & noncanonical/numeric-boundary & verify & 12851 & 18650 & 1.45$\times$ & +14 & +706 \\
go:schubfach-vs-bd & noncanonical/rfc-key-sorting & verify & 3015 & 2959 & 1.02$\times$ & +0 & +0 \\
go:schubfach-vs-bd & noncanonical/small & verify & 1250 & 1668 & 1.33$\times$ & +2 & +50 \\
go:schubfach-vs-bd & noncanonical/surrogate-pair & verify & 656 & 643 & 1.02$\times$ & +0 & +0 \\
go:schubfach-vs-bd & noncanonical/unicode & verify & 1566 & 1547 & 1.01$\times$ & +0 & +0 \\
go:schubfach-vs-bd & noncanonical/verify-whitespace & verify & 922 & 1554 & 1.69$\times$ & +4 & +82 \\
rs:schubfach-vs-bd & array-2048 & canonicalize & 1261394 & 2891090 & 2.29$\times$ & +0 & +0 \\
rs:schubfach-vs-bd & array-256 & canonicalize & 147811 & 332147 & 2.25$\times$ & +0 & +0 \\
rs:schubfach-vs-bd & canonical-minimal & canonicalize & 183 & 418 & 2.28$\times$ & +0 & +0 \\
rs:schubfach-vs-bd & control-escapes & canonicalize & 620 & 617 & 1.01$\times$ & +0 & +0 \\
rs:schubfach-vs-bd & deep & canonicalize & 26488 & 25683 & 1.03$\times$ & +0 & +0 \\
rs:schubfach-vs-bd & deep-64 & canonicalize & 9482 & 9174 & 1.03$\times$ & +0 & +0 \\
rs:schubfach-vs-bd & escaped-key-order & canonicalize & 303 & 742 & 2.45$\times$ & +0 & +0 \\
rs:schubfach-vs-bd & large & canonicalize & 5132778 & 16125242 & 3.14$\times$ & +0 & +0 \\
rs:schubfach-vs-bd & long-string & canonicalize & 17188 & 17328 & 1.01$\times$ & +0 & +0 \\
rs:schubfach-vs-bd & medium & canonicalize & 166916 & 435894 & 2.61$\times$ & +0 & +0 \\
rs:schubfach-vs-bd & mixed-prod & canonicalize & 65576 & 107081 & 1.63$\times$ & +0 & +0 \\
rs:schubfach-vs-bd & nested-mixed & canonicalize & 878 & 1698 & 1.93$\times$ & +0 & +0 \\
rs:schubfach-vs-bd & number-heavy & canonicalize & 1403 & 11464 & 8.17$\times$ & +0 & +0 \\
rs:schubfach-vs-bd & numeric-boundary & canonicalize & 955 & 8774 & 9.18$\times$ & +0 & +0 \\
rs:schubfach-vs-bd & rfc-key-sorting & canonicalize & 1278 & 1263 & 1.01$\times$ & +0 & +0 \\
rs:schubfach-vs-bd & small & canonicalize & 578 & 1026 & 1.78$\times$ & +0 & +0 \\
rs:schubfach-vs-bd & surrogate-pair & canonicalize & 273 & 276 & 1.01$\times$ & +0 & +0 \\
rs:schubfach-vs-bd & unicode & canonicalize & 1686 & 1650 & 1.02$\times$ & +0 & +0 \\
rs:schubfach-vs-bd & verify-whitespace & canonicalize & 369 & 983 & 2.66$\times$ & +0 & +0 \\
rs:schubfach-vs-bd & canonical/array-2048 & verify & 1178769 & 2920389 & 2.48$\times$ & +0 & +0 \\
rs:schubfach-vs-bd & canonical/array-256 & verify & 132345 & 316436 & 2.39$\times$ & +0 & +0 \\
rs:schubfach-vs-bd & canonical/canonical-minimal & verify & 181 & 421 & 2.32$\times$ & +0 & +0 \\
rs:schubfach-vs-bd & canonical/control-escapes & verify & 586 & 608 & 1.04$\times$ & +0 & +0 \\
rs:schubfach-vs-bd & canonical/deep & verify & 26824 & 25174 & 1.07$\times$ & +0 & +0 \\
rs:schubfach-vs-bd & canonical/deep-64 & verify & 9156 & 9266 & 1.01$\times$ & +0 & +0 \\
rs:schubfach-vs-bd & canonical/escaped-key-order & verify & 291 & 712 & 2.45$\times$ & +0 & +0 \\
rs:schubfach-vs-bd & canonical/large & verify & 5375262 & 16330345 & 3.04$\times$ & +0 & +0 \\
rs:schubfach-vs-bd & canonical/long-string & verify & 17299 & 17580 & 1.02$\times$ & +0 & +0 \\
rs:schubfach-vs-bd & canonical/medium & verify & 143326 & 437556 & 3.05$\times$ & +0 & +0 \\
rs:schubfach-vs-bd & canonical/mixed-prod & verify & 65487 & 109623 & 1.67$\times$ & +0 & +0 \\
rs:schubfach-vs-bd & canonical/nested-mixed & verify & 856 & 1666 & 1.95$\times$ & +0 & +0 \\
rs:schubfach-vs-bd & canonical/number-heavy & verify & 1388 & 11075 & 7.98$\times$ & +0 & +0 \\
rs:schubfach-vs-bd & canonical/numeric-boundary & verify & 961 & 8726 & 9.08$\times$ & +0 & +0 \\
rs:schubfach-vs-bd & canonical/rfc-key-sorting & verify & 1322 & 1360 & 1.03$\times$ & +0 & +0 \\
rs:schubfach-vs-bd & canonical/small & verify & 579 & 1016 & 1.75$\times$ & +0 & +0 \\
rs:schubfach-vs-bd & canonical/surrogate-pair & verify & 284 & 291 & 1.02$\times$ & +0 & +0 \\
rs:schubfach-vs-bd & canonical/unicode & verify & 1725 & 1648 & 1.05$\times$ & +0 & +0 \\
rs:schubfach-vs-bd & canonical/verify-whitespace & verify & 353 & 967 & 2.74$\times$ & +0 & +0 \\
schubfach:go-vs-rs & array-2048 & canonicalize & 2538189 & 1261394 & 2.01$\times$ & -49134 & -4939278 \\
schubfach:go-vs-rs & array-256 & canonicalize & 304201 & 147811 & 2.06$\times$ & -6131 & -604880 \\
schubfach:go-vs-rs & canonical-minimal & canonicalize & 351 & 183 & 1.92$\times$ & -8 & -504 \\
schubfach:go-vs-rs & control-escapes & canonicalize & 1333 & 620 & 2.15$\times$ & -21 & -2064 \\
schubfach:go-vs-rs & deep & canonicalize & 36663 & 26488 & 1.38$\times$ & -843 & -58960 \\
schubfach:go-vs-rs & deep-64 & canonicalize & 16218 & 9482 & 1.71$\times$ & -387 & -26904 \\
schubfach:go-vs-rs & escaped-key-order & canonicalize & 706 & 303 & 2.33$\times$ & -14 & -1168 \\
schubfach:go-vs-rs & large & canonicalize & 11952645 & 5132778 & 2.33$\times$ & -204776 & -20914304 \\
schubfach:go-vs-rs & long-string & canonicalize & 75536 & 17188 & 4.39$\times$ & -14 & -35948 \\
schubfach:go-vs-rs & medium & canonicalize & 341381 & 166916 & 2.05$\times$ & -6380 & -618687 \\
schubfach:go-vs-rs & mixed-prod & canonicalize & 135688 & 65576 & 2.07$\times$ & -2340 & -256989 \\
schubfach:go-vs-rs & nested-mixed & canonicalize & 2177 & 878 & 2.48$\times$ & -45 & -4128 \\
schubfach:go-vs-rs & number-heavy & canonicalize & 13812 & 1403 & 9.84$\times$ & -60 & -4576 \\
schubfach:go-vs-rs & numeric-boundary & canonicalize & 12603 & 955 & 13.19$\times$ & -38 & -3208 \\
schubfach:go-vs-rs & rfc-key-sorting & canonicalize & 2999 & 1278 & 2.35$\times$ & -47 & -4256 \\
schubfach:go-vs-rs & small & canonicalize & 1235 & 578 & 2.14$\times$ & -22 & -2152 \\
schubfach:go-vs-rs & surrogate-pair & canonicalize & 646 & 273 & 2.37$\times$ & -14 & -1168 \\
schubfach:go-vs-rs & unicode & canonicalize & 1546 & 1686 & 1.09$\times$ & -19 & -2019 \\
schubfach:go-vs-rs & verify-whitespace & canonicalize & 906 & 369 & 2.45$\times$ & -19 & -1576 \\
schubfach:go-vs-rs & canonical/array-2048 & verify & 2475365 & 1178769 & 2.10$\times$ & -49134 & -4939271 \\
schubfach:go-vs-rs & canonical/array-256 & verify & 314132 & 132345 & 2.37$\times$ & -6131 & -604881 \\
schubfach:go-vs-rs & canonical/canonical-minimal & verify & 359 & 181 & 1.98$\times$ & -8 & -504 \\
schubfach:go-vs-rs & canonical/control-escapes & verify & 1319 & 586 & 2.25$\times$ & -21 & -2064 \\
schubfach:go-vs-rs & canonical/deep & verify & 37929 & 26824 & 1.41$\times$ & -843 & -58960 \\
schubfach:go-vs-rs & canonical/deep-64 & verify & 16664 & 9156 & 1.82$\times$ & -387 & -26904 \\
schubfach:go-vs-rs & canonical/escaped-key-order & verify & 685 & 291 & 2.36$\times$ & -14 & -1160 \\
schubfach:go-vs-rs & canonical/large & verify & 11465484 & 5375262 & 2.13$\times$ & -204776 & -20914288 \\
schubfach:go-vs-rs & canonical/long-string & verify & 76827 & 17299 & 4.44$\times$ & -14 & -35948 \\
schubfach:go-vs-rs & canonical/medium & verify & 338837 & 143326 & 2.36$\times$ & -6380 & -618688 \\
schubfach:go-vs-rs & canonical/mixed-prod & verify & 132472 & 65487 & 2.02$\times$ & -2340 & -256989 \\
schubfach:go-vs-rs & canonical/nested-mixed & verify & 2203 & 856 & 2.58$\times$ & -45 & -4128 \\
schubfach:go-vs-rs & canonical/number-heavy & verify & 13838 & 1388 & 9.97$\times$ & -60 & -4576 \\
schubfach:go-vs-rs & canonical/numeric-boundary & verify & 12820 & 961 & 13.34$\times$ & -38 & -3208 \\
schubfach:go-vs-rs & canonical/rfc-key-sorting & verify & 2840 & 1322 & 2.15$\times$ & -42 & -4240 \\
schubfach:go-vs-rs & canonical/small & verify & 1214 & 579 & 2.10$\times$ & -22 & -2152 \\
schubfach:go-vs-rs & canonical/surrogate-pair & verify & 651 & 284 & 2.29$\times$ & -14 & -1160 \\
schubfach:go-vs-rs & canonical/unicode & verify & 1540 & 1725 & 1.12$\times$ & -19 & -2003 \\
schubfach:go-vs-rs & canonical/verify-whitespace & verify & 929 & 353 & 2.63$\times$ & -19 & -1576 \\
bd:go-vs-rs & array-2048 & canonicalize & 4058559 & 2891090 & 1.40$\times$ & -57344 & -5125968 \\
bd:go-vs-rs & array-256 & canonicalize & 472012 & 332147 & 1.42$\times$ & -7151 & -629856 \\
bd:go-vs-rs & canonical-minimal & canonicalize & 530 & 418 & 1.27$\times$ & -8 & -504 \\
bd:go-vs-rs & control-escapes & canonicalize & 1252 & 617 & 2.03$\times$ & -21 & -2064 \\
bd:go-vs-rs & deep & canonicalize & 36648 & 25683 & 1.43$\times$ & -843 & -58960 \\
bd:go-vs-rs & deep-64 & canonicalize & 16221 & 9174 & 1.77$\times$ & -387 & -26904 \\
bd:go-vs-rs & escaped-key-order & canonicalize & 1125 & 742 & 1.52$\times$ & -16 & -1209 \\
bd:go-vs-rs & large & canonicalize & 20608504 & 16125242 & 1.28$\times$ & -237568 & -21577561 \\
bd:go-vs-rs & long-string & canonicalize & 75430 & 17328 & 4.35$\times$ & -14 & -35948 \\
bd:go-vs-rs & medium & canonicalize & 579627 & 435894 & 1.33$\times$ & -7394 & -639818 \\
bd:go-vs-rs & mixed-prod & canonicalize & 179064 & 107081 & 1.67$\times$ & -2595 & -263226 \\
bd:go-vs-rs & nested-mixed & canonicalize & 3065 & 1698 & 1.81$\times$ & -49 & -4213 \\
bd:go-vs-rs & number-heavy & canonicalize & 21307 & 11464 & 1.86$\times$ & -80 & -5406 \\
bd:go-vs-rs & numeric-boundary & canonicalize & 18700 & 8774 & 2.13$\times$ & -52 & -3914 \\
bd:go-vs-rs & rfc-key-sorting & canonicalize & 2963 & 1263 & 2.35$\times$ & -47 & -4256 \\
bd:go-vs-rs & small & canonicalize & 1685 & 1026 & 1.64$\times$ & -24 & -2202 \\
bd:go-vs-rs & surrogate-pair & canonicalize & 638 & 276 & 2.31$\times$ & -14 & -1168 \\
bd:go-vs-rs & unicode & canonicalize & 1541 & 1650 & 1.07$\times$ & -19 & -2019 \\
bd:go-vs-rs & verify-whitespace & canonicalize & 1559 & 983 & 1.59$\times$ & -23 & -1658 \\
bd:go-vs-rs & canonical/array-2048 & verify & 4019281 & 2920389 & 1.38$\times$ & -57342 & -5125793 \\
bd:go-vs-rs & canonical/array-256 & verify & 483772 & 316436 & 1.53$\times$ & -7152 & -630089 \\
bd:go-vs-rs & canonical/canonical-minimal & verify & 548 & 421 & 1.30$\times$ & -8 & -504 \\
bd:go-vs-rs & canonical/control-escapes & verify & 1277 & 608 & 2.10$\times$ & -21 & -2064 \\
bd:go-vs-rs & canonical/deep & verify & 38260 & 25174 & 1.52$\times$ & -843 & -58960 \\
bd:go-vs-rs & canonical/deep-64 & verify & 16930 & 9266 & 1.83$\times$ & -387 & -26904 \\
bd:go-vs-rs & canonical/escaped-key-order & verify & 1134 & 712 & 1.59$\times$ & -16 & -1201 \\
bd:go-vs-rs & canonical/large & verify & 19941996 & 16330345 & 1.22$\times$ & -237565 & -21577268 \\
bd:go-vs-rs & canonical/long-string & verify & 77335 & 17580 & 4.40$\times$ & -14 & -35948 \\
bd:go-vs-rs & canonical/medium & verify & 576583 & 437556 & 1.32$\times$ & -7395 & -640059 \\
bd:go-vs-rs & canonical/mixed-prod & verify & 175149 & 109623 & 1.60$\times$ & -2595 & -263312 \\
bd:go-vs-rs & canonical/nested-mixed & verify & 3098 & 1666 & 1.86$\times$ & -49 & -4214 \\
bd:go-vs-rs & canonical/number-heavy & verify & 21562 & 11075 & 1.95$\times$ & -80 & -5410 \\
bd:go-vs-rs & canonical/numeric-boundary & verify & 18752 & 8726 & 2.15$\times$ & -52 & -3917 \\
bd:go-vs-rs & canonical/rfc-key-sorting & verify & 2831 & 1360 & 2.08$\times$ & -42 & -4240 \\
bd:go-vs-rs & canonical/small & verify & 1651 & 1016 & 1.63$\times$ & -24 & -2203 \\
bd:go-vs-rs & canonical/surrogate-pair & verify & 653 & 291 & 2.24$\times$ & -14 & -1160 \\
bd:go-vs-rs & canonical/unicode & verify & 1544 & 1648 & 1.07$\times$ & -19 & -2003 \\
bd:go-vs-rs & canonical/verify-whitespace & verify & 1588 & 967 & 1.64$\times$ & -23 & -1658 \\
\end{longtable}
\normalsize

\end{landscape}

\clearpage

% ─── Section 2: Full CLI Benchmarks ──────────────────────────────────────────

\section{Full CLI Benchmark Results}
\label{sec:supp-cli}

Table~\ref{tab:cli-benchmarks} reports the complete CLI benchmark
results for all 16~comparison pairs across all workloads and modes.
Each row reports the mean wall time of $n{=}15$ repetitions (3~warmup
discarded).  Task order was randomized with a fixed seed.

\begin{landscape}
% Auto-generated by scripts/generate-paper-tables.go
\scriptsize
\begin{longtable}{lllrrrr}
\caption{CLI benchmark results (end-to-end, $n=15$).}
\label{tab:cli-benchmarks}\\
\toprule
Pair & Workload & Mode & Impl A (ms) & Impl B (ms) & Speedup & Winner \\
\midrule
\endfirsthead
\toprule
Pair & Workload & Mode & Impl A (ms) & Impl B (ms) & Speedup & Winner \\
\midrule
\endhead
\midrule
\multicolumn{7}{r}{\textit{Continued on next page}} \\
\endfoot
\bottomrule
\multicolumn{7}{l}{\footnotesize Speedup = slower/faster mean wall time. $n=15$ repetitions, 3 warmup runs discarded.} \\
\endlastfoot
go:schubfach-vs-bd & array-2048 & canonicalize & 5.642 & 7.202 & 1.28$\times$ & schubfach \\
go:schubfach-vs-bd & array-2048 & verify & 5.694 & 7.125 & 1.25$\times$ & schubfach \\
go:schubfach-vs-bd & array-256 & canonicalize & 1.840 & 2.383 & 1.30$\times$ & schubfach \\
go:schubfach-vs-bd & array-256 & verify & 2.087 & 2.246 & 1.08$\times$ & schubfach \\
go:schubfach-vs-bd & canonical-minimal & canonicalize & 1.378 & 1.655 & 1.20$\times$ & schubfach \\
go:schubfach-vs-bd & canonical-minimal & verify & 1.443 & 1.654 & 1.15$\times$ & schubfach \\
go:schubfach-vs-bd & control-escapes & canonicalize & 1.501 & 1.569 & 1.05$\times$ & schubfach \\
go:schubfach-vs-bd & control-escapes & verify & 1.352 & 1.697 & 1.26$\times$ & schubfach \\
go:schubfach-vs-bd & deep & canonicalize & 1.643 & 1.722 & 1.05$\times$ & schubfach \\
go:schubfach-vs-bd & deep & verify & 1.704 & 1.723 & 1.01$\times$ & schubfach \\
go:schubfach-vs-bd & deep-64 & canonicalize & 1.443 & 1.828 & 1.27$\times$ & schubfach \\
go:schubfach-vs-bd & deep-64 & verify & 1.584 & 1.636 & 1.03$\times$ & schubfach \\
go:schubfach-vs-bd & escaped-key-order & canonicalize & 1.489 & 1.553 & 1.04$\times$ & schubfach \\
go:schubfach-vs-bd & escaped-key-order & verify & 1.543 & 1.546 & 1.00$\times$ & schubfach \\
go:schubfach-vs-bd & large & canonicalize & 16.530 & 24.906 & 1.51$\times$ & schubfach \\
go:schubfach-vs-bd & large & verify & 15.913 & 24.895 & 1.56$\times$ & schubfach \\
go:schubfach-vs-bd & long-string & canonicalize & 1.566 & 1.821 & 1.16$\times$ & schubfach \\
go:schubfach-vs-bd & long-string & verify & 1.791 & 1.912 & 1.07$\times$ & schubfach \\
go:schubfach-vs-bd & medium & canonicalize & 1.986 & 2.659 & 1.34$\times$ & schubfach \\
go:schubfach-vs-bd & medium & verify & 2.143 & 2.531 & 1.18$\times$ & schubfach \\
go:schubfach-vs-bd & mixed-prod & canonicalize & 1.851 & 1.907 & 1.03$\times$ & schubfach \\
go:schubfach-vs-bd & mixed-prod & verify & 1.614 & 1.877 & 1.16$\times$ & schubfach \\
go:schubfach-vs-bd & nested-mixed & canonicalize & 1.481 & 1.613 & 1.09$\times$ & schubfach \\
go:schubfach-vs-bd & nested-mixed & verify & 1.486 & 1.739 & 1.17$\times$ & schubfach \\
go:schubfach-vs-bd & number-heavy & canonicalize & 1.484 & 1.656 & 1.12$\times$ & schubfach \\
go:schubfach-vs-bd & number-heavy & verify & 1.523 & 1.763 & 1.16$\times$ & schubfach \\
go:schubfach-vs-bd & numeric-boundary & canonicalize & 1.578 & 1.556 & 1.01$\times$ & json-canon \\
go:schubfach-vs-bd & numeric-boundary & verify & 1.550 & 1.627 & 1.05$\times$ & schubfach \\
go:schubfach-vs-bd & rfc-key-sorting & canonicalize & 1.442 & 1.564 & 1.08$\times$ & schubfach \\
go:schubfach-vs-bd & rfc-key-sorting & verify & 1.499 & 1.566 & 1.04$\times$ & schubfach \\
go:schubfach-vs-bd & small & canonicalize & 1.453 & 1.588 & 1.09$\times$ & schubfach \\
go:schubfach-vs-bd & small & verify & 1.561 & 1.653 & 1.06$\times$ & schubfach \\
go:schubfach-vs-bd & surrogate-pair & canonicalize & 1.425 & 1.549 & 1.09$\times$ & schubfach \\
go:schubfach-vs-bd & surrogate-pair & verify & 1.546 & 1.469 & 1.05$\times$ & json-canon \\
go:schubfach-vs-bd & unicode & canonicalize & 1.819 & 1.474 & 1.23$\times$ & json-canon \\
go:schubfach-vs-bd & unicode & verify & 1.488 & 1.524 & 1.02$\times$ & schubfach \\
go:schubfach-vs-bd & verify-whitespace & canonicalize & 1.478 & 1.634 & 1.11$\times$ & schubfach \\
go:schubfach-vs-bd & verify-whitespace & verify & 1.457 & 1.785 & 1.23$\times$ & schubfach \\
rs:schubfach-vs-bd & array-2048 & canonicalize & 2.664 & 4.606 & 1.73$\times$ & schubfach-rs \\
rs:schubfach-vs-bd & array-2048 & verify & 2.451 & 4.380 & 1.79$\times$ & schubfach-rs \\
rs:schubfach-vs-bd & array-256 & canonicalize & 0.866 & 1.149 & 1.33$\times$ & schubfach-rs \\
rs:schubfach-vs-bd & array-256 & verify & 0.817 & 1.206 & 1.48$\times$ & schubfach-rs \\
rs:schubfach-vs-bd & canonical-minimal & canonicalize & 0.571 & 0.619 & 1.08$\times$ & schubfach-rs \\
rs:schubfach-vs-bd & canonical-minimal & verify & 0.538 & 0.610 & 1.13$\times$ & schubfach-rs \\
rs:schubfach-vs-bd & control-escapes & canonicalize & 0.559 & 0.644 & 1.15$\times$ & schubfach-rs \\
rs:schubfach-vs-bd & control-escapes & verify & 0.535 & 0.653 & 1.22$\times$ & schubfach-rs \\
rs:schubfach-vs-bd & deep & canonicalize & 0.605 & 0.682 & 1.13$\times$ & schubfach-rs \\
rs:schubfach-vs-bd & deep & verify & 0.648 & 0.694 & 1.07$\times$ & schubfach-rs \\
rs:schubfach-vs-bd & deep-64 & canonicalize & 0.614 & 0.632 & 1.03$\times$ & schubfach-rs \\
rs:schubfach-vs-bd & deep-64 & verify & 0.563 & 0.644 & 1.14$\times$ & schubfach-rs \\
rs:schubfach-vs-bd & escaped-key-order & canonicalize & 0.556 & 0.685 & 1.23$\times$ & schubfach-rs \\
rs:schubfach-vs-bd & escaped-key-order & verify & 0.554 & 0.698 & 1.26$\times$ & schubfach-rs \\
rs:schubfach-vs-bd & large & canonicalize & 7.967 & 18.705 & 2.35$\times$ & schubfach-rs \\
rs:schubfach-vs-bd & large & verify & 7.214 & 18.635 & 2.58$\times$ & schubfach-rs \\
rs:schubfach-vs-bd & long-string & canonicalize & 0.630 & 0.716 & 1.14$\times$ & schubfach-rs \\
rs:schubfach-vs-bd & long-string & verify & 0.647 & 0.667 & 1.03$\times$ & schubfach-rs \\
rs:schubfach-vs-bd & medium & canonicalize & 0.869 & 1.338 & 1.54$\times$ & schubfach-rs \\
rs:schubfach-vs-bd & medium & verify & 0.861 & 1.366 & 1.59$\times$ & schubfach-rs \\
rs:schubfach-vs-bd & mixed-prod & canonicalize & 0.694 & 0.882 & 1.27$\times$ & schubfach-rs \\
rs:schubfach-vs-bd & mixed-prod & verify & 0.717 & 0.865 & 1.21$\times$ & schubfach-rs \\
rs:schubfach-vs-bd & nested-mixed & canonicalize & 0.559 & 0.689 & 1.23$\times$ & schubfach-rs \\
rs:schubfach-vs-bd & nested-mixed & verify & 0.569 & 0.700 & 1.23$\times$ & schubfach-rs \\
rs:schubfach-vs-bd & number-heavy & canonicalize & 0.570 & 0.735 & 1.29$\times$ & schubfach-rs \\
rs:schubfach-vs-bd & number-heavy & verify & 0.568 & 0.732 & 1.29$\times$ & schubfach-rs \\
rs:schubfach-vs-bd & numeric-boundary & canonicalize & 0.562 & 0.736 & 1.31$\times$ & schubfach-rs \\
rs:schubfach-vs-bd & numeric-boundary & verify & 0.546 & 0.718 & 1.31$\times$ & schubfach-rs \\
rs:schubfach-vs-bd & rfc-key-sorting & canonicalize & 0.578 & 0.629 & 1.09$\times$ & schubfach-rs \\
rs:schubfach-vs-bd & rfc-key-sorting & verify & 0.561 & 0.617 & 1.10$\times$ & schubfach-rs \\
rs:schubfach-vs-bd & small & canonicalize & 0.545 & 0.673 & 1.23$\times$ & schubfach-rs \\
rs:schubfach-vs-bd & small & verify & 0.588 & 0.704 & 1.20$\times$ & schubfach-rs \\
rs:schubfach-vs-bd & surrogate-pair & canonicalize & 0.550 & 0.599 & 1.09$\times$ & schubfach-rs \\
rs:schubfach-vs-bd & surrogate-pair & verify & 0.563 & 0.623 & 1.11$\times$ & schubfach-rs \\
rs:schubfach-vs-bd & unicode & canonicalize & 0.564 & 0.617 & 1.09$\times$ & schubfach-rs \\
rs:schubfach-vs-bd & unicode & verify & 0.543 & 0.644 & 1.19$\times$ & schubfach-rs \\
rs:schubfach-vs-bd & verify-whitespace & canonicalize & 0.547 & 0.719 & 1.31$\times$ & schubfach-rs \\
rs:schubfach-vs-bd & verify-whitespace & verify & 0.550 & 0.678 & 1.23$\times$ & schubfach-rs \\
schubfach:go-vs-rs & array-2048 & canonicalize & 5.642 & 2.664 & 2.12$\times$ & schubfach-rs \\
schubfach:go-vs-rs & array-2048 & verify & 5.694 & 2.451 & 2.32$\times$ & schubfach-rs \\
schubfach:go-vs-rs & array-256 & canonicalize & 1.840 & 0.866 & 2.13$\times$ & schubfach-rs \\
schubfach:go-vs-rs & array-256 & verify & 2.087 & 0.817 & 2.55$\times$ & schubfach-rs \\
schubfach:go-vs-rs & canonical-minimal & canonicalize & 1.378 & 0.571 & 2.41$\times$ & schubfach-rs \\
schubfach:go-vs-rs & canonical-minimal & verify & 1.443 & 0.538 & 2.68$\times$ & schubfach-rs \\
schubfach:go-vs-rs & control-escapes & canonicalize & 1.501 & 0.559 & 2.68$\times$ & schubfach-rs \\
schubfach:go-vs-rs & control-escapes & verify & 1.352 & 0.535 & 2.53$\times$ & schubfach-rs \\
schubfach:go-vs-rs & deep & canonicalize & 1.643 & 0.605 & 2.72$\times$ & schubfach-rs \\
schubfach:go-vs-rs & deep & verify & 1.704 & 0.648 & 2.63$\times$ & schubfach-rs \\
schubfach:go-vs-rs & deep-64 & canonicalize & 1.443 & 0.614 & 2.35$\times$ & schubfach-rs \\
schubfach:go-vs-rs & deep-64 & verify & 1.584 & 0.563 & 2.81$\times$ & schubfach-rs \\
schubfach:go-vs-rs & escaped-key-order & canonicalize & 1.489 & 0.556 & 2.68$\times$ & schubfach-rs \\
schubfach:go-vs-rs & escaped-key-order & verify & 1.543 & 0.554 & 2.78$\times$ & schubfach-rs \\
schubfach:go-vs-rs & large & canonicalize & 16.530 & 7.967 & 2.07$\times$ & schubfach-rs \\
schubfach:go-vs-rs & large & verify & 15.913 & 7.214 & 2.21$\times$ & schubfach-rs \\
schubfach:go-vs-rs & long-string & canonicalize & 1.566 & 0.630 & 2.48$\times$ & schubfach-rs \\
schubfach:go-vs-rs & long-string & verify & 1.791 & 0.647 & 2.77$\times$ & schubfach-rs \\
schubfach:go-vs-rs & medium & canonicalize & 1.986 & 0.869 & 2.29$\times$ & schubfach-rs \\
schubfach:go-vs-rs & medium & verify & 2.143 & 0.861 & 2.49$\times$ & schubfach-rs \\
schubfach:go-vs-rs & mixed-prod & canonicalize & 1.851 & 0.694 & 2.67$\times$ & schubfach-rs \\
schubfach:go-vs-rs & mixed-prod & verify & 1.614 & 0.717 & 2.25$\times$ & schubfach-rs \\
schubfach:go-vs-rs & nested-mixed & canonicalize & 1.481 & 0.559 & 2.65$\times$ & schubfach-rs \\
schubfach:go-vs-rs & nested-mixed & verify & 1.486 & 0.569 & 2.61$\times$ & schubfach-rs \\
schubfach:go-vs-rs & number-heavy & canonicalize & 1.484 & 0.570 & 2.61$\times$ & schubfach-rs \\
schubfach:go-vs-rs & number-heavy & verify & 1.523 & 0.568 & 2.68$\times$ & schubfach-rs \\
schubfach:go-vs-rs & numeric-boundary & canonicalize & 1.578 & 0.562 & 2.81$\times$ & schubfach-rs \\
schubfach:go-vs-rs & numeric-boundary & verify & 1.550 & 0.546 & 2.84$\times$ & schubfach-rs \\
schubfach:go-vs-rs & rfc-key-sorting & canonicalize & 1.442 & 0.578 & 2.49$\times$ & schubfach-rs \\
schubfach:go-vs-rs & rfc-key-sorting & verify & 1.499 & 0.561 & 2.67$\times$ & schubfach-rs \\
schubfach:go-vs-rs & small & canonicalize & 1.453 & 0.545 & 2.67$\times$ & schubfach-rs \\
schubfach:go-vs-rs & small & verify & 1.561 & 0.588 & 2.65$\times$ & schubfach-rs \\
schubfach:go-vs-rs & surrogate-pair & canonicalize & 1.425 & 0.550 & 2.59$\times$ & schubfach-rs \\
schubfach:go-vs-rs & surrogate-pair & verify & 1.546 & 0.563 & 2.75$\times$ & schubfach-rs \\
schubfach:go-vs-rs & unicode & canonicalize & 1.819 & 0.564 & 3.23$\times$ & schubfach-rs \\
schubfach:go-vs-rs & unicode & verify & 1.488 & 0.543 & 2.74$\times$ & schubfach-rs \\
schubfach:go-vs-rs & verify-whitespace & canonicalize & 1.478 & 0.547 & 2.70$\times$ & schubfach-rs \\
schubfach:go-vs-rs & verify-whitespace & verify & 1.457 & 0.550 & 2.65$\times$ & schubfach-rs \\
bd:go-vs-rs & array-2048 & canonicalize & 7.202 & 4.606 & 1.56$\times$ & json-canon-rs \\
bd:go-vs-rs & array-2048 & verify & 7.125 & 4.380 & 1.63$\times$ & json-canon-rs \\
bd:go-vs-rs & array-256 & canonicalize & 2.383 & 1.149 & 2.07$\times$ & json-canon-rs \\
bd:go-vs-rs & array-256 & verify & 2.246 & 1.206 & 1.86$\times$ & json-canon-rs \\
bd:go-vs-rs & canonical-minimal & canonicalize & 1.655 & 0.619 & 2.67$\times$ & json-canon-rs \\
bd:go-vs-rs & canonical-minimal & verify & 1.654 & 0.610 & 2.71$\times$ & json-canon-rs \\
bd:go-vs-rs & control-escapes & canonicalize & 1.569 & 0.644 & 2.44$\times$ & json-canon-rs \\
bd:go-vs-rs & control-escapes & verify & 1.697 & 0.653 & 2.60$\times$ & json-canon-rs \\
bd:go-vs-rs & deep & canonicalize & 1.722 & 0.682 & 2.53$\times$ & json-canon-rs \\
bd:go-vs-rs & deep & verify & 1.723 & 0.694 & 2.48$\times$ & json-canon-rs \\
bd:go-vs-rs & deep-64 & canonicalize & 1.828 & 0.632 & 2.89$\times$ & json-canon-rs \\
bd:go-vs-rs & deep-64 & verify & 1.636 & 0.644 & 2.54$\times$ & json-canon-rs \\
bd:go-vs-rs & escaped-key-order & canonicalize & 1.553 & 0.685 & 2.27$\times$ & json-canon-rs \\
bd:go-vs-rs & escaped-key-order & verify & 1.546 & 0.698 & 2.22$\times$ & json-canon-rs \\
bd:go-vs-rs & large & canonicalize & 24.906 & 18.705 & 1.33$\times$ & json-canon-rs \\
bd:go-vs-rs & large & verify & 24.895 & 18.635 & 1.34$\times$ & json-canon-rs \\
bd:go-vs-rs & long-string & canonicalize & 1.821 & 0.716 & 2.54$\times$ & json-canon-rs \\
bd:go-vs-rs & long-string & verify & 1.912 & 0.667 & 2.87$\times$ & json-canon-rs \\
bd:go-vs-rs & medium & canonicalize & 2.659 & 1.338 & 1.99$\times$ & json-canon-rs \\
bd:go-vs-rs & medium & verify & 2.531 & 1.366 & 1.85$\times$ & json-canon-rs \\
bd:go-vs-rs & mixed-prod & canonicalize & 1.907 & 0.882 & 2.16$\times$ & json-canon-rs \\
bd:go-vs-rs & mixed-prod & verify & 1.877 & 0.865 & 2.17$\times$ & json-canon-rs \\
bd:go-vs-rs & nested-mixed & canonicalize & 1.613 & 0.689 & 2.34$\times$ & json-canon-rs \\
bd:go-vs-rs & nested-mixed & verify & 1.739 & 0.700 & 2.48$\times$ & json-canon-rs \\
bd:go-vs-rs & number-heavy & canonicalize & 1.656 & 0.735 & 2.25$\times$ & json-canon-rs \\
bd:go-vs-rs & number-heavy & verify & 1.763 & 0.732 & 2.41$\times$ & json-canon-rs \\
bd:go-vs-rs & numeric-boundary & canonicalize & 1.556 & 0.736 & 2.11$\times$ & json-canon-rs \\
bd:go-vs-rs & numeric-boundary & verify & 1.627 & 0.718 & 2.27$\times$ & json-canon-rs \\
bd:go-vs-rs & rfc-key-sorting & canonicalize & 1.564 & 0.629 & 2.48$\times$ & json-canon-rs \\
bd:go-vs-rs & rfc-key-sorting & verify & 1.566 & 0.617 & 2.54$\times$ & json-canon-rs \\
bd:go-vs-rs & small & canonicalize & 1.588 & 0.673 & 2.36$\times$ & json-canon-rs \\
bd:go-vs-rs & small & verify & 1.653 & 0.704 & 2.35$\times$ & json-canon-rs \\
bd:go-vs-rs & surrogate-pair & canonicalize & 1.549 & 0.599 & 2.58$\times$ & json-canon-rs \\
bd:go-vs-rs & surrogate-pair & verify & 1.469 & 0.623 & 2.36$\times$ & json-canon-rs \\
bd:go-vs-rs & unicode & canonicalize & 1.474 & 0.617 & 2.39$\times$ & json-canon-rs \\
bd:go-vs-rs & unicode & verify & 1.524 & 0.644 & 2.37$\times$ & json-canon-rs \\
bd:go-vs-rs & verify-whitespace & canonicalize & 1.634 & 0.719 & 2.27$\times$ & json-canon-rs \\
bd:go-vs-rs & verify-whitespace & verify & 1.785 & 0.678 & 2.63$\times$ & json-canon-rs \\
\end{longtable}
\normalsize

\end{landscape}

\clearpage

% ─── Section 3: Full Statistical Comparisons ─────────────────────────────────

\section{Full Statistical Comparisons}
\label{sec:supp-stats}

Table~\ref{tab:stats-comparisons} reports all 1{,}644 statistical
comparisons across 16~pairs, 3~tracks (API, CLI end-to-end, CLI
algorithmic), 2~modes (canonicalize, verify), and all workloads.
Each comparison includes:

\begin{itemize}
\item \textbf{Speedup:} ratio of slower to faster mean.
\item \textbf{CI$_{95}$:} BCa bootstrap 95\% confidence interval on
      the speedup ratio (5{,}000 resamples).
\item \textbf{$p$:} raw permutation test $p$-value (5{,}000 resamples).
\item \textbf{$p_\text{adj}$:} Benjamini--Hochberg adjusted $p$-value
      controlling false discovery rate at $q = 0.05$ across all
      1{,}644 comparisons.
\item \textbf{$d$:} Cohen's $d$ effect size
      ($(\bar{x}_A - \bar{x}_B)/s_p$; negative = implementation A
      faster).
\item \textbf{Noise:} standard deviation of the faster implementation,
      quantifying measurement noise in the same units as the metric.
\item \textbf{MDE:} minimum detectable effect at 80\% power, expressed
      as a percentage of the faster implementation's mean.
\end{itemize}

\noindent
$p$-values at the resolution floor of $1/(R+1) = 0.0002$ (where
$R = 5{,}000$ resamples) should be interpreted as $p < 0.0002$.

\begin{landscape}
% Auto-generated by scripts/generate-paper-tables.go
\scriptsize
\begin{longtable}{lllllllrlrrrrr}
\caption{Statistical comparisons with Benjamini-Hochberg corrected $p$-values.}
\label{tab:stats-comparisons}\\
\toprule
Pair & Track & Mode & Workload & Impl A & Impl B & Winner & Speedup & CI$_{95}$ & $p$ & $p_\text{adj}$ & $d$ & Noise (ms) & MDE (\%) \\
\midrule
\endfirsthead
\toprule
Pair & Track & Mode & Workload & Impl A & Impl B & Winner & Speedup & CI$_{95}$ & $p$ & $p_\text{adj}$ & $d$ & Noise (ms) & MDE (\%) \\
\midrule
\endhead
\midrule
\multicolumn{14}{r}{\textit{Continued on next page}} \\
\endfoot
\bottomrule
\multicolumn{14}{l}{\footnotesize $^*$Significant after Benjamini-Hochberg FDR correction at $q = 0.05$.} \\
\multicolumn{14}{l}{\footnotesize $d$ = Cohen's $d$: $(\bar{x}_A - \bar{x}_B)/s_p$; negative = impl A faster. Track: e2e = end-to-end CLI.} \\
\multicolumn{14}{l}{\footnotesize Noise = CV $\times$ mean of faster impl. MDE = minimum detectable effect at 80\% power ($z$-approx).} \\
\multicolumn{14}{l}{\footnotesize Rust allocs/op and B/op measured via counting global allocator (\S\ref{sec:discussion}).} \\
\endlastfoot
go:schubfach-vs-bd & e2e & canonicalize & array-2048 & schubfach & json-canon & schubfach$^*$ & 1.25$\times$ & [1.183, 1.331] & 0.0002 & 0.0004 & -2.69 & 0.5247 & 9.4 \\
go:schubfach-vs-ryu & e2e & canonicalize & array-2048 & schubfach & ryu & ryu & 1.05$\times$ & [0.985, 1.108] & 0.1200 & 0.1896 & 0.57 & 0.4264 & 8.0 \\
go:schubfach-vs-dragonbox & e2e & canonicalize & array-2048 & schubfach & dragonbox & dragonbox & 1.07$\times$ & [0.993, 1.139] & 0.0886 & 0.1451 & 0.65 & 0.5575 & 10.6 \\
go:ryu-vs-dragonbox & e2e & canonicalize & array-2048 & ryu & dragonbox & dragonbox & 1.01$\times$ & [0.952, 1.081] & 0.6629 & 0.7459 & 0.16 & 0.5575 & 10.6 \\
go:ryu-vs-bd & e2e & canonicalize & array-2048 & ryu & json-canon & ryu$^*$ & 1.31$\times$ & [1.243, 1.384] & 0.0002 & 0.0004 & -3.51 & 0.4264 & 8.0 \\
go:dragonbox-vs-bd & e2e & canonicalize & array-2048 & dragonbox & json-canon & dragonbox$^*$ & 1.33$\times$ & [1.258, 1.422] & 0.0002 & 0.0004 & -3.25 & 0.5575 & 10.6 \\
rs:schubfach-vs-bd & e2e & canonicalize & array-2048 & schubfach-rs & json-canon-rs & schubfach-rs$^*$ & 1.59$\times$ & [1.436, 1.802] & 0.0002 & 0.0004 & -2.64 & 0.3842 & 15.5 \\
rs:schubfach-vs-ryu & e2e & canonicalize & array-2048 & schubfach-rs & ryu-rs & schubfach-rs & 1.03$\times$ & [0.914, 1.148] & 0.6731 & 0.7528 & -0.16 & 0.3842 & 15.5 \\
rs:schubfach-vs-dragonbox & e2e & canonicalize & array-2048 & schubfach-rs & dragonbox-rs & dragonbox-rs & 1.03$\times$ & [0.909, 1.153] & 0.6665 & 0.7478 & 0.16 & 0.4617 & 19.1 \\
rs:ryu-vs-dragonbox & e2e & canonicalize & array-2048 & ryu-rs & dragonbox-rs & dragonbox-rs & 1.05$\times$ & [0.925, 1.187] & 0.4471 & 0.5568 & 0.28 & 0.4617 & 19.1 \\
rs:ryu-vs-bd & e2e & canonicalize & array-2048 & ryu-rs & json-canon-rs & ryu-rs$^*$ & 1.55$\times$ & [1.388, 1.771] & 0.0002 & 0.0004 & -2.40 & 0.4683 & 18.4 \\
rs:dragonbox-vs-bd & e2e & canonicalize & array-2048 & dragonbox-rs & json-canon-rs & dragonbox-rs$^*$ & 1.64$\times$ & [1.449, 1.855] & 0.0002 & 0.0004 & -2.63 & 0.4617 & 19.1 \\
schubfach:go-vs-rs & e2e & canonicalize & array-2048 & schubfach & schubfach-rs & schubfach-rs$^*$ & 2.25$\times$ & [2.060, 2.450] & 0.0002 & 0.0004 & 6.91 & 0.3842 & 15.5 \\
bd:go-vs-rs & e2e & canonicalize & array-2048 & json-canon & json-canon-rs & json-canon-rs$^*$ & 1.77$\times$ & [1.591, 1.916] & 0.0002 & 0.0004 & 4.93 & 0.7084 & 17.9 \\
ryu:go-vs-rs & e2e & canonicalize & array-2048 & ryu & ryu-rs & ryu-rs$^*$ & 2.09$\times$ & [1.913, 2.316] & 0.0002 & 0.0004 & 6.34 & 0.4683 & 18.4 \\
dragonbox:go-vs-rs & e2e & canonicalize & array-2048 & dragonbox & dragonbox-rs & dragonbox-rs$^*$ & 2.17$\times$ & [1.947, 2.396] & 0.0002 & 0.0004 & 5.65 & 0.4617 & 19.1 \\
go:schubfach-vs-bd & e2e & canonicalize & array-256 & schubfach & json-canon & schubfach & 1.16$\times$ & [1.024, 1.320] & 0.0366 & 0.0650 & -0.80 & 0.3590 & 18.3 \\
go:schubfach-vs-ryu & e2e & canonicalize & array-256 & schubfach & ryu & schubfach & 1.04$\times$ & [0.924, 1.173] & 0.5405 & 0.6387 & -0.22 & 0.3590 & 18.3 \\
go:schubfach-vs-dragonbox & e2e & canonicalize & array-256 & schubfach & dragonbox & dragonbox & 1.09$\times$ & [0.950, 1.236] & 0.2196 & 0.3119 & 0.45 & 0.3744 & 20.8 \\
go:ryu-vs-dragonbox & e2e & canonicalize & array-256 & ryu & dragonbox & dragonbox & 1.13$\times$ & [0.991, 1.275] & 0.0762 & 0.1269 & 0.68 & 0.3744 & 20.8 \\
go:ryu-vs-bd & e2e & canonicalize & array-256 & ryu & json-canon & ryu & 1.12$\times$ & [0.991, 1.267] & 0.0996 & 0.1597 & -0.63 & 0.3352 & 16.4 \\
go:dragonbox-vs-bd & e2e & canonicalize & array-256 & dragonbox & json-canon & dragonbox$^*$ & 1.26$\times$ & [1.093, 1.440] & 0.0040 & 0.0079 & -1.20 & 0.3744 & 20.8 \\
rs:schubfach-vs-bd & e2e & canonicalize & array-256 & schubfach-rs & json-canon-rs & schubfach-rs$^*$ & 1.36$\times$ & [1.234, 1.484] & 0.0002 & 0.0004 & -2.21 & 0.1020 & 12.7 \\
rs:schubfach-vs-ryu & e2e & canonicalize & array-256 & schubfach-rs & ryu-rs & schubfach-rs & 1.02$\times$ & [0.943, 1.096] & 0.6153 & 0.7077 & -0.18 & 0.1020 & 12.7 \\
rs:schubfach-vs-dragonbox & e2e & canonicalize & array-256 & schubfach-rs & dragonbox-rs & schubfach-rs & 1.04$\times$ & [0.956, 1.120] & 0.3907 & 0.4991 & -0.32 & 0.1020 & 12.7 \\
rs:ryu-vs-dragonbox & e2e & canonicalize & array-256 & ryu-rs & dragonbox-rs & ryu-rs & 1.02$\times$ & [0.951, 1.088] & 0.6533 & 0.7394 & -0.17 & 0.0772 & 9.4 \\
rs:ryu-vs-bd & e2e & canonicalize & array-256 & ryu-rs & json-canon-rs & ryu-rs$^*$ & 1.33$\times$ & [1.222, 1.443] & 0.0002 & 0.0004 & -2.23 & 0.0772 & 9.4 \\
rs:dragonbox-vs-bd & e2e & canonicalize & array-256 & dragonbox-rs & json-canon-rs & dragonbox-rs$^*$ & 1.31$\times$ & [1.201, 1.423] & 0.0002 & 0.0004 & -2.09 & 0.0835 & 10.0 \\
schubfach:go-vs-rs & e2e & canonicalize & array-256 & schubfach & schubfach-rs & schubfach-rs$^*$ & 2.45$\times$ & [2.190, 2.709] & 0.0002 & 0.0004 & 4.50 & 0.1020 & 12.7 \\
bd:go-vs-rs & e2e & canonicalize & array-256 & json-canon & json-canon-rs & json-canon-rs$^*$ & 2.09$\times$ & [1.876, 2.353] & 0.0002 & 0.0004 & 3.73 & 0.1566 & 14.4 \\
ryu:go-vs-rs & e2e & canonicalize & array-256 & ryu & ryu-rs & ryu-rs$^*$ & 2.49$\times$ & [2.277, 2.729] & 0.0002 & 0.0004 & 5.13 & 0.0772 & 9.4 \\
dragonbox:go-vs-rs & e2e & canonicalize & array-256 & dragonbox & dragonbox-rs & dragonbox-rs$^*$ & 2.17$\times$ & [1.959, 2.430] & 0.0002 & 0.0004 & 3.66 & 0.0835 & 10.0 \\
go:schubfach-vs-bd & e2e & canonicalize & canonical-minimal & schubfach & json-canon & schubfach & 1.02$\times$ & [0.901, 1.171] & 0.7495 & 0.8143 & -0.12 & 0.2626 & 18.0 \\
go:schubfach-vs-ryu & e2e & canonicalize & canonical-minimal & schubfach & ryu & ryu & 1.08$\times$ & [0.937, 1.251] & 0.2647 & 0.3652 & 0.40 & 0.3153 & 23.4 \\
go:schubfach-vs-dragonbox & e2e & canonicalize & canonical-minimal & schubfach & dragonbox & dragonbox & 1.03$\times$ & [0.910, 1.158] & 0.6647 & 0.7472 & 0.16 & 0.2496 & 17.5 \\
go:ryu-vs-dragonbox & e2e & canonicalize & canonical-minimal & ryu & dragonbox & ryu & 1.06$\times$ & [0.919, 1.211] & 0.4737 & 0.5777 & -0.27 & 0.3153 & 23.4 \\
go:ryu-vs-bd & e2e & canonicalize & canonical-minimal & ryu & json-canon & ryu & 1.11$\times$ & [0.959, 1.288] & 0.1990 & 0.2873 & -0.48 & 0.3153 & 23.4 \\
go:dragonbox-vs-bd & e2e & canonicalize & canonical-minimal & dragonbox & json-canon & dragonbox & 1.05$\times$ & [0.920, 1.193] & 0.4673 & 0.5743 & -0.26 & 0.2496 & 17.5 \\
rs:schubfach-vs-bd & e2e & canonicalize & canonical-minimal & schubfach-rs & json-canon-rs & schubfach-rs$^*$ & 1.10$\times$ & [1.044, 1.146] & 0.0002 & 0.0004 & -1.43 & 0.0351 & 6.4 \\
rs:schubfach-vs-ryu & e2e & canonicalize & canonical-minimal & schubfach-rs & ryu-rs & schubfach-rs & 1.03$\times$ & [0.962, 1.096] & 0.4939 & 0.5956 & -0.27 & 0.0351 & 6.4 \\
rs:schubfach-vs-dragonbox & e2e & canonicalize & canonical-minimal & schubfach-rs & dragonbox-rs & dragonbox-rs & 1.04$\times$ & [0.980, 1.119] & 0.2813 & 0.3804 & 0.40 & 0.0633 & 12.0 \\
rs:ryu-vs-dragonbox & e2e & canonicalize & canonical-minimal & ryu-rs & dragonbox-rs & dragonbox-rs & 1.06$\times$ & [0.988, 1.163] & 0.1534 & 0.2343 & 0.54 & 0.0633 & 12.0 \\
rs:ryu-vs-bd & e2e & canonicalize & canonical-minimal & ryu-rs & json-canon-rs & ryu-rs & 1.07$\times$ & [0.999, 1.138] & 0.0474 & 0.0818 & -0.74 & 0.0654 & 11.6 \\
rs:dragonbox-vs-bd & e2e & canonicalize & canonical-minimal & dragonbox-rs & json-canon-rs & dragonbox-rs$^*$ & 1.14$\times$ & [1.073, 1.230] & 0.0002 & 0.0004 & -1.41 & 0.0633 & 12.0 \\
schubfach:go-vs-rs & e2e & canonicalize & canonical-minimal & schubfach & schubfach-rs & schubfach-rs$^*$ & 2.67$\times$ & [2.385, 2.885] & 0.0002 & 0.0004 & 4.99 & 0.0351 & 6.4 \\
bd:go-vs-rs & e2e & canonicalize & canonical-minimal & json-canon & json-canon-rs & json-canon-rs$^*$ & 2.48$\times$ & [2.239, 2.750] & 0.0002 & 0.0004 & 4.22 & 0.0404 & 6.7 \\
ryu:go-vs-rs & e2e & canonicalize & canonical-minimal & ryu & ryu-rs & ryu-rs$^*$ & 2.40$\times$ & [2.097, 2.694] & 0.0002 & 0.0004 & 3.53 & 0.0654 & 11.6 \\
dragonbox:go-vs-rs & e2e & canonicalize & canonical-minimal & dragonbox & dragonbox-rs & dragonbox-rs$^*$ & 2.69$\times$ & [2.439, 2.992] & 0.0002 & 0.0004 & 5.03 & 0.0633 & 12.0 \\
go:schubfach-vs-bd & e2e & canonicalize & control-escapes & schubfach & json-canon & schubfach & 1.09$\times$ & [0.968, 1.232] & 0.1678 & 0.2516 & -0.51 & 0.2523 & 17.6 \\
go:schubfach-vs-ryu & e2e & canonicalize & control-escapes & schubfach & ryu & ryu & 1.06$\times$ & [0.933, 1.199] & 0.3803 & 0.4877 & 0.32 & 0.2659 & 19.7 \\
go:schubfach-vs-dragonbox & e2e & canonicalize & control-escapes & schubfach & dragonbox & schubfach & 1.01$\times$ & [0.889, 1.135] & 0.9060 & 0.9342 & -0.04 & 0.2523 & 17.6 \\
go:ryu-vs-dragonbox & e2e & canonicalize & control-escapes & ryu & dragonbox & ryu & 1.07$\times$ & [0.933, 1.205] & 0.3201 & 0.4214 & -0.36 & 0.2659 & 19.7 \\
go:ryu-vs-bd & e2e & canonicalize & control-escapes & ryu & json-canon & ryu & 1.16$\times$ & [1.018, 1.319] & 0.0392 & 0.0692 & -0.80 & 0.2659 & 19.7 \\
go:dragonbox-vs-bd & e2e & canonicalize & control-escapes & dragonbox & json-canon & dragonbox & 1.09$\times$ & [0.961, 1.238] & 0.2164 & 0.3086 & -0.46 & 0.2680 & 18.6 \\
rs:schubfach-vs-bd & e2e & canonicalize & control-escapes & schubfach-rs & json-canon-rs & schubfach-rs$^*$ & 1.11$\times$ & [1.046, 1.186] & 0.0038 & 0.0075 & -1.12 & 0.0375 & 6.9 \\
rs:schubfach-vs-ryu & e2e & canonicalize & control-escapes & schubfach-rs & ryu-rs & schubfach-rs & 1.01$\times$ & [0.946, 1.063] & 0.8016 & 0.8538 & -0.09 & 0.0375 & 6.9 \\
rs:schubfach-vs-dragonbox & e2e & canonicalize & control-escapes & schubfach-rs & dragonbox-rs & schubfach-rs & 1.03$\times$ & [0.950, 1.092] & 0.4979 & 0.5998 & -0.26 & 0.0375 & 6.9 \\
rs:ryu-vs-dragonbox & e2e & canonicalize & control-escapes & ryu-rs & dragonbox-rs & ryu-rs & 1.02$\times$ & [0.938, 1.095] & 0.6675 & 0.7478 & -0.16 & 0.0532 & 9.8 \\
rs:ryu-vs-bd & e2e & canonicalize & control-escapes & ryu-rs & json-canon-rs & ryu-rs$^*$ & 1.10$\times$ & [1.027, 1.186] & 0.0130 & 0.0245 & -0.94 & 0.0532 & 9.8 \\
rs:dragonbox-vs-bd & e2e & canonicalize & control-escapes & dragonbox-rs & json-canon-rs & dragonbox-rs & 1.08$\times$ & [1.001, 1.183] & 0.0616 & 0.1050 & -0.69 & 0.0700 & 12.6 \\
schubfach:go-vs-rs & e2e & canonicalize & control-escapes & schubfach & schubfach-rs & schubfach-rs$^*$ & 2.65$\times$ & [2.396, 2.890] & 0.0002 & 0.0004 & 5.06 & 0.0375 & 6.9 \\
bd:go-vs-rs & e2e & canonicalize & control-escapes & json-canon & json-canon-rs & json-canon-rs$^*$ & 2.61$\times$ & [2.327, 2.872] & 0.0002 & 0.0004 & 4.71 & 0.0688 & 11.5 \\
ryu:go-vs-rs & e2e & canonicalize & control-escapes & ryu & ryu-rs & ryu-rs$^*$ & 2.48$\times$ & [2.240, 2.774] & 0.0002 & 0.0004 & 4.30 & 0.0532 & 9.8 \\
dragonbox:go-vs-rs & e2e & canonicalize & control-escapes & dragonbox & dragonbox-rs & dragonbox-rs$^*$ & 2.60$\times$ & [2.330, 2.886] & 0.0002 & 0.0004 & 4.64 & 0.0700 & 12.6 \\
go:schubfach-vs-bd & e2e & canonicalize & deep-64 & schubfach & json-canon & schubfach & 1.14$\times$ & [0.988, 1.321] & 0.0814 & 0.1347 & -0.66 & 0.3537 & 23.8 \\
go:schubfach-vs-ryu & e2e & canonicalize & deep-64 & schubfach & ryu & schubfach & 1.06$\times$ & [0.912, 1.228] & 0.4779 & 0.5807 & -0.26 & 0.3537 & 23.8 \\
go:schubfach-vs-dragonbox & e2e & canonicalize & deep-64 & schubfach & dragonbox & schubfach & 1.06$\times$ & [0.929, 1.212] & 0.3929 & 0.5004 & -0.32 & 0.3537 & 23.8 \\
go:ryu-vs-dragonbox & e2e & canonicalize & deep-64 & ryu & dragonbox & ryu & 1.01$\times$ & [0.893, 1.129] & 0.9354 & 0.9569 & -0.03 & 0.3071 & 19.5 \\
go:ryu-vs-bd & e2e & canonicalize & deep-64 & ryu & json-canon & ryu & 1.08$\times$ & [0.946, 1.214] & 0.2486 & 0.3485 & -0.43 & 0.3071 & 19.5 \\
go:dragonbox-vs-bd & e2e & canonicalize & deep-64 & dragonbox & json-canon & dragonbox & 1.08$\times$ & [0.954, 1.197] & 0.2162 & 0.3086 & -0.45 & 0.2336 & 14.8 \\
rs:schubfach-vs-bd & e2e & canonicalize & deep-64 & schubfach-rs & json-canon-rs & schubfach-rs & 1.09$\times$ & [0.977, 1.229] & 0.1444 & 0.2232 & -0.54 & 0.1136 & 20.1 \\
rs:schubfach-vs-ryu & e2e & canonicalize & deep-64 & schubfach-rs & ryu-rs & schubfach-rs & 1.03$\times$ & [0.929, 1.154] & 0.6057 & 0.6985 & -0.20 & 0.1136 & 20.1 \\
rs:schubfach-vs-dragonbox & e2e & canonicalize & deep-64 & schubfach-rs & dragonbox-rs & schubfach-rs & 1.02$\times$ & [0.904, 1.153] & 0.8100 & 0.8613 & -0.09 & 0.1136 & 20.1 \\
rs:ryu-vs-dragonbox & e2e & canonicalize & deep-64 & ryu-rs & dragonbox-rs & dragonbox-rs & 1.01$\times$ & [0.927, 1.104] & 0.7812 & 0.8389 & 0.11 & 0.0906 & 15.8 \\
rs:ryu-vs-bd & e2e & canonicalize & deep-64 & ryu-rs & json-canon-rs & ryu-rs & 1.06$\times$ & [0.974, 1.130] & 0.1554 & 0.2368 & -0.54 & 0.0522 & 9.0 \\
rs:dragonbox-vs-bd & e2e & canonicalize & deep-64 & dragonbox-rs & json-canon-rs & dragonbox-rs & 1.07$\times$ & [0.971, 1.175] & 0.1812 & 0.2658 & -0.51 & 0.0906 & 15.8 \\
schubfach:go-vs-rs & e2e & canonicalize & deep-64 & schubfach & schubfach-rs & schubfach-rs$^*$ & 2.64$\times$ & [2.275, 3.052] & 0.0002 & 0.0004 & 3.60 & 0.1136 & 20.1 \\
bd:go-vs-rs & e2e & canonicalize & deep-64 & json-canon & json-canon-rs & json-canon-rs$^*$ & 2.76$\times$ & [2.468, 3.036] & 0.0002 & 0.0004 & 5.00 & 0.0749 & 12.2 \\
ryu:go-vs-rs & e2e & canonicalize & deep-64 & ryu & ryu-rs & ryu-rs$^*$ & 2.70$\times$ & [2.436, 2.991] & 0.0002 & 0.0004 & 4.60 & 0.0522 & 9.0 \\
dragonbox:go-vs-rs & e2e & canonicalize & deep-64 & dragonbox & dragonbox-rs & dragonbox-rs$^*$ & 2.76$\times$ & [2.468, 3.038] & 0.0002 & 0.0004 & 5.81 & 0.0906 & 15.8 \\
go:schubfach-vs-bd & e2e & canonicalize & deep & schubfach & json-canon & schubfach & 1.17$\times$ & [1.026, 1.336] & 0.0326 & 0.0586 & -0.85 & 0.2718 & 18.7 \\
go:schubfach-vs-ryu & e2e & canonicalize & deep & schubfach & ryu & schubfach & 1.10$\times$ & [0.971, 1.238] & 0.1194 & 0.1888 & -0.58 & 0.2718 & 18.7 \\
go:schubfach-vs-dragonbox & e2e & canonicalize & deep & schubfach & dragonbox & schubfach$^*$ & 1.15$\times$ & [1.022, 1.273] & 0.0190 & 0.0353 & -0.89 & 0.2718 & 18.7 \\
go:ryu-vs-dragonbox & e2e & canonicalize & deep & ryu & dragonbox & ryu & 1.04$\times$ & [0.933, 1.154] & 0.4619 & 0.5698 & -0.27 & 0.2663 & 16.6 \\
go:ryu-vs-bd & e2e & canonicalize & deep & ryu & json-canon & ryu & 1.06$\times$ & [0.936, 1.201] & 0.3487 & 0.4531 & -0.34 & 0.2663 & 16.6 \\
go:dragonbox-vs-bd & e2e & canonicalize & deep & dragonbox & json-canon & dragonbox & 1.02$\times$ & [0.907, 1.146] & 0.7333 & 0.7994 & -0.13 & 0.2243 & 13.4 \\
rs:schubfach-vs-bd & e2e & canonicalize & deep & schubfach-rs & json-canon-rs & schubfach-rs$^*$ & 1.07$\times$ & [1.011, 1.121] & 0.0124 & 0.0235 & -0.96 & 0.0432 & 6.7 \\
rs:schubfach-vs-ryu & e2e & canonicalize & deep & schubfach-rs & ryu-rs & ryu-rs$^*$ & 1.09$\times$ & [1.025, 1.180] & 0.0184 & 0.0342 & 0.90 & 0.0758 & 12.8 \\
rs:schubfach-vs-dragonbox & e2e & canonicalize & deep & schubfach-rs & dragonbox-rs & dragonbox-rs & 1.06$\times$ & [1.012, 1.144] & 0.0494 & 0.0849 & 0.73 & 0.0618 & 10.2 \\
rs:ryu-vs-dragonbox & e2e & canonicalize & deep & ryu-rs & dragonbox-rs & ryu-rs & 1.03$\times$ & [0.949, 1.110] & 0.5207 & 0.6215 & -0.24 & 0.0758 & 12.8 \\
rs:ryu-vs-bd & e2e & canonicalize & deep & ryu-rs & json-canon-rs & ryu-rs$^*$ & 1.17$\times$ & [1.087, 1.259] & 0.0002 & 0.0004 & -1.55 & 0.0758 & 12.8 \\
rs:dragonbox-vs-bd & e2e & canonicalize & deep & dragonbox-rs & json-canon-rs & dragonbox-rs$^*$ & 1.14$\times$ & [1.077, 1.222] & 0.0004 & 0.0008 & -1.47 & 0.0618 & 10.2 \\
schubfach:go-vs-rs & e2e & canonicalize & deep & schubfach & schubfach-rs & schubfach-rs$^*$ & 2.26$\times$ & [2.058, 2.487] & 0.0002 & 0.0004 & 4.26 & 0.0432 & 6.7 \\
bd:go-vs-rs & e2e & canonicalize & deep & json-canon & json-canon-rs & json-canon-rs$^*$ & 2.47$\times$ & [2.217, 2.723] & 0.0002 & 0.0004 & 4.29 & 0.0566 & 8.2 \\
ryu:go-vs-rs & e2e & canonicalize & deep & ryu & ryu-rs & ryu-rs$^*$ & 2.72$\times$ & [2.461, 3.006] & 0.0002 & 0.0004 & 5.32 & 0.0758 & 12.8 \\
dragonbox:go-vs-rs & e2e & canonicalize & deep & dragonbox & dragonbox-rs & dragonbox-rs$^*$ & 2.76$\times$ & [2.528, 2.968] & 0.0002 & 0.0004 & 6.63 & 0.0618 & 10.2 \\
go:schubfach-vs-bd & e2e & canonicalize & escaped-key-order & schubfach & json-canon & schubfach & 1.18$\times$ & [0.996, 1.716] & 0.2511 & 0.3506 & -0.44 & 0.2547 & 17.1 \\
go:schubfach-vs-ryu & e2e & canonicalize & escaped-key-order & schubfach & ryu & schubfach & 1.07$\times$ & [0.965, 1.205] & 0.2691 & 0.3705 & -0.42 & 0.2547 & 17.1 \\
go:schubfach-vs-dragonbox & e2e & canonicalize & escaped-key-order & schubfach & dragonbox & dragonbox & 1.03$\times$ & [0.898, 1.173] & 0.6953 & 0.7670 & 0.14 & 0.3239 & 22.3 \\
go:ryu-vs-dragonbox & e2e & canonicalize & escaped-key-order & ryu & dragonbox & dragonbox & 1.10$\times$ & [0.969, 1.260] & 0.1730 & 0.2564 & 0.51 & 0.3239 & 22.3 \\
go:ryu-vs-bd & e2e & canonicalize & escaped-key-order & ryu & json-canon & ryu & 1.10$\times$ & [0.930, 1.564] & 0.6523 & 0.7394 & -0.27 & 0.2519 & 15.8 \\
go:dragonbox-vs-bd & e2e & canonicalize & escaped-key-order & dragonbox & json-canon & dragonbox & 1.21$\times$ & [1.007, 1.710] & 0.1700 & 0.2534 & -0.49 & 0.3239 & 22.3 \\
rs:schubfach-vs-bd & e2e & canonicalize & escaped-key-order & schubfach-rs & json-canon-rs & schubfach-rs$^*$ & 1.25$\times$ & [1.160, 1.366] & 0.0002 & 0.0004 & -1.98 & 0.0680 & 12.8 \\
rs:schubfach-vs-ryu & e2e & canonicalize & escaped-key-order & schubfach-rs & ryu-rs & ryu-rs & 1.01$\times$ & [0.922, 1.103] & 0.8738 & 0.9098 & 0.06 & 0.0728 & 13.9 \\
rs:schubfach-vs-dragonbox & e2e & canonicalize & escaped-key-order & schubfach-rs & dragonbox-rs & schubfach-rs & 1.01$\times$ & [0.944, 1.104] & 0.7840 & 0.8394 & -0.10 & 0.0680 & 12.8 \\
rs:ryu-vs-dragonbox & e2e & canonicalize & escaped-key-order & ryu-rs & dragonbox-rs & ryu-rs & 1.02$\times$ & [0.945, 1.109] & 0.6605 & 0.7443 & -0.16 & 0.0728 & 13.9 \\
rs:ryu-vs-bd & e2e & canonicalize & escaped-key-order & ryu-rs & json-canon-rs & ryu-rs$^*$ & 1.26$\times$ & [1.165, 1.376] & 0.0002 & 0.0004 & -1.97 & 0.0728 & 13.9 \\
rs:dragonbox-vs-bd & e2e & canonicalize & escaped-key-order & dragonbox-rs & json-canon-rs & dragonbox-rs$^*$ & 1.24$\times$ & [1.158, 1.324] & 0.0004 & 0.0008 & -2.13 & 0.0502 & 9.4 \\
schubfach:go-vs-rs & e2e & canonicalize & escaped-key-order & schubfach & schubfach-rs & schubfach-rs$^*$ & 2.82$\times$ & [2.537, 3.108] & 0.0002 & 0.0004 & 5.28 & 0.0680 & 12.8 \\
bd:go-vs-rs & e2e & canonicalize & escaped-key-order & json-canon & json-canon-rs & json-canon-rs$^*$ & 2.64$\times$ & [2.256, 3.747] & 0.0002 & 0.0004 & 1.89 & 0.0713 & 10.7 \\
ryu:go-vs-rs & e2e & canonicalize & escaped-key-order & ryu & ryu-rs & ryu-rs$^*$ & 3.03$\times$ & [2.744, 3.369] & 0.0002 & 0.0004 & 5.90 & 0.0728 & 13.9 \\
dragonbox:go-vs-rs & e2e & canonicalize & escaped-key-order & dragonbox & dragonbox-rs & dragonbox-rs$^*$ & 2.71$\times$ & [2.406, 3.021] & 0.0002 & 0.0004 & 4.04 & 0.0502 & 9.4 \\
go:schubfach-vs-bd & e2e & canonicalize & large & schubfach & json-canon & schubfach$^*$ & 1.55$\times$ & [1.500, 1.594] & 0.0002 & 0.0004 & -9.17 & 0.6642 & 4.2 \\
go:schubfach-vs-ryu & e2e & canonicalize & large & schubfach & ryu & ryu & 1.02$\times$ & [0.990, 1.053] & 0.2320 & 0.3280 & 0.45 & 0.7737 & 4.9 \\
go:schubfach-vs-dragonbox & e2e & canonicalize & large & schubfach & dragonbox & dragonbox & 1.03$\times$ & [0.995, 1.059] & 0.1206 & 0.1903 & 0.59 & 0.7887 & 5.1 \\
go:ryu-vs-dragonbox & e2e & canonicalize & large & ryu & dragonbox & dragonbox & 1.01$\times$ & [0.970, 1.040] & 0.6911 & 0.7643 & 0.14 & 0.7887 & 5.1 \\
go:ryu-vs-bd & e2e & canonicalize & large & ryu & json-canon & ryu$^*$ & 1.58$\times$ & [1.528, 1.632] & 0.0002 & 0.0004 & -9.13 & 0.7737 & 4.9 \\
go:dragonbox-vs-bd & e2e & canonicalize & large & dragonbox & json-canon & dragonbox$^*$ & 1.59$\times$ & [1.536, 1.642] & 0.0002 & 0.0004 & -9.19 & 0.7887 & 5.1 \\
rs:schubfach-vs-bd & e2e & canonicalize & large & schubfach-rs & json-canon-rs & schubfach-rs$^*$ & 1.96$\times$ & [1.842, 2.072] & 0.0002 & 0.0004 & -9.98 & 0.8980 & 11.6 \\
rs:schubfach-vs-ryu & e2e & canonicalize & large & schubfach-rs & ryu-rs & ryu-rs & 1.01$\times$ & [0.931, 1.088] & 0.8348 & 0.8790 & 0.08 & 0.8828 & 11.5 \\
rs:schubfach-vs-dragonbox & e2e & canonicalize & large & schubfach-rs & dragonbox-rs & schubfach-rs & 1.02$\times$ & [0.935, 1.151] & 0.7315 & 0.7987 & -0.14 & 0.8980 & 11.6 \\
rs:ryu-vs-dragonbox & e2e & canonicalize & large & ryu-rs & dragonbox-rs & ryu-rs & 1.03$\times$ & [0.946, 1.152] & 0.6053 & 0.6985 & -0.20 & 0.8828 & 11.5 \\
rs:ryu-vs-bd & e2e & canonicalize & large & ryu-rs & json-canon-rs & ryu-rs$^*$ & 1.98$\times$ & [1.847, 2.079] & 0.0002 & 0.0004 & -10.19 & 0.8828 & 11.5 \\
rs:dragonbox-vs-bd & e2e & canonicalize & large & dragonbox-rs & json-canon-rs & dragonbox-rs$^*$ & 1.92$\times$ & [1.718, 2.052] & 0.0002 & 0.0004 & -7.11 & 1.3562 & 17.1 \\
schubfach:go-vs-rs & e2e & canonicalize & large & schubfach & schubfach-rs & schubfach-rs$^*$ & 2.06$\times$ & [1.938, 2.177] & 0.0002 & 0.0004 & 10.65 & 0.8980 & 11.6 \\
bd:go-vs-rs & e2e & canonicalize & large & json-canon & json-canon-rs & json-canon-rs$^*$ & 1.63$\times$ & [1.578, 1.676] & 0.0002 & 0.0004 & 10.23 & 0.5949 & 3.9 \\
ryu:go-vs-rs & e2e & canonicalize & large & ryu & ryu-rs & ryu-rs$^*$ & 2.04$\times$ & [1.899, 2.145] & 0.0002 & 0.0004 & 9.83 & 0.8828 & 11.5 \\
dragonbox:go-vs-rs & e2e & canonicalize & large & dragonbox & dragonbox-rs & dragonbox-rs$^*$ & 1.97$\times$ & [1.755, 2.102] & 0.0002 & 0.0004 & 7.05 & 1.3562 & 17.1 \\
go:schubfach-vs-bd & e2e & canonicalize & long-string & schubfach & json-canon & schubfach & 1.14$\times$ & [0.979, 1.340] & 0.1378 & 0.2135 & -0.56 & 0.3637 & 23.4 \\
go:schubfach-vs-ryu & e2e & canonicalize & long-string & schubfach & ryu & ryu & 1.00$\times$ & [0.861, 1.176] & 0.9738 & 0.9795 & 0.01 & 0.3311 & 21.3 \\
go:schubfach-vs-dragonbox & e2e & canonicalize & long-string & schubfach & dragonbox & schubfach & 1.09$\times$ & [0.966, 1.244] & 0.1940 & 0.2814 & -0.49 & 0.3637 & 23.4 \\
go:ryu-vs-dragonbox & e2e & canonicalize & long-string & ryu & dragonbox & ryu & 1.09$\times$ & [0.981, 1.246] & 0.1620 & 0.2451 & -0.53 & 0.3311 & 21.3 \\
go:ryu-vs-bd & e2e & canonicalize & long-string & ryu & json-canon & ryu & 1.14$\times$ & [0.993, 1.340] & 0.1136 & 0.1808 & -0.59 & 0.3311 & 21.3 \\
go:dragonbox-vs-bd & e2e & canonicalize & long-string & dragonbox & json-canon & dragonbox & 1.04$\times$ & [0.919, 1.198] & 0.5565 & 0.6534 & -0.22 & 0.2167 & 12.8 \\
rs:schubfach-vs-bd & e2e & canonicalize & long-string & schubfach-rs & json-canon-rs & schubfach-rs & 1.04$\times$ & [0.966, 1.108] & 0.2450 & 0.3449 & -0.43 & 0.0419 & 6.7 \\
rs:schubfach-vs-ryu & e2e & canonicalize & long-string & schubfach-rs & ryu-rs & ryu-rs & 1.02$\times$ & [0.984, 1.077] & 0.4059 & 0.5129 & 0.32 & 0.0395 & 6.4 \\
rs:schubfach-vs-dragonbox & e2e & canonicalize & long-string & schubfach-rs & dragonbox-rs & dragonbox-rs & 1.02$\times$ & [0.974, 1.064] & 0.4993 & 0.6005 & 0.25 & 0.0401 & 6.5 \\
rs:ryu-vs-dragonbox & e2e & canonicalize & long-string & ryu-rs & dragonbox-rs & ryu-rs & 1.00$\times$ & [0.967, 1.059] & 0.8624 & 0.9014 & -0.06 & 0.0395 & 6.4 \\
rs:ryu-vs-bd & e2e & canonicalize & long-string & ryu-rs & json-canon-rs & ryu-rs & 1.07$\times$ & [0.993, 1.136] & 0.0892 & 0.1459 & -0.64 & 0.0395 & 6.4 \\
rs:dragonbox-vs-bd & e2e & canonicalize & long-string & dragonbox-rs & json-canon-rs & dragonbox-rs & 1.06$\times$ & [0.982, 1.126] & 0.1218 & 0.1917 & -0.60 & 0.0401 & 6.5 \\
schubfach:go-vs-rs & e2e & canonicalize & long-string & schubfach & schubfach-rs & schubfach-rs$^*$ & 2.48$\times$ & [2.207, 2.771] & 0.0002 & 0.0004 & 3.67 & 0.0419 & 6.7 \\
bd:go-vs-rs & e2e & canonicalize & long-string & json-canon & json-canon-rs & json-canon-rs$^*$ & 2.70$\times$ & [2.392, 3.085] & 0.0002 & 0.0004 & 3.76 & 0.0827 & 12.6 \\
ryu:go-vs-rs & e2e & canonicalize & long-string & ryu & ryu-rs & ryu-rs$^*$ & 2.52$\times$ & [2.249, 2.779] & 0.0002 & 0.0004 & 4.06 & 0.0395 & 6.4 \\
dragonbox:go-vs-rs & e2e & canonicalize & long-string & dragonbox & dragonbox-rs & dragonbox-rs$^*$ & 2.75$\times$ & [2.548, 2.926] & 0.0002 & 0.0004 & 7.09 & 0.0401 & 6.5 \\
go:schubfach-vs-bd & e2e & canonicalize & medium & schubfach & json-canon & schubfach$^*$ & 1.22$\times$ & [1.070, 1.376] & 0.0058 & 0.0113 & -1.08 & 0.3476 & 16.5 \\
go:schubfach-vs-ryu & e2e & canonicalize & medium & schubfach & ryu & ryu & 1.09$\times$ & [0.956, 1.220] & 0.1744 & 0.2582 & 0.50 & 0.3712 & 19.2 \\
go:schubfach-vs-dragonbox & e2e & canonicalize & medium & schubfach & dragonbox & dragonbox & 1.07$\times$ & [0.942, 1.202] & 0.3135 & 0.4153 & 0.37 & 0.3909 & 19.8 \\
go:ryu-vs-dragonbox & e2e & canonicalize & medium & ryu & dragonbox & ryu & 1.02$\times$ & [0.888, 1.158] & 0.7720 & 0.8311 & -0.11 & 0.3712 & 19.2 \\
go:ryu-vs-bd & e2e & canonicalize & medium & ryu & json-canon & ryu$^*$ & 1.33$\times$ & [1.159, 1.503] & 0.0008 & 0.0017 & -1.46 & 0.3712 & 19.2 \\
go:dragonbox-vs-bd & e2e & canonicalize & medium & dragonbox & json-canon & dragonbox$^*$ & 1.30$\times$ & [1.127, 1.473] & 0.0020 & 0.0041 & -1.34 & 0.3909 & 19.8 \\
rs:schubfach-vs-bd & e2e & canonicalize & medium & schubfach-rs & json-canon-rs & schubfach-rs$^*$ & 1.58$\times$ & [1.471, 1.711] & 0.0002 & 0.0004 & -4.25 & 0.1019 & 12.0 \\
rs:schubfach-vs-ryu & e2e & canonicalize & medium & schubfach-rs & ryu-rs & ryu-rs & 1.10$\times$ & [1.001, 1.204] & 0.0702 & 0.1178 & 0.69 & 0.1198 & 15.4 \\
rs:schubfach-vs-dragonbox & e2e & canonicalize & medium & schubfach-rs & dragonbox-rs & dragonbox-rs & 1.05$\times$ & [0.963, 1.144] & 0.2699 & 0.3708 & 0.41 & 0.1062 & 13.1 \\
rs:ryu-vs-dragonbox & e2e & canonicalize & medium & ryu-rs & dragonbox-rs & ryu-rs & 1.04$\times$ & [0.950, 1.154] & 0.4163 & 0.5235 & -0.30 & 0.1198 & 15.4 \\
rs:ryu-vs-bd & e2e & canonicalize & medium & ryu-rs & json-canon-rs & ryu-rs$^*$ & 1.74$\times$ & [1.589, 1.902] & 0.0002 & 0.0004 & -4.59 & 0.1198 & 15.4 \\
rs:dragonbox-vs-bd & e2e & canonicalize & medium & dragonbox-rs & json-canon-rs & dragonbox-rs$^*$ & 1.66$\times$ & [1.538, 1.798] & 0.0002 & 0.0004 & -4.54 & 0.1062 & 13.1 \\
schubfach:go-vs-rs & e2e & canonicalize & medium & schubfach & schubfach-rs & schubfach-rs$^*$ & 2.48$\times$ & [2.240, 2.732] & 0.0002 & 0.0004 & 5.02 & 0.1019 & 12.0 \\
bd:go-vs-rs & e2e & canonicalize & medium & json-canon & json-canon-rs & json-canon-rs$^*$ & 1.91$\times$ & [1.700, 2.098] & 0.0002 & 0.0004 & 3.36 & 0.1344 & 10.0 \\
ryu:go-vs-rs & e2e & canonicalize & medium & ryu & ryu-rs & ryu-rs$^*$ & 2.49$\times$ & [2.228, 2.820] & 0.0002 & 0.0004 & 4.29 & 0.1198 & 15.4 \\
dragonbox:go-vs-rs & e2e & canonicalize & medium & dragonbox & dragonbox-rs & dragonbox-rs$^*$ & 2.44$\times$ & [2.172, 2.727] & 0.0002 & 0.0004 & 4.16 & 0.1062 & 13.1 \\
go:schubfach-vs-bd & e2e & canonicalize & mixed-prod & schubfach & json-canon & schubfach$^*$ & 1.18$\times$ & [1.041, 1.340] & 0.0182 & 0.0339 & -0.91 & 0.3032 & 18.6 \\
go:schubfach-vs-ryu & e2e & canonicalize & mixed-prod & schubfach & ryu & schubfach & 1.06$\times$ & [0.946, 1.192] & 0.3195 & 0.4214 & -0.37 & 0.3032 & 18.6 \\
go:schubfach-vs-dragonbox & e2e & canonicalize & mixed-prod & schubfach & dragonbox & dragonbox & 1.00$\times$ & [0.878, 1.154] & 0.9408 & 0.9601 & 0.02 & 0.3383 & 20.8 \\
go:ryu-vs-dragonbox & e2e & canonicalize & mixed-prod & ryu & dragonbox & dragonbox & 1.07$\times$ & [0.945, 1.210] & 0.3197 & 0.4214 & 0.37 & 0.3383 & 20.8 \\
go:ryu-vs-bd & e2e & canonicalize & mixed-prod & ryu & json-canon & ryu & 1.11$\times$ & [0.988, 1.249] & 0.1158 & 0.1838 & -0.60 & 0.2801 & 16.1 \\
go:dragonbox-vs-bd & e2e & canonicalize & mixed-prod & dragonbox & json-canon & dragonbox$^*$ & 1.18$\times$ & [1.042, 1.359] & 0.0208 & 0.0384 & -0.89 & 0.3383 & 20.8 \\
rs:schubfach-vs-bd & e2e & canonicalize & mixed-prod & schubfach-rs & json-canon-rs & schubfach-rs$^*$ & 1.21$\times$ & [1.115, 1.310] & 0.0004 & 0.0008 & -1.65 & 0.0784 & 11.2 \\
rs:schubfach-vs-ryu & e2e & canonicalize & mixed-prod & schubfach-rs & ryu-rs & ryu-rs$^*$ & 1.11$\times$ & [1.027, 1.209] & 0.0176 & 0.0329 & 0.90 & 0.0804 & 12.7 \\
rs:schubfach-vs-dragonbox & e2e & canonicalize & mixed-prod & schubfach-rs & dragonbox-rs & dragonbox-rs & 1.05$\times$ & [0.986, 1.125] & 0.1590 & 0.2412 & 0.54 & 0.0542 & 8.1 \\
rs:ryu-vs-dragonbox & e2e & canonicalize & mixed-prod & ryu-rs & dragonbox-rs & ryu-rs & 1.05$\times$ & [0.985, 1.144] & 0.1704 & 0.2534 & -0.51 & 0.0804 & 12.7 \\
rs:ryu-vs-bd & e2e & canonicalize & mixed-prod & ryu-rs & json-canon-rs & ryu-rs$^*$ & 1.35$\times$ & [1.239, 1.468] & 0.0002 & 0.0004 & -2.39 & 0.0804 & 12.7 \\
rs:dragonbox-vs-bd & e2e & canonicalize & mixed-prod & dragonbox-rs & json-canon-rs & dragonbox-rs$^*$ & 1.28$\times$ & [1.180, 1.365] & 0.0002 & 0.0004 & -2.25 & 0.0542 & 8.1 \\
schubfach:go-vs-rs & e2e & canonicalize & mixed-prod & schubfach & schubfach-rs & schubfach-rs$^*$ & 2.32$\times$ & [2.107, 2.585] & 0.0002 & 0.0004 & 4.29 & 0.0784 & 11.2 \\
bd:go-vs-rs & e2e & canonicalize & mixed-prod & json-canon & json-canon-rs & json-canon-rs$^*$ & 2.25$\times$ & [2.038, 2.516] & 0.0002 & 0.0004 & 4.24 & 0.1061 & 12.4 \\
ryu:go-vs-rs & e2e & canonicalize & mixed-prod & ryu & ryu-rs & ryu-rs$^*$ & 2.74$\times$ & [2.501, 3.045] & 0.0002 & 0.0004 & 5.47 & 0.0804 & 12.7 \\
dragonbox:go-vs-rs & e2e & canonicalize & mixed-prod & dragonbox & dragonbox-rs & dragonbox-rs$^*$ & 2.43$\times$ & [2.183, 2.705] & 0.0002 & 0.0004 & 4.03 & 0.0542 & 8.1 \\
go:schubfach-vs-bd & e2e & canonicalize & nested-mixed & schubfach & json-canon & schubfach & 1.02$\times$ & [0.914, 1.148] & 0.6707 & 0.7507 & -0.15 & 0.2108 & 13.2 \\
go:schubfach-vs-ryu & e2e & canonicalize & nested-mixed & schubfach & ryu & ryu & 1.11$\times$ & [1.005, 1.246] & 0.0778 & 0.1290 & 0.69 & 0.2533 & 17.6 \\
go:schubfach-vs-dragonbox & e2e & canonicalize & nested-mixed & schubfach & dragonbox & dragonbox & 1.08$\times$ & [0.965, 1.203] & 0.1984 & 0.2871 & 0.48 & 0.2768 & 18.7 \\
go:ryu-vs-dragonbox & e2e & canonicalize & nested-mixed & ryu & dragonbox & ryu & 1.03$\times$ & [0.918, 1.172] & 0.6609 & 0.7443 & -0.16 & 0.2533 & 17.6 \\
go:ryu-vs-bd & e2e & canonicalize & nested-mixed & ryu & json-canon & ryu & 1.14$\times$ & [1.007, 1.288] & 0.0684 & 0.1154 & -0.71 & 0.2533 & 17.6 \\
go:dragonbox-vs-bd & e2e & canonicalize & nested-mixed & dragonbox & json-canon & dragonbox & 1.10$\times$ & [0.968, 1.245] & 0.1476 & 0.2273 & -0.54 & 0.2768 & 18.7 \\
rs:schubfach-vs-bd & e2e & canonicalize & nested-mixed & schubfach-rs & json-canon-rs & schubfach-rs$^*$ & 1.18$\times$ & [1.070, 1.295] & 0.0030 & 0.0061 & -1.18 & 0.0753 & 13.8 \\
rs:schubfach-vs-ryu & e2e & canonicalize & nested-mixed & schubfach-rs & ryu-rs & ryu-rs & 1.02$\times$ & [0.924, 1.130] & 0.6853 & 0.7614 & 0.15 & 0.0832 & 15.6 \\
rs:schubfach-vs-dragonbox & e2e & canonicalize & nested-mixed & schubfach-rs & dragonbox-rs & schubfach-rs & 1.01$\times$ & [0.953, 1.094] & 0.8548 & 0.8949 & -0.08 & 0.0753 & 13.8 \\
rs:ryu-vs-dragonbox & e2e & canonicalize & nested-mixed & ryu-rs & dragonbox-rs & ryu-rs & 1.03$\times$ & [0.964, 1.118] & 0.4629 & 0.5705 & -0.27 & 0.0832 & 15.6 \\
rs:ryu-vs-bd & e2e & canonicalize & nested-mixed & ryu-rs & json-canon-rs & ryu-rs$^*$ & 1.20$\times$ & [1.086, 1.326] & 0.0018 & 0.0037 & -1.27 & 0.0832 & 15.6 \\
rs:dragonbox-vs-bd & e2e & canonicalize & nested-mixed & dragonbox-rs & json-canon-rs & dragonbox-rs$^*$ & 1.17$\times$ & [1.071, 1.233] & 0.0006 & 0.0013 & -1.46 & 0.0150 & 2.7 \\
schubfach:go-vs-rs & e2e & canonicalize & nested-mixed & schubfach & schubfach-rs & schubfach-rs$^*$ & 2.92$\times$ & [2.668, 3.209] & 0.0002 & 0.0004 & 6.77 & 0.0753 & 13.8 \\
bd:go-vs-rs & e2e & canonicalize & nested-mixed & json-canon & json-canon-rs & json-canon-rs$^*$ & 2.55$\times$ & [2.257, 2.836] & 0.0002 & 0.0004 & 4.45 & 0.0897 & 14.0 \\
ryu:go-vs-rs & e2e & canonicalize & nested-mixed & ryu & ryu-rs & ryu-rs$^*$ & 2.69$\times$ & [2.385, 2.993] & 0.0002 & 0.0004 & 4.90 & 0.0832 & 15.6 \\
dragonbox:go-vs-rs & e2e & canonicalize & nested-mixed & dragonbox & dragonbox-rs & dragonbox-rs$^*$ & 2.69$\times$ & [2.437, 2.924] & 0.0002 & 0.0004 & 4.85 & 0.0150 & 2.7 \\
go:schubfach-vs-bd & e2e & canonicalize & number-heavy & schubfach & json-canon & schubfach & 1.11$\times$ & [0.998, 1.237] & 0.0934 & 0.1512 & -0.64 & 0.2289 & 15.7 \\
go:schubfach-vs-ryu & e2e & canonicalize & number-heavy & schubfach & ryu & schubfach & 1.04$\times$ & [0.919, 1.160] & 0.5237 & 0.6229 & -0.23 & 0.2289 & 15.7 \\
go:schubfach-vs-dragonbox & e2e & canonicalize & number-heavy & schubfach & dragonbox & schubfach & 1.06$\times$ & [0.949, 1.175] & 0.3169 & 0.4193 & -0.36 & 0.2289 & 15.7 \\
go:ryu-vs-dragonbox & e2e & canonicalize & number-heavy & ryu & dragonbox & ryu & 1.02$\times$ & [0.910, 1.147] & 0.7826 & 0.8391 & -0.10 & 0.2811 & 18.5 \\
go:ryu-vs-bd & e2e & canonicalize & number-heavy & ryu & json-canon & ryu & 1.06$\times$ & [0.951, 1.208] & 0.3211 & 0.4223 & -0.36 & 0.2811 & 18.5 \\
go:dragonbox-vs-bd & e2e & canonicalize & number-heavy & dragonbox & json-canon & dragonbox & 1.05$\times$ & [0.939, 1.167] & 0.4101 & 0.5172 & -0.30 & 0.2348 & 15.2 \\
rs:schubfach-vs-bd & e2e & canonicalize & number-heavy & schubfach-rs & json-canon-rs & schubfach-rs$^*$ & 1.34$\times$ & [1.242, 1.449] & 0.0002 & 0.0004 & -2.64 & 0.0593 & 10.9 \\
rs:schubfach-vs-ryu & e2e & canonicalize & number-heavy & schubfach-rs & ryu-rs & ryu-rs & 1.03$\times$ & [0.946, 1.101] & 0.4927 & 0.5953 & 0.26 & 0.0602 & 11.4 \\
rs:schubfach-vs-dragonbox & e2e & canonicalize & number-heavy & schubfach-rs & dragonbox-rs & schubfach-rs & 1.01$\times$ & [0.957, 1.094] & 0.6999 & 0.7701 & -0.15 & 0.0593 & 10.9 \\
rs:ryu-vs-dragonbox & e2e & canonicalize & number-heavy & ryu-rs & dragonbox-rs & ryu-rs & 1.04$\times$ & [0.977, 1.122] & 0.2567 & 0.3557 & -0.43 & 0.0602 & 11.4 \\
rs:ryu-vs-bd & e2e & canonicalize & number-heavy & ryu-rs & json-canon-rs & ryu-rs$^*$ & 1.38$\times$ & [1.273, 1.487] & 0.0002 & 0.0004 & -2.84 & 0.0602 & 11.4 \\
rs:dragonbox-vs-bd & e2e & canonicalize & number-heavy & dragonbox-rs & json-canon-rs & dragonbox-rs$^*$ & 1.32$\times$ & [1.233, 1.411] & 0.0002 & 0.0004 & -2.72 & 0.0463 & 8.4 \\
schubfach:go-vs-rs & e2e & canonicalize & number-heavy & schubfach & schubfach-rs & schubfach-rs$^*$ & 2.68$\times$ & [2.463, 2.956] & 0.0002 & 0.0004 & 5.61 & 0.0593 & 10.9 \\
bd:go-vs-rs & e2e & canonicalize & number-heavy & json-canon & json-canon-rs & json-canon-rs$^*$ & 2.22$\times$ & [2.018, 2.448] & 0.0002 & 0.0004 & 4.54 & 0.0824 & 11.3 \\
ryu:go-vs-rs & e2e & canonicalize & number-heavy & ryu & ryu-rs & ryu-rs$^*$ & 2.87$\times$ & [2.562, 3.166] & 0.0002 & 0.0004 & 4.98 & 0.0602 & 11.4 \\
dragonbox:go-vs-rs & e2e & canonicalize & number-heavy & dragonbox & dragonbox-rs & dragonbox-rs$^*$ & 2.80$\times$ & [2.567, 3.041] & 0.0002 & 0.0004 & 5.99 & 0.0463 & 8.4 \\
go:schubfach-vs-bd & e2e & canonicalize & numeric-boundary & schubfach & json-canon & schubfach$^*$ & 1.17$\times$ & [1.028, 1.316] & 0.0262 & 0.0479 & -0.87 & 0.2519 & 17.3 \\
go:schubfach-vs-ryu & e2e & canonicalize & numeric-boundary & schubfach & ryu & schubfach & 1.07$\times$ & [0.957, 1.206] & 0.2368 & 0.3345 & -0.43 & 0.2519 & 17.3 \\
go:schubfach-vs-dragonbox & e2e & canonicalize & numeric-boundary & schubfach & dragonbox & dragonbox & 1.03$\times$ & [0.896, 1.184] & 0.7199 & 0.7887 & 0.13 & 0.3257 & 23.0 \\
go:ryu-vs-dragonbox & e2e & canonicalize & numeric-boundary & ryu & dragonbox & dragonbox & 1.10$\times$ & [0.962, 1.265] & 0.1706 & 0.2534 & 0.51 & 0.3257 & 23.0 \\
go:ryu-vs-bd & e2e & canonicalize & numeric-boundary & ryu & json-canon & ryu & 1.09$\times$ & [0.967, 1.225] & 0.1854 & 0.2707 & -0.49 & 0.2511 & 16.1 \\
go:dragonbox-vs-bd & e2e & canonicalize & numeric-boundary & dragonbox & json-canon & dragonbox$^*$ & 1.20$\times$ & [1.039, 1.382] & 0.0226 & 0.0417 & -0.89 & 0.3257 & 23.0 \\
rs:schubfach-vs-bd & e2e & canonicalize & numeric-boundary & schubfach-rs & json-canon-rs & schubfach-rs$^*$ & 1.28$\times$ & [1.187, 1.381] & 0.0002 & 0.0004 & -2.27 & 0.0718 & 13.3 \\
rs:schubfach-vs-ryu & e2e & canonicalize & numeric-boundary & schubfach-rs & ryu-rs & schubfach-rs & 1.01$\times$ & [0.913, 1.117] & 0.8406 & 0.8837 & -0.08 & 0.0718 & 13.3 \\
rs:schubfach-vs-dragonbox & e2e & canonicalize & numeric-boundary & schubfach-rs & dragonbox-rs & schubfach-rs & 1.05$\times$ & [0.971, 1.140] & 0.2741 & 0.3749 & -0.41 & 0.0718 & 13.3 \\
rs:ryu-vs-dragonbox & e2e & canonicalize & numeric-boundary & ryu-rs & dragonbox-rs & ryu-rs & 1.03$\times$ & [0.948, 1.141] & 0.4933 & 0.5955 & -0.26 & 0.0900 & 16.5 \\
rs:ryu-vs-bd & e2e & canonicalize & numeric-boundary & ryu-rs & json-canon-rs & ryu-rs$^*$ & 1.26$\times$ & [1.157, 1.386] & 0.0002 & 0.0004 & -1.89 & 0.0900 & 16.5 \\
rs:dragonbox-vs-bd & e2e & canonicalize & numeric-boundary & dragonbox-rs & json-canon-rs & dragonbox-rs$^*$ & 1.22$\times$ & [1.140, 1.291] & 0.0002 & 0.0004 & -2.18 & 0.0537 & 9.5 \\
schubfach:go-vs-rs & e2e & canonicalize & numeric-boundary & schubfach & schubfach-rs & schubfach-rs$^*$ & 2.69$\times$ & [2.427, 2.978] & 0.0002 & 0.0004 & 5.06 & 0.0718 & 13.3 \\
bd:go-vs-rs & e2e & canonicalize & numeric-boundary & json-canon & json-canon-rs & json-canon-rs$^*$ & 2.46$\times$ & [2.229, 2.706] & 0.0002 & 0.0004 & 4.52 & 0.0623 & 9.0 \\
ryu:go-vs-rs & e2e & canonicalize & numeric-boundary & ryu & ryu-rs & ryu-rs$^*$ & 2.86$\times$ & [2.554, 3.191] & 0.0002 & 0.0004 & 5.51 & 0.0900 & 16.5 \\
dragonbox:go-vs-rs & e2e & canonicalize & numeric-boundary & dragonbox & dragonbox-rs & dragonbox-rs$^*$ & 2.51$\times$ & [2.206, 2.803] & 0.0002 & 0.0004 & 3.74 & 0.0537 & 9.5 \\
go:schubfach-vs-bd & e2e & canonicalize & rfc-key-sorting & schubfach & json-canon & schubfach & 1.07$\times$ & [0.957, 1.208] & 0.2789 & 0.3784 & -0.40 & 0.2528 & 17.7 \\
go:schubfach-vs-ryu & e2e & canonicalize & rfc-key-sorting & schubfach & ryu & schubfach & 1.02$\times$ & [0.900, 1.147] & 0.7886 & 0.8427 & -0.10 & 0.2528 & 17.7 \\
go:schubfach-vs-dragonbox & e2e & canonicalize & rfc-key-sorting & schubfach & dragonbox & schubfach & 1.00$\times$ & [0.889, 1.128] & 0.9718 & 0.9790 & -0.01 & 0.2528 & 17.7 \\
go:ryu-vs-dragonbox & e2e & canonicalize & rfc-key-sorting & ryu & dragonbox & dragonbox & 1.01$\times$ & [0.895, 1.135] & 0.8158 & 0.8660 & 0.08 & 0.2433 & 17.0 \\
go:ryu-vs-bd & e2e & canonicalize & rfc-key-sorting & ryu & json-canon & ryu & 1.05$\times$ & [0.943, 1.191] & 0.4143 & 0.5215 & -0.30 & 0.2565 & 17.6 \\
go:dragonbox-vs-bd & e2e & canonicalize & rfc-key-sorting & dragonbox & json-canon & dragonbox & 1.07$\times$ & [0.954, 1.197] & 0.2879 & 0.3873 & -0.39 & 0.2433 & 17.0 \\
rs:schubfach-vs-bd & e2e & canonicalize & rfc-key-sorting & schubfach-rs & json-canon-rs & schubfach-rs$^*$ & 1.12$\times$ & [1.037, 1.194] & 0.0040 & 0.0079 & -1.08 & 0.0530 & 9.9 \\
rs:schubfach-vs-ryu & e2e & canonicalize & rfc-key-sorting & schubfach-rs & ryu-rs & schubfach-rs & 1.09$\times$ & [1.011, 1.164] & 0.0320 & 0.0576 & -0.81 & 0.0530 & 9.9 \\
rs:schubfach-vs-dragonbox & e2e & canonicalize & rfc-key-sorting & schubfach-rs & dragonbox-rs & dragonbox-rs & 1.00$\times$ & [0.925, 1.111] & 0.9384 & 0.9592 & 0.03 & 0.0820 & 15.4 \\
rs:ryu-vs-dragonbox & e2e & canonicalize & rfc-key-sorting & ryu-rs & dragonbox-rs & dragonbox-rs & 1.09$\times$ & [1.002, 1.197] & 0.0774 & 0.1285 & 0.67 & 0.0820 & 15.4 \\
rs:ryu-vs-bd & e2e & canonicalize & rfc-key-sorting & ryu-rs & json-canon-rs & ryu-rs & 1.03$\times$ & [0.955, 1.102] & 0.4767 & 0.5807 & -0.26 & 0.0627 & 10.8 \\
rs:dragonbox-vs-bd & e2e & canonicalize & rfc-key-sorting & dragonbox-rs & json-canon-rs & dragonbox-rs$^*$ & 1.12$\times$ & [1.030, 1.238] & 0.0206 & 0.0381 & -0.89 & 0.0820 & 15.4 \\
schubfach:go-vs-rs & e2e & canonicalize & rfc-key-sorting & schubfach & schubfach-rs & schubfach-rs$^*$ & 2.67$\times$ & [2.430, 2.951] & 0.0002 & 0.0004 & 5.01 & 0.0530 & 9.9 \\
bd:go-vs-rs & e2e & canonicalize & rfc-key-sorting & json-canon & json-canon-rs & json-canon-rs$^*$ & 2.56$\times$ & [2.334, 2.813] & 0.0002 & 0.0004 & 5.14 & 0.0638 & 10.7 \\
ryu:go-vs-rs & e2e & canonicalize & rfc-key-sorting & ryu & ryu-rs & ryu-rs$^*$ & 2.50$\times$ & [2.246, 2.745] & 0.0002 & 0.0004 & 4.79 & 0.0627 & 10.8 \\
dragonbox:go-vs-rs & e2e & canonicalize & rfc-key-sorting & dragonbox & dragonbox-rs & dragonbox-rs$^*$ & 2.69$\times$ & [2.424, 3.011] & 0.0002 & 0.0004 & 5.07 & 0.0820 & 15.4 \\
go:schubfach-vs-bd & e2e & canonicalize & small & schubfach & json-canon & schubfach & 1.08$\times$ & [0.955, 1.218] & 0.2412 & 0.3399 & -0.44 & 0.2672 & 18.8 \\
go:schubfach-vs-ryu & e2e & canonicalize & small & schubfach & ryu & schubfach & 1.05$\times$ & [0.930, 1.180] & 0.4723 & 0.5766 & -0.26 & 0.2672 & 18.8 \\
go:schubfach-vs-dragonbox & e2e & canonicalize & small & schubfach & dragonbox & schubfach & 1.05$\times$ & [0.932, 1.175] & 0.4057 & 0.5129 & -0.30 & 0.2672 & 18.8 \\
go:ryu-vs-dragonbox & e2e & canonicalize & small & ryu & dragonbox & ryu & 1.01$\times$ & [0.902, 1.116] & 0.9336 & 0.9558 & -0.03 & 0.2472 & 16.6 \\
go:ryu-vs-bd & e2e & canonicalize & small & ryu & json-canon & ryu & 1.03$\times$ & [0.921, 1.161] & 0.6071 & 0.6994 & -0.19 & 0.2472 & 16.6 \\
go:dragonbox-vs-bd & e2e & canonicalize & small & dragonbox & json-canon & dragonbox & 1.03$\times$ & [0.924, 1.151] & 0.6673 & 0.7478 & -0.16 & 0.2278 & 15.3 \\
rs:schubfach-vs-bd & e2e & canonicalize & small & schubfach-rs & json-canon-rs & schubfach-rs$^*$ & 1.23$\times$ & [1.125, 1.336] & 0.0008 & 0.0017 & -1.66 & 0.0683 & 12.8 \\
rs:schubfach-vs-ryu & e2e & canonicalize & small & schubfach-rs & ryu-rs & schubfach-rs & 1.04$\times$ & [0.978, 1.130] & 0.2751 & 0.3752 & -0.40 & 0.0683 & 12.8 \\
rs:schubfach-vs-dragonbox & e2e & canonicalize & small & schubfach-rs & dragonbox-rs & dragonbox-rs & 1.02$\times$ & [0.914, 1.120] & 0.7399 & 0.8052 & 0.12 & 0.0874 & 16.6 \\
rs:ryu-vs-dragonbox & e2e & canonicalize & small & ryu-rs & dragonbox-rs & dragonbox-rs & 1.06$\times$ & [0.961, 1.149] & 0.2154 & 0.3079 & 0.47 & 0.0874 & 16.6 \\
rs:ryu-vs-bd & e2e & canonicalize & small & ryu-rs & json-canon-rs & ryu-rs$^*$ & 1.18$\times$ & [1.092, 1.253] & 0.0004 & 0.0008 & -1.58 & 0.0407 & 7.3 \\
rs:dragonbox-vs-bd & e2e & canonicalize & small & dragonbox-rs & json-canon-rs & dragonbox-rs$^*$ & 1.25$\times$ & [1.119, 1.370] & 0.0006 & 0.0013 & -1.58 & 0.0874 & 16.6 \\
schubfach:go-vs-rs & e2e & canonicalize & small & schubfach & schubfach-rs & schubfach-rs$^*$ & 2.66$\times$ & [2.383, 2.958] & 0.0002 & 0.0004 & 4.64 & 0.0683 & 12.8 \\
bd:go-vs-rs & e2e & canonicalize & small & json-canon & json-canon-rs & json-canon-rs$^*$ & 2.33$\times$ & [2.107, 2.579] & 0.0002 & 0.0004 & 4.64 & 0.0814 & 12.4 \\
ryu:go-vs-rs & e2e & canonicalize & small & ryu & ryu-rs & ryu-rs$^*$ & 2.67$\times$ & [2.437, 2.880] & 0.0002 & 0.0004 & 5.36 & 0.0407 & 7.3 \\
dragonbox:go-vs-rs & e2e & canonicalize & small & dragonbox & dragonbox-rs & dragonbox-rs$^*$ & 2.84$\times$ & [2.522, 3.142] & 0.0002 & 0.0004 & 5.74 & 0.0874 & 16.6 \\
go:schubfach-vs-bd & e2e & canonicalize & surrogate-pair & schubfach & json-canon & schubfach & 1.05$\times$ & [0.954, 1.177] & 0.3617 & 0.4665 & -0.33 & 0.2586 & 17.9 \\
go:schubfach-vs-ryu & e2e & canonicalize & surrogate-pair & schubfach & ryu & ryu & 1.02$\times$ & [0.900, 1.166] & 0.7331 & 0.7994 & 0.12 & 0.2987 & 21.1 \\
go:schubfach-vs-dragonbox & e2e & canonicalize & surrogate-pair & schubfach & dragonbox & schubfach & 1.06$\times$ & [0.965, 1.193] & 0.2603 & 0.3599 & -0.42 & 0.2586 & 17.9 \\
go:ryu-vs-dragonbox & e2e & canonicalize & surrogate-pair & ryu & dragonbox & ryu & 1.09$\times$ & [0.974, 1.227] & 0.1764 & 0.2602 & -0.51 & 0.2987 & 21.1 \\
go:ryu-vs-bd & e2e & canonicalize & surrogate-pair & ryu & json-canon & ryu & 1.08$\times$ & [0.954, 1.216] & 0.2408 & 0.3397 & -0.43 & 0.2987 & 21.1 \\
go:dragonbox-vs-bd & e2e & canonicalize & surrogate-pair & dragonbox & json-canon & json-canon & 1.01$\times$ & [0.926, 1.099] & 0.8198 & 0.8696 & 0.08 & 0.2063 & 13.6 \\
rs:schubfach-vs-bd & e2e & canonicalize & surrogate-pair & schubfach-rs & json-canon-rs & schubfach-rs$^*$ & 1.13$\times$ & [1.061, 1.206] & 0.0002 & 0.0004 & -1.37 & 0.0486 & 8.8 \\
rs:schubfach-vs-ryu & e2e & canonicalize & surrogate-pair & schubfach-rs & ryu-rs & ryu-rs & 1.04$\times$ & [0.973, 1.102] & 0.2563 & 0.3555 & 0.43 & 0.0495 & 9.3 \\
rs:schubfach-vs-dragonbox & e2e & canonicalize & surrogate-pair & schubfach-rs & dragonbox-rs & schubfach-rs & 1.00$\times$ & [0.930, 1.076] & 0.9566 & 0.9707 & -0.02 & 0.0486 & 8.8 \\
rs:ryu-vs-dragonbox & e2e & canonicalize & surrogate-pair & ryu-rs & dragonbox-rs & ryu-rs & 1.04$\times$ & [0.967, 1.121] & 0.3229 & 0.4242 & -0.37 & 0.0495 & 9.3 \\
rs:ryu-vs-bd & e2e & canonicalize & surrogate-pair & ryu-rs & json-canon-rs & ryu-rs$^*$ & 1.18$\times$ & [1.102, 1.251] & 0.0002 & 0.0004 & -1.75 & 0.0495 & 9.3 \\
rs:dragonbox-vs-bd & e2e & canonicalize & surrogate-pair & dragonbox-rs & json-canon-rs & dragonbox-rs$^*$ & 1.13$\times$ & [1.046, 1.216] & 0.0030 & 0.0061 & -1.14 & 0.0685 & 12.4 \\
schubfach:go-vs-rs & e2e & canonicalize & surrogate-pair & schubfach & schubfach-rs & schubfach-rs$^*$ & 2.61$\times$ & [2.345, 2.848] & 0.0002 & 0.0004 & 4.91 & 0.0486 & 8.8 \\
bd:go-vs-rs & e2e & canonicalize & surrogate-pair & json-canon & json-canon-rs & json-canon-rs$^*$ & 2.43$\times$ & [2.239, 2.620] & 0.0002 & 0.0004 & 6.03 & 0.0595 & 9.5 \\
ryu:go-vs-rs & e2e & canonicalize & surrogate-pair & ryu & ryu-rs & ryu-rs$^*$ & 2.65$\times$ & [2.358, 2.943] & 0.0002 & 0.0004 & 4.21 & 0.0495 & 9.3 \\
dragonbox:go-vs-rs & e2e & canonicalize & surrogate-pair & dragonbox & dragonbox-rs & dragonbox-rs$^*$ & 2.77$\times$ & [2.540, 2.997] & 0.0002 & 0.0004 & 7.31 & 0.0685 & 12.4 \\
go:schubfach-vs-bd & e2e & canonicalize & unicode & schubfach & json-canon & schubfach & 1.10$\times$ & [0.983, 1.219] & 0.1160 & 0.1839 & -0.60 & 0.2324 & 16.0 \\
go:schubfach-vs-ryu & e2e & canonicalize & unicode & schubfach & ryu & schubfach & 1.01$\times$ & [0.898, 1.132] & 0.8220 & 0.8707 & -0.08 & 0.2324 & 16.0 \\
go:schubfach-vs-dragonbox & e2e & canonicalize & unicode & schubfach & dragonbox & schubfach & 1.00$\times$ & [0.899, 1.118] & 0.9328 & 0.9557 & -0.03 & 0.2324 & 16.0 \\
go:ryu-vs-dragonbox & e2e & canonicalize & unicode & ryu & dragonbox & dragonbox & 1.01$\times$ & [0.898, 1.132] & 0.8834 & 0.9167 & 0.05 & 0.2401 & 16.5 \\
go:ryu-vs-bd & e2e & canonicalize & unicode & ryu & json-canon & ryu & 1.08$\times$ & [0.967, 1.216] & 0.1944 & 0.2816 & -0.49 & 0.2578 & 17.5 \\
go:dragonbox-vs-bd & e2e & canonicalize & unicode & dragonbox & json-canon & dragonbox & 1.09$\times$ & [0.982, 1.215] & 0.1424 & 0.2204 & -0.56 & 0.2401 & 16.5 \\
rs:schubfach-vs-bd & e2e & canonicalize & unicode & schubfach-rs & json-canon-rs & schubfach-rs$^*$ & 1.17$\times$ & [1.093, 1.254] & 0.0004 & 0.0008 & -1.52 & 0.0423 & 7.8 \\
rs:schubfach-vs-ryu & e2e & canonicalize & unicode & schubfach-rs & ryu-rs & schubfach-rs & 1.05$\times$ & [0.995, 1.106] & 0.1232 & 0.1930 & -0.59 & 0.0423 & 7.8 \\
rs:schubfach-vs-dragonbox & e2e & canonicalize & unicode & schubfach-rs & dragonbox-rs & schubfach-rs & 1.01$\times$ & [0.928, 1.106] & 0.8300 & 0.8772 & -0.08 & 0.0423 & 7.8 \\
rs:ryu-vs-dragonbox & e2e & canonicalize & unicode & ryu-rs & dragonbox-rs & dragonbox-rs & 1.03$\times$ & [0.949, 1.131] & 0.4577 & 0.5662 & 0.27 & 0.0899 & 16.5 \\
rs:ryu-vs-bd & e2e & canonicalize & unicode & ryu-rs & json-canon-rs & ryu-rs$^*$ & 1.11$\times$ & [1.043, 1.194] & 0.0054 & 0.0106 & -1.09 & 0.0431 & 7.7 \\
rs:dragonbox-vs-bd & e2e & canonicalize & unicode & dragonbox-rs & json-canon-rs & dragonbox-rs$^*$ & 1.15$\times$ & [1.051, 1.271] & 0.0078 & 0.0151 & -1.04 & 0.0899 & 16.5 \\
schubfach:go-vs-rs & e2e & canonicalize & unicode & schubfach & schubfach-rs & schubfach-rs$^*$ & 2.69$\times$ & [2.465, 2.924] & 0.0002 & 0.0004 & 5.57 & 0.0423 & 7.8 \\
bd:go-vs-rs & e2e & canonicalize & unicode & json-canon & json-canon-rs & json-canon-rs$^*$ & 2.53$\times$ & [2.287, 2.767] & 0.0002 & 0.0004 & 5.36 & 0.0738 & 11.7 \\
ryu:go-vs-rs & e2e & canonicalize & unicode & ryu & ryu-rs & ryu-rs$^*$ & 2.61$\times$ & [2.353, 2.827] & 0.0002 & 0.0004 & 5.01 & 0.0431 & 7.7 \\
dragonbox:go-vs-rs & e2e & canonicalize & unicode & dragonbox & dragonbox-rs & dragonbox-rs$^*$ & 2.67$\times$ & [2.371, 2.964] & 0.0002 & 0.0004 & 5.14 & 0.0899 & 16.5 \\
go:schubfach-vs-bd & e2e & canonicalize & verify-whitespace & schubfach & json-canon & schubfach & 1.06$\times$ & [0.939, 1.210] & 0.3623 & 0.4665 & -0.33 & 0.2703 & 17.9 \\
go:schubfach-vs-ryu & e2e & canonicalize & verify-whitespace & schubfach & ryu & schubfach & 1.01$\times$ & [0.897, 1.138] & 0.8316 & 0.8778 & -0.08 & 0.2703 & 17.9 \\
go:schubfach-vs-dragonbox & e2e & canonicalize & verify-whitespace & schubfach & dragonbox & dragonbox & 1.03$\times$ & [0.918, 1.163] & 0.5927 & 0.6878 & 0.20 & 0.2451 & 16.8 \\
go:ryu-vs-dragonbox & e2e & canonicalize & verify-whitespace & ryu & dragonbox & dragonbox & 1.05$\times$ & [0.937, 1.169] & 0.4461 & 0.5567 & 0.29 & 0.2451 & 16.8 \\
go:ryu-vs-bd & e2e & canonicalize & verify-whitespace & ryu & json-canon & ryu & 1.05$\times$ & [0.928, 1.195] & 0.4697 & 0.5756 & -0.26 & 0.2586 & 16.9 \\
go:dragonbox-vs-bd & e2e & canonicalize & verify-whitespace & dragonbox & json-canon & dragonbox & 1.10$\times$ & [0.972, 1.247] & 0.1706 & 0.2534 & -0.52 & 0.2451 & 16.8 \\
rs:schubfach-vs-bd & e2e & canonicalize & verify-whitespace & schubfach-rs & json-canon-rs & schubfach-rs$^*$ & 1.19$\times$ & [1.101, 1.298] & 0.0004 & 0.0008 & -1.45 & 0.0657 & 11.5 \\
rs:schubfach-vs-ryu & e2e & canonicalize & verify-whitespace & schubfach-rs & ryu-rs & ryu-rs & 1.05$\times$ & [0.966, 1.154] & 0.2953 & 0.3956 & 0.39 & 0.0783 & 14.4 \\
rs:schubfach-vs-dragonbox & e2e & canonicalize & verify-whitespace & schubfach-rs & dragonbox-rs & dragonbox-rs & 1.09$\times$ & [1.016, 1.182] & 0.0330 & 0.0591 & 0.79 & 0.0572 & 11.0 \\
rs:ryu-vs-dragonbox & e2e & canonicalize & verify-whitespace & ryu-rs & dragonbox-rs & dragonbox-rs & 1.04$\times$ & [0.949, 1.130] & 0.4173 & 0.5243 & 0.30 & 0.0572 & 11.0 \\
rs:ryu-vs-bd & e2e & canonicalize & verify-whitespace & ryu-rs & json-canon-rs & ryu-rs$^*$ & 1.26$\times$ & [1.150, 1.392] & 0.0002 & 0.0004 & -1.69 & 0.0783 & 14.4 \\
rs:dragonbox-vs-bd & e2e & canonicalize & verify-whitespace & dragonbox-rs & json-canon-rs & dragonbox-rs$^*$ & 1.30$\times$ & [1.209, 1.423] & 0.0002 & 0.0004 & -2.17 & 0.0572 & 11.0 \\
schubfach:go-vs-rs & e2e & canonicalize & verify-whitespace & schubfach & schubfach-rs & schubfach-rs$^*$ & 2.65$\times$ & [2.366, 2.911] & 0.0002 & 0.0004 & 4.89 & 0.0657 & 11.5 \\
bd:go-vs-rs & e2e & canonicalize & verify-whitespace & json-canon & json-canon-rs & json-canon-rs$^*$ & 2.36$\times$ & [2.091, 2.636] & 0.0002 & 0.0004 & 4.04 & 0.0891 & 13.1 \\
ryu:go-vs-rs & e2e & canonicalize & verify-whitespace & ryu & ryu-rs & ryu-rs$^*$ & 2.83$\times$ & [2.528, 3.131] & 0.0002 & 0.0004 & 5.30 & 0.0783 & 14.4 \\
dragonbox:go-vs-rs & e2e & canonicalize & verify-whitespace & dragonbox & dragonbox-rs & dragonbox-rs$^*$ & 2.80$\times$ & [2.523, 3.059] & 0.0002 & 0.0004 & 5.40 & 0.0572 & 11.0 \\
go:schubfach-vs-bd & e2e & verify & array-2048 & schubfach & json-canon & schubfach$^*$ & 1.31$\times$ & [1.221, 1.398] & 0.0002 & 0.0004 & -2.76 & 0.6324 & 11.5 \\
go:schubfach-vs-ryu & e2e & verify & array-2048 & schubfach & ryu & ryu & 1.00$\times$ & [0.922, 1.084] & 0.9448 & 0.9627 & 0.03 & 0.6812 & 12.5 \\
go:schubfach-vs-dragonbox & e2e & verify & array-2048 & schubfach & dragonbox & schubfach & 1.04$\times$ & [0.969, 1.111] & 0.2991 & 0.3995 & -0.39 & 0.6324 & 11.5 \\
go:ryu-vs-dragonbox & e2e & verify & array-2048 & ryu & dragonbox & ryu & 1.04$\times$ & [0.960, 1.110] & 0.2885 & 0.3877 & -0.40 & 0.6812 & 12.5 \\
go:ryu-vs-bd & e2e & verify & array-2048 & ryu & json-canon & ryu$^*$ & 1.31$\times$ & [1.215, 1.404] & 0.0002 & 0.0004 & -2.68 & 0.6812 & 12.5 \\
go:dragonbox-vs-bd & e2e & verify & array-2048 & dragonbox & json-canon & dragonbox$^*$ & 1.26$\times$ & [1.190, 1.338] & 0.0002 & 0.0004 & -2.70 & 0.4921 & 8.6 \\
rs:schubfach-vs-bd & e2e & verify & array-2048 & schubfach-rs & json-canon-rs & schubfach-rs$^*$ & 1.50$\times$ & [1.347, 1.713] & 0.0002 & 0.0004 & -2.33 & 0.4919 & 19.2 \\
rs:schubfach-vs-ryu & e2e & verify & array-2048 & schubfach-rs & ryu-rs & ryu-rs & 1.14$\times$ & [1.001, 1.271] & 0.0420 & 0.0733 & 0.77 & 0.3432 & 15.3 \\
rs:schubfach-vs-dragonbox & e2e & verify & array-2048 & schubfach-rs & dragonbox-rs & dragonbox-rs & 1.12$\times$ & [0.972, 1.263] & 0.1220 & 0.1918 & 0.58 & 0.4487 & 19.6 \\
rs:ryu-vs-dragonbox & e2e & verify & array-2048 & ryu-rs & dragonbox-rs & ryu-rs & 1.02$\times$ & [0.908, 1.155] & 0.7147 & 0.7839 & -0.13 & 0.3432 & 15.3 \\
rs:ryu-vs-bd & e2e & verify & array-2048 & ryu-rs & json-canon-rs & ryu-rs$^*$ & 1.71$\times$ & [1.556, 1.921] & 0.0002 & 0.0004 & -3.25 & 0.3432 & 15.3 \\
rs:dragonbox-vs-bd & e2e & verify & array-2048 & dragonbox-rs & json-canon-rs & dragonbox-rs$^*$ & 1.67$\times$ & [1.489, 1.903] & 0.0002 & 0.0004 & -2.91 & 0.4487 & 19.6 \\
schubfach:go-vs-rs & e2e & verify & array-2048 & schubfach & schubfach-rs & schubfach-rs$^*$ & 2.14$\times$ & [1.946, 2.412] & 0.0002 & 0.0004 & 5.28 & 0.4919 & 19.2 \\
bd:go-vs-rs & e2e & verify & array-2048 & json-canon & json-canon-rs & json-canon-rs$^*$ & 1.87$\times$ & [1.679, 2.011] & 0.0002 & 0.0004 & 5.52 & 0.6210 & 16.2 \\
ryu:go-vs-rs & e2e & verify & array-2048 & ryu & ryu-rs & ryu-rs$^*$ & 2.44$\times$ & [2.221, 2.682] & 0.0002 & 0.0004 & 6.12 & 0.3432 & 15.3 \\
dragonbox:go-vs-rs & e2e & verify & array-2048 & dragonbox & dragonbox-rs & dragonbox-rs$^*$ & 2.49$\times$ & [2.230, 2.741] & 0.0002 & 0.0004 & 7.40 & 0.4487 & 19.6 \\
go:schubfach-vs-bd & e2e & verify & array-256 & schubfach & json-canon & schubfach & 1.14$\times$ & [1.013, 1.281] & 0.0420 & 0.0733 & -0.77 & 0.3551 & 17.8 \\
go:schubfach-vs-ryu & e2e & verify & array-256 & schubfach & ryu & schubfach & 1.04$\times$ & [0.927, 1.151] & 0.5491 & 0.6476 & -0.22 & 0.3551 & 17.8 \\
go:schubfach-vs-dragonbox & e2e & verify & array-256 & schubfach & dragonbox & dragonbox & 1.10$\times$ & [0.955, 1.229] & 0.1780 & 0.2620 & 0.50 & 0.3518 & 19.3 \\
go:ryu-vs-dragonbox & e2e & verify & array-256 & ryu & dragonbox & dragonbox & 1.13$\times$ & [0.992, 1.258] & 0.0414 & 0.0725 & 0.77 & 0.3518 & 19.3 \\
go:ryu-vs-bd & e2e & verify & array-256 & ryu & json-canon & ryu & 1.10$\times$ & [0.994, 1.228] & 0.0956 & 0.1542 & -0.62 & 0.2934 & 14.2 \\
go:dragonbox-vs-bd & e2e & verify & array-256 & dragonbox & json-canon & dragonbox$^*$ & 1.25$\times$ & [1.092, 1.395] & 0.0018 & 0.0037 & -1.25 & 0.3518 & 19.3 \\
rs:schubfach-vs-bd & e2e & verify & array-256 & schubfach-rs & json-canon-rs & schubfach-rs$^*$ & 1.31$\times$ & [1.193, 1.447] & 0.0002 & 0.0004 & -1.87 & 0.0928 & 11.5 \\
rs:schubfach-vs-ryu & e2e & verify & array-256 & schubfach-rs & ryu-rs & ryu-rs & 1.02$\times$ & [0.943, 1.101] & 0.5641 & 0.6606 & 0.22 & 0.0892 & 11.3 \\
rs:schubfach-vs-dragonbox & e2e & verify & array-256 & schubfach-rs & dragonbox-rs & dragonbox-rs & 1.02$\times$ & [0.936, 1.101] & 0.7005 & 0.7701 & 0.14 & 0.0975 & 12.3 \\
rs:ryu-vs-dragonbox & e2e & verify & array-256 & ryu-rs & dragonbox-rs & ryu-rs & 1.01$\times$ & [0.928, 1.088] & 0.8744 & 0.9098 & -0.06 & 0.0892 & 11.3 \\
rs:ryu-vs-bd & e2e & verify & array-256 & ryu-rs & json-canon-rs & ryu-rs$^*$ & 1.34$\times$ & [1.219, 1.471] & 0.0002 & 0.0004 & -2.03 & 0.0892 & 11.3 \\
rs:dragonbox-vs-bd & e2e & verify & array-256 & dragonbox-rs & json-canon-rs & dragonbox-rs$^*$ & 1.33$\times$ & [1.209, 1.468] & 0.0004 & 0.0008 & -1.95 & 0.0975 & 12.3 \\
schubfach:go-vs-rs & e2e & verify & array-256 & schubfach & schubfach-rs & schubfach-rs$^*$ & 2.47$\times$ & [2.247, 2.756] & 0.0002 & 0.0004 & 4.68 & 0.0928 & 11.5 \\
bd:go-vs-rs & e2e & verify & array-256 & json-canon & json-canon-rs & json-canon-rs$^*$ & 2.15$\times$ & [1.939, 2.410] & 0.0002 & 0.0004 & 4.21 & 0.1679 & 15.9 \\
ryu:go-vs-rs & e2e & verify & array-256 & ryu & ryu-rs & ryu-rs$^*$ & 2.63$\times$ & [2.413, 2.876] & 0.0002 & 0.0004 & 6.03 & 0.0892 & 11.3 \\
dragonbox:go-vs-rs & e2e & verify & array-256 & dragonbox & dragonbox-rs & dragonbox-rs$^*$ & 2.30$\times$ & [2.092, 2.612] & 0.0002 & 0.0004 & 4.07 & 0.0975 & 12.3 \\
go:schubfach-vs-bd & e2e & verify & canonical-minimal & schubfach & json-canon & schubfach$^*$ & 1.13$\times$ & [1.041, 1.257] & 0.0178 & 0.0332 & -0.92 & 0.2287 & 14.9 \\
go:schubfach-vs-ryu & e2e & verify & canonical-minimal & schubfach & ryu & schubfach & 1.04$\times$ & [0.935, 1.154] & 0.5345 & 0.6334 & -0.23 & 0.2287 & 14.9 \\
go:schubfach-vs-dragonbox & e2e & verify & canonical-minimal & schubfach & dragonbox & schubfach & 1.02$\times$ & [0.938, 1.137] & 0.6221 & 0.7109 & -0.18 & 0.2287 & 14.9 \\
go:ryu-vs-dragonbox & e2e & verify & canonical-minimal & ryu & dragonbox & dragonbox & 1.01$\times$ & [0.911, 1.110] & 0.8378 & 0.8814 & 0.08 & 0.1990 & 12.7 \\
go:ryu-vs-bd & e2e & verify & canonical-minimal & ryu & json-canon & ryu & 1.09$\times$ & [0.991, 1.205] & 0.0924 & 0.1500 & -0.63 & 0.2567 & 16.2 \\
go:dragonbox-vs-bd & e2e & verify & canonical-minimal & dragonbox & json-canon & dragonbox & 1.10$\times$ & [1.012, 1.203] & 0.0362 & 0.0644 & -0.80 & 0.1990 & 12.7 \\
rs:schubfach-vs-bd & e2e & verify & canonical-minimal & schubfach-rs & json-canon-rs & schubfach-rs & 1.04$\times$ & [0.942, 1.128] & 0.4535 & 0.5621 & -0.28 & 0.0655 & 12.1 \\
rs:schubfach-vs-ryu & e2e & verify & canonical-minimal & schubfach-rs & ryu-rs & schubfach-rs & 1.02$\times$ & [0.944, 1.103] & 0.6367 & 0.7252 & -0.18 & 0.0655 & 12.1 \\
rs:schubfach-vs-dragonbox & e2e & verify & canonical-minimal & schubfach-rs & dragonbox-rs & schubfach-rs & 1.02$\times$ & [0.953, 1.093] & 0.6177 & 0.7084 & -0.18 & 0.0655 & 12.1 \\
rs:ryu-vs-dragonbox & e2e & verify & canonical-minimal & ryu-rs & dragonbox-rs & dragonbox-rs & 1.00$\times$ & [0.940, 1.070] & 0.9550 & 0.9707 & 0.02 & 0.0430 & 7.8 \\
rs:ryu-vs-bd & e2e & verify & canonical-minimal & ryu-rs & json-canon-rs & ryu-rs & 1.02$\times$ & [0.925, 1.105] & 0.7273 & 0.7962 & -0.13 & 0.0617 & 11.2 \\
rs:dragonbox-vs-bd & e2e & verify & canonical-minimal & dragonbox-rs & json-canon-rs & dragonbox-rs & 1.02$\times$ & [0.937, 1.096] & 0.6547 & 0.7399 & -0.16 & 0.0430 & 7.8 \\
schubfach:go-vs-rs & e2e & verify & canonical-minimal & schubfach & schubfach-rs & schubfach-rs$^*$ & 2.84$\times$ & [2.555, 3.073] & 0.0002 & 0.0004 & 6.04 & 0.0655 & 12.1 \\
bd:go-vs-rs & e2e & verify & canonical-minimal & json-canon & json-canon-rs & json-canon-rs$^*$ & 3.09$\times$ & [2.826, 3.398] & 0.0002 & 0.0004 & 7.32 & 0.0811 & 14.5 \\
ryu:go-vs-rs & e2e & verify & canonical-minimal & ryu & ryu-rs & ryu-rs$^*$ & 2.88$\times$ & [2.604, 3.144] & 0.0002 & 0.0004 & 5.68 & 0.0617 & 11.2 \\
dragonbox:go-vs-rs & e2e & verify & canonical-minimal & dragonbox & dragonbox-rs & dragonbox-rs$^*$ & 2.85$\times$ & [2.638, 3.051] & 0.0002 & 0.0004 & 7.26 & 0.0430 & 7.8 \\
go:schubfach-vs-bd & e2e & verify & control-escapes & schubfach & json-canon & schubfach & 1.01$\times$ & [0.907, 1.141] & 0.8986 & 0.9280 & -0.05 & 0.2775 & 18.1 \\
go:schubfach-vs-ryu & e2e & verify & control-escapes & schubfach & ryu & ryu & 1.05$\times$ & [0.918, 1.167] & 0.4593 & 0.5677 & 0.27 & 0.2425 & 16.5 \\
go:schubfach-vs-dragonbox & e2e & verify & control-escapes & schubfach & dragonbox & dragonbox & 1.11$\times$ & [0.967, 1.254] & 0.1494 & 0.2298 & 0.54 & 0.2838 & 20.5 \\
go:ryu-vs-dragonbox & e2e & verify & control-escapes & ryu & dragonbox & dragonbox & 1.06$\times$ & [0.938, 1.206] & 0.3957 & 0.5033 & 0.31 & 0.2838 & 20.5 \\
go:ryu-vs-bd & e2e & verify & control-escapes & ryu & json-canon & ryu & 1.06$\times$ & [0.952, 1.182] & 0.3385 & 0.4411 & -0.36 & 0.2425 & 16.5 \\
go:dragonbox-vs-bd & e2e & verify & control-escapes & dragonbox & json-canon & dragonbox & 1.12$\times$ & [0.990, 1.272] & 0.0904 & 0.1475 & -0.64 & 0.2838 & 20.5 \\
rs:schubfach-vs-bd & e2e & verify & control-escapes & schubfach-rs & json-canon-rs & schubfach-rs & 1.08$\times$ & [1.006, 1.184] & 0.0724 & 0.1210 & -0.69 & 0.0681 & 12.5 \\
rs:schubfach-vs-ryu & e2e & verify & control-escapes & schubfach-rs & ryu-rs & schubfach-rs & 1.01$\times$ & [0.949, 1.103] & 0.7646 & 0.8266 & -0.13 & 0.0681 & 12.5 \\
rs:schubfach-vs-dragonbox & e2e & verify & control-escapes & schubfach-rs & dragonbox-rs & schubfach-rs & 1.01$\times$ & [0.943, 1.113] & 0.7962 & 0.8487 & -0.10 & 0.0681 & 12.5 \\
rs:ryu-vs-dragonbox & e2e & verify & control-escapes & ryu-rs & dragonbox-rs & dragonbox-rs & 1.00$\times$ & [0.927, 1.065] & 0.9674 & 0.9772 & 0.01 & 0.0632 & 11.5 \\
rs:ryu-vs-bd & e2e & verify & control-escapes & ryu-rs & json-canon-rs & ryu-rs & 1.07$\times$ & [0.996, 1.141] & 0.0686 & 0.1156 & -0.68 & 0.0472 & 8.6 \\
rs:dragonbox-vs-bd & e2e & verify & control-escapes & dragonbox-rs & json-canon-rs & dragonbox-rs & 1.07$\times$ & [0.985, 1.148] & 0.1058 & 0.1690 & -0.62 & 0.0632 & 11.5 \\
schubfach:go-vs-rs & e2e & verify & control-escapes & schubfach & schubfach-rs & schubfach-rs$^*$ & 2.82$\times$ & [2.513, 3.115] & 0.0002 & 0.0004 & 5.01 & 0.0681 & 12.5 \\
bd:go-vs-rs & e2e & verify & control-escapes & json-canon & json-canon-rs & json-canon-rs$^*$ & 2.63$\times$ & [2.399, 2.872] & 0.0002 & 0.0004 & 5.79 & 0.0659 & 11.2 \\
ryu:go-vs-rs & e2e & verify & control-escapes & ryu & ryu-rs & ryu-rs$^*$ & 2.66$\times$ & [2.423, 2.901] & 0.0002 & 0.0004 & 5.35 & 0.0472 & 8.6 \\
dragonbox:go-vs-rs & e2e & verify & control-escapes & dragonbox & dragonbox-rs & dragonbox-rs$^*$ & 2.52$\times$ & [2.230, 2.787] & 0.0002 & 0.0004 & 4.15 & 0.0632 & 11.5 \\
go:schubfach-vs-bd & e2e & verify & deep-64 & schubfach & json-canon & json-canon & 1.01$\times$ & [0.872, 1.163] & 0.8512 & 0.8933 & 0.07 & 0.3102 & 20.0 \\
go:schubfach-vs-ryu & e2e & verify & deep-64 & schubfach & ryu & ryu & 1.05$\times$ & [0.907, 1.208] & 0.5305 & 0.6298 & 0.24 & 0.2957 & 19.7 \\
go:schubfach-vs-dragonbox & e2e & verify & deep-64 & schubfach & dragonbox & schubfach & 1.02$\times$ & [0.894, 1.173] & 0.7730 & 0.8315 & -0.11 & 0.3548 & 22.5 \\
go:ryu-vs-dragonbox & e2e & verify & deep-64 & ryu & dragonbox & ryu & 1.07$\times$ & [0.952, 1.220] & 0.2935 & 0.3936 & -0.40 & 0.2957 & 19.7 \\
go:ryu-vs-bd & e2e & verify & deep-64 & ryu & json-canon & ryu & 1.04$\times$ & [0.907, 1.188] & 0.6197 & 0.7096 & -0.18 & 0.2957 & 19.7 \\
go:dragonbox-vs-bd & e2e & verify & deep-64 & dragonbox & json-canon & json-canon & 1.03$\times$ & [0.913, 1.171] & 0.5925 & 0.6878 & 0.19 & 0.3102 & 20.0 \\
rs:schubfach-vs-bd & e2e & verify & deep-64 & schubfach-rs & json-canon-rs & schubfach-rs$^*$ & 1.08$\times$ & [1.025, 1.137] & 0.0074 & 0.0143 & -1.03 & 0.0346 & 6.0 \\
rs:schubfach-vs-ryu & e2e & verify & deep-64 & schubfach-rs & ryu-rs & ryu-rs & 1.02$\times$ & [0.972, 1.086] & 0.4283 & 0.5360 & 0.30 & 0.0559 & 9.9 \\
rs:schubfach-vs-dragonbox & e2e & verify & deep-64 & schubfach-rs & dragonbox-rs & dragonbox-rs & 1.02$\times$ & [0.969, 1.110] & 0.4825 & 0.5846 & 0.26 & 0.0686 & 12.1 \\
rs:ryu-vs-dragonbox & e2e & verify & deep-64 & ryu-rs & dragonbox-rs & dragonbox-rs & 1.00$\times$ & [0.935, 1.089] & 0.9926 & 0.9949 & 0.00 & 0.0686 & 12.1 \\
rs:ryu-vs-bd & e2e & verify & deep-64 & ryu-rs & json-canon-rs & ryu-rs$^*$ & 1.11$\times$ & [1.040, 1.183] & 0.0046 & 0.0090 & -1.11 & 0.0559 & 9.9 \\
rs:dragonbox-vs-bd & e2e & verify & deep-64 & dragonbox-rs & json-canon-rs & dragonbox-rs$^*$ & 1.11$\times$ & [1.038, 1.204] & 0.0100 & 0.0191 & -0.99 & 0.0686 & 12.1 \\
schubfach:go-vs-rs & e2e & verify & deep-64 & schubfach & schubfach-rs & schubfach-rs$^*$ & 2.72$\times$ & [2.414, 3.018] & 0.0002 & 0.0004 & 4.04 & 0.0346 & 6.0 \\
bd:go-vs-rs & e2e & verify & deep-64 & json-canon & json-canon-rs & json-canon-rs$^*$ & 2.48$\times$ & [2.208, 2.736] & 0.0002 & 0.0004 & 4.25 & 0.0566 & 9.0 \\
ryu:go-vs-rs & e2e & verify & deep-64 & ryu & ryu-rs & ryu-rs$^*$ & 2.65$\times$ & [2.368, 2.928] & 0.0002 & 0.0004 & 4.49 & 0.0559 & 9.9 \\
dragonbox:go-vs-rs & e2e & verify & deep-64 & dragonbox & dragonbox-rs & dragonbox-rs$^*$ & 2.85$\times$ & [2.573, 3.130] & 0.0002 & 0.0004 & 5.68 & 0.0686 & 12.1 \\
go:schubfach-vs-bd & e2e & verify & deep & schubfach & json-canon & schubfach & 1.07$\times$ & [0.960, 1.203] & 0.2857 & 0.3852 & -0.40 & 0.2449 & 15.1 \\
go:schubfach-vs-ryu & e2e & verify & deep & schubfach & ryu & ryu & 1.06$\times$ & [0.937, 1.185] & 0.3843 & 0.4924 & 0.32 & 0.3033 & 19.7 \\
go:schubfach-vs-dragonbox & e2e & verify & deep & schubfach & dragonbox & dragonbox & 1.03$\times$ & [0.913, 1.154] & 0.6225 & 0.7109 & 0.18 & 0.3065 & 19.4 \\
go:ryu-vs-dragonbox & e2e & verify & deep & ryu & dragonbox & ryu & 1.02$\times$ & [0.900, 1.173] & 0.7357 & 0.8013 & -0.13 & 0.3033 & 19.7 \\
go:ryu-vs-bd & e2e & verify & deep & ryu & json-canon & ryu & 1.13$\times$ & [0.992, 1.285] & 0.0848 & 0.1398 & -0.65 & 0.3033 & 19.7 \\
go:dragonbox-vs-bd & e2e & verify & deep & dragonbox & json-canon & dragonbox & 1.10$\times$ & [0.970, 1.250] & 0.1680 & 0.2516 & -0.52 & 0.3065 & 19.4 \\
rs:schubfach-vs-bd & e2e & verify & deep & schubfach-rs & json-canon-rs & schubfach-rs$^*$ & 1.14$\times$ & [1.064, 1.240] & 0.0034 & 0.0068 & -1.15 & 0.0785 & 12.8 \\
rs:schubfach-vs-ryu & e2e & verify & deep & schubfach-rs & ryu-rs & ryu-rs & 1.00$\times$ & [0.916, 1.065] & 0.9572 & 0.9707 & 0.02 & 0.0506 & 8.3 \\
rs:schubfach-vs-dragonbox & e2e & verify & deep & schubfach-rs & dragonbox-rs & schubfach-rs & 1.05$\times$ & [0.992, 1.147] & 0.2456 & 0.3454 & -0.45 & 0.0785 & 12.8 \\
rs:ryu-vs-dragonbox & e2e & verify & deep & ryu-rs & dragonbox-rs & ryu-rs & 1.05$\times$ & [1.003, 1.115] & 0.0982 & 0.1576 & -0.64 & 0.0506 & 8.3 \\
rs:ryu-vs-bd & e2e & verify & deep & ryu-rs & json-canon-rs & ryu-rs$^*$ & 1.14$\times$ & [1.069, 1.208] & 0.0006 & 0.0013 & -1.43 & 0.0506 & 8.3 \\
rs:dragonbox-vs-bd & e2e & verify & deep & dragonbox-rs & json-canon-rs & dragonbox-rs$^*$ & 1.09$\times$ & [1.015, 1.139] & 0.0156 & 0.0292 & -0.98 & 0.0433 & 6.8 \\
schubfach:go-vs-rs & e2e & verify & deep & schubfach & schubfach-rs & schubfach-rs$^*$ & 2.65$\times$ & [2.411, 2.921] & 0.0002 & 0.0004 & 5.70 & 0.0785 & 12.8 \\
bd:go-vs-rs & e2e & verify & deep & json-canon & json-canon-rs & json-canon-rs$^*$ & 2.50$\times$ & [2.259, 2.757] & 0.0002 & 0.0004 & 4.62 & 0.0686 & 9.9 \\
ryu:go-vs-rs & e2e & verify & deep & ryu & ryu-rs & ryu-rs$^*$ & 2.52$\times$ & [2.249, 2.769] & 0.0002 & 0.0004 & 4.36 & 0.0506 & 8.3 \\
dragonbox:go-vs-rs & e2e & verify & deep & dragonbox & dragonbox-rs & dragonbox-rs$^*$ & 2.46$\times$ & [2.234, 2.715] & 0.0002 & 0.0004 & 4.37 & 0.0433 & 6.8 \\
go:schubfach-vs-bd & e2e & verify & escaped-key-order & schubfach & json-canon & json-canon & 1.06$\times$ & [0.939, 1.190] & 0.3371 & 0.4398 & 0.35 & 0.2973 & 20.7 \\
go:schubfach-vs-ryu & e2e & verify & escaped-key-order & schubfach & ryu & ryu & 1.07$\times$ & [0.965, 1.189] & 0.2138 & 0.3060 & 0.46 & 0.2576 & 18.1 \\
go:schubfach-vs-dragonbox & e2e & verify & escaped-key-order & schubfach & dragonbox & dragonbox & 1.09$\times$ & [0.972, 1.200] & 0.1372 & 0.2128 & 0.57 & 0.2464 & 17.6 \\
go:ryu-vs-dragonbox & e2e & verify & escaped-key-order & ryu & dragonbox & dragonbox & 1.02$\times$ & [0.897, 1.142] & 0.8098 & 0.8613 & 0.09 & 0.2464 & 17.6 \\
go:ryu-vs-bd & e2e & verify & escaped-key-order & ryu & json-canon & ryu & 1.01$\times$ & [0.890, 1.152] & 0.8746 & 0.9098 & -0.06 & 0.2576 & 18.1 \\
go:dragonbox-vs-bd & e2e & verify & escaped-key-order & dragonbox & json-canon & dragonbox & 1.03$\times$ & [0.905, 1.171] & 0.6969 & 0.7681 & -0.14 & 0.2464 & 17.6 \\
rs:schubfach-vs-bd & e2e & verify & escaped-key-order & schubfach-rs & json-canon-rs & schubfach-rs$^*$ & 1.33$\times$ & [1.272, 1.444] & 0.0002 & 0.0004 & -3.79 & 0.0575 & 10.8 \\
rs:schubfach-vs-ryu & e2e & verify & escaped-key-order & schubfach-rs & ryu-rs & schubfach-rs & 1.09$\times$ & [0.991, 1.317] & 0.2482 & 0.3485 & -0.45 & 0.0575 & 10.8 \\
rs:schubfach-vs-dragonbox & e2e & verify & escaped-key-order & schubfach-rs & dragonbox-rs & schubfach-rs & 1.02$\times$ & [0.950, 1.105] & 0.6669 & 0.7478 & -0.16 & 0.0575 & 10.8 \\
rs:ryu-vs-dragonbox & e2e & verify & escaped-key-order & ryu-rs & dragonbox-rs & dragonbox-rs & 1.07$\times$ & [0.973, 1.283] & 0.3989 & 0.5063 & 0.36 & 0.0603 & 11.1 \\
rs:ryu-vs-bd & e2e & verify & escaped-key-order & ryu-rs & json-canon-rs & ryu-rs$^*$ & 1.22$\times$ & [1.016, 1.329] & 0.0008 & 0.0017 & -1.27 & 0.1411 & 24.3 \\
rs:dragonbox-vs-bd & e2e & verify & escaped-key-order & dragonbox-rs & json-canon-rs & dragonbox-rs$^*$ & 1.31$\times$ & [1.239, 1.394] & 0.0002 & 0.0004 & -3.46 & 0.0603 & 11.1 \\
schubfach:go-vs-rs & e2e & verify & escaped-key-order & schubfach & schubfach-rs & schubfach-rs$^*$ & 2.86$\times$ & [2.631, 3.116] & 0.0002 & 0.0004 & 6.75 & 0.0575 & 10.8 \\
bd:go-vs-rs & e2e & verify & escaped-key-order & json-canon & json-canon-rs & json-canon-rs$^*$ & 2.03$\times$ & [1.838, 2.254] & 0.0002 & 0.0004 & 3.52 & 0.0341 & 4.8 \\
ryu:go-vs-rs & e2e & verify & escaped-key-order & ryu & ryu-rs & ryu-rs$^*$ & 2.45$\times$ & [2.022, 2.740] & 0.0002 & 0.0004 & 4.14 & 0.1411 & 24.3 \\
dragonbox:go-vs-rs & e2e & verify & escaped-key-order & dragonbox & dragonbox-rs & dragonbox-rs$^*$ & 2.58$\times$ & [2.347, 2.870] & 0.0002 & 0.0004 & 4.88 & 0.0603 & 11.1 \\
go:schubfach-vs-bd & e2e & verify & large & schubfach & json-canon & schubfach$^*$ & 1.48$\times$ & [1.433, 1.519] & 0.0002 & 0.0004 & -10.93 & 0.8835 & 5.4 \\
go:schubfach-vs-ryu & e2e & verify & large & schubfach & ryu & ryu & 1.04$\times$ & [0.999, 1.084] & 0.0774 & 0.1285 & 0.69 & 1.0578 & 6.7 \\
go:schubfach-vs-dragonbox & e2e & verify & large & schubfach & dragonbox & dragonbox & 1.04$\times$ & [1.000, 1.088] & 0.0692 & 0.1164 & 0.68 & 1.1001 & 7.0 \\
go:ryu-vs-dragonbox & e2e & verify & large & ryu & dragonbox & dragonbox & 1.00$\times$ & [0.957, 1.050] & 0.9954 & 0.9969 & 0.00 & 1.1001 & 7.0 \\
go:ryu-vs-bd & e2e & verify & large & ryu & json-canon & ryu$^*$ & 1.55$\times$ & [1.487, 1.593] & 0.0002 & 0.0004 & -10.35 & 1.0578 & 6.7 \\
go:dragonbox-vs-bd & e2e & verify & large & dragonbox & json-canon & dragonbox$^*$ & 1.55$\times$ & [1.492, 1.601] & 0.0002 & 0.0004 & -10.04 & 1.1001 & 7.0 \\
rs:schubfach-vs-bd & e2e & verify & large & schubfach-rs & json-canon-rs & schubfach-rs$^*$ & 2.05$\times$ & [1.893, 2.171] & 0.0002 & 0.0004 & -9.13 & 1.0081 & 13.7 \\
rs:schubfach-vs-ryu & e2e & verify & large & schubfach-rs & ryu-rs & ryu-rs & 1.08$\times$ & [0.995, 1.177] & 0.0924 & 0.1500 & 0.63 & 0.7636 & 11.2 \\
rs:schubfach-vs-dragonbox & e2e & verify & large & schubfach-rs & dragonbox-rs & dragonbox-rs & 1.01$\times$ & [0.929, 1.102] & 0.7612 & 0.8250 & 0.11 & 0.8942 & 12.3 \\
rs:ryu-vs-dragonbox & e2e & verify & large & ryu-rs & dragonbox-rs & ryu-rs & 1.07$\times$ & [0.988, 1.156] & 0.1530 & 0.2340 & -0.55 & 0.7636 & 11.2 \\
rs:ryu-vs-bd & e2e & verify & large & ryu-rs & json-canon-rs & ryu-rs$^*$ & 2.21$\times$ & [2.044, 2.309] & 0.0002 & 0.0004 & -11.60 & 0.7636 & 11.2 \\
rs:dragonbox-vs-bd & e2e & verify & large & dragonbox-rs & json-canon-rs & dragonbox-rs$^*$ & 2.08$\times$ & [1.910, 2.178] & 0.0002 & 0.0004 & -10.00 & 0.8942 & 12.3 \\
schubfach:go-vs-rs & e2e & verify & large & schubfach & schubfach-rs & schubfach-rs$^*$ & 2.22$\times$ & [2.046, 2.357] & 0.0002 & 0.0004 & 9.71 & 1.0081 & 13.7 \\
bd:go-vs-rs & e2e & verify & large & json-canon & json-canon-rs & json-canon-rs$^*$ & 1.61$\times$ & [1.566, 1.645] & 0.0002 & 0.0004 & 14.86 & 0.6950 & 4.6 \\
ryu:go-vs-rs & e2e & verify & large & ryu & ryu-rs & ryu-rs$^*$ & 2.30$\times$ & [2.124, 2.425] & 0.0002 & 0.0004 & 9.86 & 0.7636 & 11.2 \\
dragonbox:go-vs-rs & e2e & verify & large & dragonbox & dragonbox-rs & dragonbox-rs$^*$ & 2.16$\times$ & [1.984, 2.281] & 0.0002 & 0.0004 & 8.61 & 0.8942 & 12.3 \\
go:schubfach-vs-bd & e2e & verify & long-string & schubfach & json-canon & schubfach & 1.14$\times$ & [0.987, 1.303] & 0.0874 & 0.1433 & -0.64 & 0.3676 & 23.9 \\
go:schubfach-vs-ryu & e2e & verify & long-string & schubfach & ryu & schubfach & 1.14$\times$ & [0.994, 1.290] & 0.0660 & 0.1116 & -0.69 & 0.3676 & 23.9 \\
go:schubfach-vs-dragonbox & e2e & verify & long-string & schubfach & dragonbox & schubfach & 1.11$\times$ & [0.964, 1.264] & 0.1350 & 0.2097 & -0.55 & 0.3676 & 23.9 \\
go:ryu-vs-dragonbox & e2e & verify & long-string & ryu & dragonbox & dragonbox & 1.02$\times$ & [0.924, 1.131] & 0.6537 & 0.7394 & 0.17 & 0.2634 & 15.4 \\
go:ryu-vs-bd & e2e & verify & long-string & ryu & json-canon & ryu & 1.00$\times$ & [0.894, 1.118] & 0.9994 & 0.9994 & -0.00 & 0.2541 & 14.5 \\
go:dragonbox-vs-bd & e2e & verify & long-string & dragonbox & json-canon & dragonbox & 1.02$\times$ & [0.920, 1.154] & 0.6821 & 0.7595 & -0.15 & 0.2634 & 15.4 \\
rs:schubfach-vs-bd & e2e & verify & long-string & schubfach-rs & json-canon-rs & schubfach-rs & 1.06$\times$ & [0.948, 1.163] & 0.2539 & 0.3532 & -0.43 & 0.1017 & 16.6 \\
rs:schubfach-vs-ryu & e2e & verify & long-string & schubfach-rs & ryu-rs & ryu-rs & 1.01$\times$ & [0.926, 1.132] & 0.9234 & 0.9498 & 0.04 & 0.0713 & 11.7 \\
rs:schubfach-vs-dragonbox & e2e & verify & long-string & schubfach-rs & dragonbox-rs & schubfach-rs & 1.00$\times$ & [0.897, 1.085] & 0.9652 & 0.9758 & -0.02 & 0.1017 & 16.6 \\
rs:ryu-vs-dragonbox & e2e & verify & long-string & ryu-rs & dragonbox-rs & ryu-rs & 1.01$\times$ & [0.940, 1.090] & 0.8538 & 0.8946 & -0.07 & 0.0713 & 11.7 \\
rs:ryu-vs-bd & e2e & verify & long-string & ryu-rs & json-canon-rs & ryu-rs & 1.07$\times$ & [0.986, 1.165] & 0.1336 & 0.2082 & -0.56 & 0.0713 & 11.7 \\
rs:dragonbox-vs-bd & e2e & verify & long-string & dragonbox-rs & json-canon-rs & dragonbox-rs & 1.06$\times$ & [0.986, 1.150] & 0.1642 & 0.2476 & -0.53 & 0.0613 & 10.0 \\
schubfach:go-vs-rs & e2e & verify & long-string & schubfach & schubfach-rs & schubfach-rs$^*$ & 2.51$\times$ & [2.175, 2.876] & 0.0002 & 0.0004 & 3.51 & 0.1017 & 16.6 \\
bd:go-vs-rs & e2e & verify & long-string & json-canon & json-canon-rs & json-canon-rs$^*$ & 2.69$\times$ & [2.416, 2.992] & 0.0002 & 0.0004 & 4.89 & 0.0830 & 12.7 \\
ryu:go-vs-rs & e2e & verify & long-string & ryu & ryu-rs & ryu-rs$^*$ & 2.87$\times$ & [2.612, 3.128] & 0.0002 & 0.0004 & 6.26 & 0.0713 & 11.7 \\
dragonbox:go-vs-rs & e2e & verify & long-string & dragonbox & dragonbox-rs & dragonbox-rs$^*$ & 2.78$\times$ & [2.538, 3.023] & 0.0002 & 0.0004 & 5.86 & 0.0613 & 10.0 \\
go:schubfach-vs-bd & e2e & verify & medium & schubfach & json-canon & schubfach & 1.13$\times$ & [0.995, 1.270] & 0.0718 & 0.1202 & -0.68 & 0.3754 & 18.4 \\
go:schubfach-vs-ryu & e2e & verify & medium & schubfach & ryu & schubfach & 1.06$\times$ & [0.937, 1.185] & 0.3345 & 0.4373 & -0.34 & 0.3754 & 18.4 \\
go:schubfach-vs-dragonbox & e2e & verify & medium & schubfach & dragonbox & dragonbox & 1.11$\times$ & [0.974, 1.243] & 0.1230 & 0.1929 & 0.58 & 0.3105 & 16.8 \\
go:ryu-vs-dragonbox & e2e & verify & medium & ryu & dragonbox & dragonbox$^*$ & 1.17$\times$ & [1.044, 1.303] & 0.0128 & 0.0242 & 0.98 & 0.3105 & 16.8 \\
go:ryu-vs-bd & e2e & verify & medium & ryu & json-canon & ryu & 1.06$\times$ & [0.953, 1.198] & 0.3201 & 0.4214 & -0.37 & 0.3504 & 16.2 \\
go:dragonbox-vs-bd & e2e & verify & medium & dragonbox & json-canon & dragonbox$^*$ & 1.24$\times$ & [1.110, 1.395] & 0.0024 & 0.0049 & -1.29 & 0.3105 & 16.8 \\
rs:schubfach-vs-bd & e2e & verify & medium & schubfach-rs & json-canon-rs & schubfach-rs$^*$ & 1.64$\times$ & [1.475, 1.814] & 0.0002 & 0.0004 & -3.32 & 0.1279 & 16.3 \\
rs:schubfach-vs-ryu & e2e & verify & medium & schubfach-rs & ryu-rs & schubfach-rs & 1.06$\times$ & [0.956, 1.173] & 0.3025 & 0.4034 & -0.38 & 0.1279 & 16.3 \\
rs:schubfach-vs-dragonbox & e2e & verify & medium & schubfach-rs & dragonbox-rs & schubfach-rs & 1.05$\times$ & [0.945, 1.152] & 0.3661 & 0.4709 & -0.33 & 0.1279 & 16.3 \\
rs:ryu-vs-dragonbox & e2e & verify & medium & ryu-rs & dragonbox-rs & dragonbox-rs & 1.01$\times$ & [0.928, 1.106] & 0.8328 & 0.8783 & 0.08 & 0.0984 & 12.0 \\
rs:ryu-vs-bd & e2e & verify & medium & ryu-rs & json-canon-rs & ryu-rs$^*$ & 1.55$\times$ & [1.403, 1.704] & 0.0002 & 0.0004 & -3.07 & 0.1206 & 14.5 \\
rs:dragonbox-vs-bd & e2e & verify & medium & dragonbox-rs & json-canon-rs & dragonbox-rs$^*$ & 1.57$\times$ & [1.434, 1.715] & 0.0002 & 0.0004 & -3.31 & 0.0984 & 12.0 \\
schubfach:go-vs-rs & e2e & verify & medium & schubfach & schubfach-rs & schubfach-rs$^*$ & 2.60$\times$ & [2.317, 2.928] & 0.0002 & 0.0004 & 4.59 & 0.1279 & 16.3 \\
bd:go-vs-rs & e2e & verify & medium & json-canon & json-canon-rs & json-canon-rs$^*$ & 1.79$\times$ & [1.615, 1.995] & 0.0002 & 0.0004 & 3.35 & 0.1774 & 13.8 \\
ryu:go-vs-rs & e2e & verify & medium & ryu & ryu-rs & ryu-rs$^*$ & 2.61$\times$ & [2.338, 2.889] & 0.0002 & 0.0004 & 5.20 & 0.1206 & 14.5 \\
dragonbox:go-vs-rs & e2e & verify & medium & dragonbox & dragonbox-rs & dragonbox-rs$^*$ & 2.25$\times$ & [2.059, 2.514] & 0.0002 & 0.0004 & 4.55 & 0.0984 & 12.0 \\
go:schubfach-vs-bd & e2e & verify & mixed-prod & schubfach & json-canon & schubfach & 1.05$\times$ & [0.936, 1.182] & 0.4447 & 0.5555 & -0.28 & 0.3485 & 19.0 \\
go:schubfach-vs-ryu & e2e & verify & mixed-prod & schubfach & ryu & ryu & 1.06$\times$ & [0.932, 1.209] & 0.3607 & 0.4659 & 0.34 & 0.3258 & 18.9 \\
go:schubfach-vs-dragonbox & e2e & verify & mixed-prod & schubfach & dragonbox & dragonbox & 1.07$\times$ & [0.947, 1.211] & 0.3125 & 0.4146 & 0.37 & 0.3040 & 17.7 \\
go:ryu-vs-dragonbox & e2e & verify & mixed-prod & ryu & dragonbox & dragonbox & 1.00$\times$ & [0.886, 1.133] & 0.9472 & 0.9644 & 0.03 & 0.3040 & 17.7 \\
go:ryu-vs-bd & e2e & verify & mixed-prod & ryu & json-canon & ryu & 1.12$\times$ & [1.000, 1.264] & 0.0894 & 0.1460 & -0.65 & 0.3258 & 18.9 \\
go:dragonbox-vs-bd & e2e & verify & mixed-prod & dragonbox & json-canon & dragonbox & 1.12$\times$ & [1.004, 1.252] & 0.0660 & 0.1116 & -0.70 & 0.3040 & 17.7 \\
rs:schubfach-vs-bd & e2e & verify & mixed-prod & schubfach-rs & json-canon-rs & schubfach-rs$^*$ & 1.24$\times$ & [1.125, 1.353] & 0.0002 & 0.0004 & -1.56 & 0.0799 & 12.0 \\
rs:schubfach-vs-ryu & e2e & verify & mixed-prod & schubfach-rs & ryu-rs & ryu-rs & 1.04$\times$ & [0.959, 1.126] & 0.3913 & 0.4994 & 0.31 & 0.0783 & 12.2 \\
rs:schubfach-vs-dragonbox & e2e & verify & mixed-prod & schubfach-rs & dragonbox-rs & schubfach-rs & 1.01$\times$ & [0.928, 1.090] & 0.8304 & 0.8772 & -0.08 & 0.0799 & 12.0 \\
rs:ryu-vs-dragonbox & e2e & verify & mixed-prod & ryu-rs & dragonbox-rs & ryu-rs & 1.05$\times$ & [0.963, 1.130] & 0.2805 & 0.3797 & -0.40 & 0.0783 & 12.2 \\
rs:ryu-vs-bd & e2e & verify & mixed-prod & ryu-rs & json-canon-rs & ryu-rs$^*$ & 1.28$\times$ & [1.166, 1.408] & 0.0004 & 0.0008 & -1.81 & 0.0783 & 12.2 \\
rs:dragonbox-vs-bd & e2e & verify & mixed-prod & dragonbox-rs & json-canon-rs & dragonbox-rs$^*$ & 1.23$\times$ & [1.114, 1.341] & 0.0004 & 0.0008 & -1.51 & 0.0781 & 11.7 \\
schubfach:go-vs-rs & e2e & verify & mixed-prod & schubfach & schubfach-rs & schubfach-rs$^*$ & 2.77$\times$ & [2.479, 3.062] & 0.0002 & 0.0004 & 4.75 & 0.0799 & 12.0 \\
bd:go-vs-rs & e2e & verify & mixed-prod & json-canon & json-canon-rs & json-canon-rs$^*$ & 2.34$\times$ & [2.131, 2.621] & 0.0002 & 0.0004 & 4.91 & 0.1228 & 14.9 \\
ryu:go-vs-rs & e2e & verify & mixed-prod & ryu & ryu-rs & ryu-rs$^*$ & 2.70$\times$ & [2.411, 2.991] & 0.0002 & 0.0004 & 4.69 & 0.0783 & 12.2 \\
dragonbox:go-vs-rs & e2e & verify & mixed-prod & dragonbox & dragonbox-rs & dragonbox-rs$^*$ & 2.56$\times$ & [2.311, 2.832] & 0.0002 & 0.0004 & 4.83 & 0.0781 & 11.7 \\
go:schubfach-vs-bd & e2e & verify & nested-mixed & schubfach & json-canon & schubfach & 1.19$\times$ & [1.022, 1.383] & 0.0364 & 0.0647 & -0.79 & 0.3175 & 23.3 \\
go:schubfach-vs-ryu & e2e & verify & nested-mixed & schubfach & ryu & schubfach & 1.16$\times$ & [1.008, 1.316] & 0.0472 & 0.0815 & -0.76 & 0.3175 & 23.3 \\
go:schubfach-vs-dragonbox & e2e & verify & nested-mixed & schubfach & dragonbox & schubfach & 1.13$\times$ & [0.990, 1.293] & 0.0872 & 0.1432 & -0.65 & 0.3175 & 23.3 \\
go:ryu-vs-dragonbox & e2e & verify & nested-mixed & ryu & dragonbox & dragonbox & 1.02$\times$ & [0.911, 1.144] & 0.7291 & 0.7975 & 0.13 & 0.2604 & 16.8 \\
go:ryu-vs-bd & e2e & verify & nested-mixed & ryu & json-canon & ryu & 1.03$\times$ & [0.902, 1.181] & 0.6859 & 0.7614 & -0.15 & 0.2650 & 16.8 \\
go:dragonbox-vs-bd & e2e & verify & nested-mixed & dragonbox & json-canon & dragonbox & 1.05$\times$ & [0.925, 1.205] & 0.4775 & 0.5807 & -0.26 & 0.2604 & 16.8 \\
rs:schubfach-vs-bd & e2e & verify & nested-mixed & schubfach-rs & json-canon-rs & schubfach-rs$^*$ & 1.30$\times$ & [1.202, 1.408] & 0.0002 & 0.0004 & -2.38 & 0.0750 & 13.8 \\
rs:schubfach-vs-ryu & e2e & verify & nested-mixed & schubfach-rs & ryu-rs & schubfach-rs & 1.01$\times$ & [0.933, 1.095] & 0.8426 & 0.8850 & -0.07 & 0.0750 & 13.8 \\
rs:schubfach-vs-dragonbox & e2e & verify & nested-mixed & schubfach-rs & dragonbox-rs & dragonbox-rs & 1.02$\times$ & [0.929, 1.136] & 0.7067 & 0.7762 & 0.14 & 0.0827 & 15.6 \\
rs:ryu-vs-dragonbox & e2e & verify & nested-mixed & ryu-rs & dragonbox-rs & dragonbox-rs & 1.03$\times$ & [0.950, 1.133] & 0.5581 & 0.6541 & 0.23 & 0.0827 & 15.6 \\
rs:ryu-vs-bd & e2e & verify & nested-mixed & ryu-rs & json-canon-rs & ryu-rs$^*$ & 1.28$\times$ & [1.217, 1.371] & 0.0002 & 0.0004 & -2.84 & 0.0496 & 9.1 \\
rs:dragonbox-vs-bd & e2e & verify & nested-mixed & dragonbox-rs & json-canon-rs & dragonbox-rs$^*$ & 1.32$\times$ & [1.229, 1.464] & 0.0002 & 0.0004 & -2.39 & 0.0827 & 15.6 \\
schubfach:go-vs-rs & e2e & verify & nested-mixed & schubfach & schubfach-rs & schubfach-rs$^*$ & 2.51$\times$ & [2.220, 2.856] & 0.0002 & 0.0004 & 3.63 & 0.0750 & 13.8 \\
bd:go-vs-rs & e2e & verify & nested-mixed & json-canon & json-canon-rs & json-canon-rs$^*$ & 2.32$\times$ & [2.057, 2.588] & 0.0002 & 0.0004 & 3.62 & 0.0620 & 8.8 \\
ryu:go-vs-rs & e2e & verify & nested-mixed & ryu & ryu-rs & ryu-rs$^*$ & 2.89$\times$ & [2.613, 3.150] & 0.0002 & 0.0004 & 5.54 & 0.0496 & 9.1 \\
dragonbox:go-vs-rs & e2e & verify & nested-mixed & dragonbox & dragonbox-rs & dragonbox-rs$^*$ & 2.91$\times$ & [2.624, 3.257] & 0.0002 & 0.0004 & 5.37 & 0.0827 & 15.6 \\
go:schubfach-vs-bd & e2e & verify & number-heavy & schubfach & json-canon & schubfach & 1.04$\times$ & [0.922, 1.209] & 0.5581 & 0.6541 & -0.21 & 0.2725 & 19.4 \\
go:schubfach-vs-ryu & e2e & verify & number-heavy & schubfach & ryu & schubfach & 1.09$\times$ & [0.951, 1.221] & 0.2060 & 0.2961 & -0.48 & 0.2725 & 19.4 \\
go:schubfach-vs-dragonbox & e2e & verify & number-heavy & schubfach & dragonbox & schubfach & 1.10$\times$ & [0.966, 1.234] & 0.1688 & 0.2522 & -0.51 & 0.2725 & 19.4 \\
go:ryu-vs-dragonbox & e2e & verify & number-heavy & ryu & dragonbox & ryu & 1.01$\times$ & [0.903, 1.134] & 0.9060 & 0.9342 & -0.04 & 0.2601 & 17.0 \\
go:ryu-vs-bd & e2e & verify & number-heavy & ryu & json-canon & json-canon & 1.04$\times$ & [0.904, 1.174] & 0.5345 & 0.6334 & 0.23 & 0.3012 & 20.5 \\
go:dragonbox-vs-bd & e2e & verify & number-heavy & dragonbox & json-canon & json-canon & 1.05$\times$ & [0.918, 1.188] & 0.4637 & 0.5709 & 0.27 & 0.3012 & 20.5 \\
rs:schubfach-vs-bd & e2e & verify & number-heavy & schubfach-rs & json-canon-rs & schubfach-rs$^*$ & 1.23$\times$ & [1.143, 1.320] & 0.0004 & 0.0008 & -1.96 & 0.0567 & 10.1 \\
rs:schubfach-vs-ryu & e2e & verify & number-heavy & schubfach-rs & ryu-rs & ryu-rs & 1.03$\times$ & [0.960, 1.121] & 0.4081 & 0.5152 & 0.31 & 0.0658 & 12.2 \\
rs:schubfach-vs-dragonbox & e2e & verify & number-heavy & schubfach-rs & dragonbox-rs & dragonbox-rs & 1.02$\times$ & [0.930, 1.117] & 0.6861 & 0.7614 & 0.15 & 0.0906 & 16.6 \\
rs:ryu-vs-dragonbox & e2e & verify & number-heavy & ryu-rs & dragonbox-rs & ryu-rs & 1.01$\times$ & [0.924, 1.122] & 0.7942 & 0.8473 & -0.10 & 0.0658 & 12.2 \\
rs:ryu-vs-bd & e2e & verify & number-heavy & ryu-rs & json-canon-rs & ryu-rs$^*$ & 1.27$\times$ & [1.182, 1.378] & 0.0002 & 0.0004 & -2.11 & 0.0658 & 12.2 \\
rs:dragonbox-vs-bd & e2e & verify & number-heavy & dragonbox-rs & json-canon-rs & dragonbox-rs$^*$ & 1.25$\times$ & [1.133, 1.373] & 0.0004 & 0.0008 & -1.70 & 0.0906 & 16.6 \\
schubfach:go-vs-rs & e2e & verify & number-heavy & schubfach & schubfach-rs & schubfach-rs$^*$ & 2.52$\times$ & [2.263, 2.791] & 0.0002 & 0.0004 & 4.40 & 0.0567 & 10.1 \\
bd:go-vs-rs & e2e & verify & number-heavy & json-canon & json-canon-rs & json-canon-rs$^*$ & 2.13$\times$ & [1.918, 2.403] & 0.0002 & 0.0004 & 3.63 & 0.0762 & 11.1 \\
ryu:go-vs-rs & e2e & verify & number-heavy & ryu & ryu-rs & ryu-rs$^*$ & 2.83$\times$ & [2.542, 3.116] & 0.0002 & 0.0004 & 5.33 & 0.0658 & 12.2 \\
dragonbox:go-vs-rs & e2e & verify & number-heavy & dragonbox & dragonbox-rs & dragonbox-rs$^*$ & 2.81$\times$ & [2.490, 3.134] & 0.0002 & 0.0004 & 5.07 & 0.0906 & 16.6 \\
go:schubfach-vs-bd & e2e & verify & numeric-boundary & schubfach & json-canon & schubfach & 1.02$\times$ & [0.924, 1.127] & 0.6593 & 0.7438 & -0.16 & 0.2155 & 14.0 \\
go:schubfach-vs-ryu & e2e & verify & numeric-boundary & schubfach & ryu & ryu & 1.06$\times$ & [0.955, 1.188] & 0.2773 & 0.3774 & 0.41 & 0.2530 & 17.5 \\
go:schubfach-vs-dragonbox & e2e & verify & numeric-boundary & schubfach & dragonbox & dragonbox & 1.07$\times$ & [0.941, 1.203] & 0.2601 & 0.3599 & 0.41 & 0.3104 & 21.7 \\
go:ryu-vs-dragonbox & e2e & verify & numeric-boundary & ryu & dragonbox & dragonbox & 1.01$\times$ & [0.883, 1.145] & 0.8910 & 0.9216 & 0.05 & 0.3104 & 21.7 \\
go:ryu-vs-bd & e2e & verify & numeric-boundary & ryu & json-canon & ryu & 1.09$\times$ & [0.973, 1.216] & 0.1640 & 0.2476 & -0.53 & 0.2530 & 17.5 \\
go:dragonbox-vs-bd & e2e & verify & numeric-boundary & dragonbox & json-canon & dragonbox & 1.10$\times$ & [0.960, 1.240] & 0.1720 & 0.2552 & -0.52 & 0.3104 & 21.7 \\
rs:schubfach-vs-bd & e2e & verify & numeric-boundary & schubfach-rs & json-canon-rs & schubfach-rs$^*$ & 1.26$\times$ & [1.163, 1.352] & 0.0002 & 0.0004 & -1.99 & 0.0512 & 9.6 \\
rs:schubfach-vs-ryu & e2e & verify & numeric-boundary & schubfach-rs & ryu-rs & ryu-rs & 1.01$\times$ & [0.936, 1.094] & 0.7792 & 0.8375 & 0.10 & 0.0689 & 13.0 \\
rs:schubfach-vs-dragonbox & e2e & verify & numeric-boundary & schubfach-rs & dragonbox-rs & schubfach-rs & 1.03$\times$ & [0.971, 1.102] & 0.3593 & 0.4645 & -0.34 & 0.0512 & 9.6 \\
rs:ryu-vs-dragonbox & e2e & verify & numeric-boundary & ryu-rs & dragonbox-rs & ryu-rs & 1.04$\times$ & [0.971, 1.131] & 0.3027 & 0.4034 & -0.39 & 0.0689 & 13.0 \\
rs:ryu-vs-bd & e2e & verify & numeric-boundary & ryu-rs & json-canon-rs & ryu-rs$^*$ & 1.27$\times$ & [1.164, 1.392] & 0.0004 & 0.0008 & -1.88 & 0.0689 & 13.0 \\
rs:dragonbox-vs-bd & e2e & verify & numeric-boundary & dragonbox-rs & json-canon-rs & dragonbox-rs$^*$ & 1.22$\times$ & [1.125, 1.308] & 0.0002 & 0.0004 & -1.75 & 0.0503 & 9.1 \\
schubfach:go-vs-rs & e2e & verify & numeric-boundary & schubfach & schubfach-rs & schubfach-rs$^*$ & 2.87$\times$ & [2.653, 3.133] & 0.0002 & 0.0004 & 6.54 & 0.0512 & 9.6 \\
bd:go-vs-rs & e2e & verify & numeric-boundary & json-canon & json-canon-rs & json-canon-rs$^*$ & 2.34$\times$ & [2.116, 2.562] & 0.0002 & 0.0004 & 5.07 & 0.0857 & 12.7 \\
ryu:go-vs-rs & e2e & verify & numeric-boundary & ryu & ryu-rs & ryu-rs$^*$ & 2.73$\times$ & [2.454, 3.027] & 0.0002 & 0.0004 & 5.04 & 0.0689 & 13.0 \\
dragonbox:go-vs-rs & e2e & verify & numeric-boundary & dragonbox & dragonbox-rs & dragonbox-rs$^*$ & 2.59$\times$ & [2.326, 2.915] & 0.0002 & 0.0004 & 4.04 & 0.0503 & 9.1 \\
go:schubfach-vs-bd & e2e & verify & rfc-key-sorting & schubfach & json-canon & schubfach & 1.09$\times$ & [0.948, 1.258] & 0.2699 & 0.3708 & -0.41 & 0.3236 & 22.6 \\
go:schubfach-vs-ryu & e2e & verify & rfc-key-sorting & schubfach & ryu & schubfach & 1.03$\times$ & [0.889, 1.190] & 0.6987 & 0.7694 & -0.14 & 0.3236 & 22.6 \\
go:schubfach-vs-dragonbox & e2e & verify & rfc-key-sorting & schubfach & dragonbox & schubfach & 1.07$\times$ & [0.931, 1.222] & 0.3701 & 0.4751 & -0.33 & 0.3236 & 22.6 \\
go:ryu-vs-dragonbox & e2e & verify & rfc-key-sorting & ryu & dragonbox & ryu & 1.04$\times$ & [0.903, 1.190] & 0.6169 & 0.7083 & -0.18 & 0.3166 & 21.4 \\
go:ryu-vs-bd & e2e & verify & rfc-key-sorting & ryu & json-canon & ryu & 1.06$\times$ & [0.912, 1.218] & 0.4619 & 0.5698 & -0.27 & 0.3166 & 21.4 \\
go:dragonbox-vs-bd & e2e & verify & rfc-key-sorting & dragonbox & json-canon & dragonbox & 1.02$\times$ & [0.898, 1.161] & 0.7852 & 0.8398 & -0.10 & 0.2784 & 18.2 \\
rs:schubfach-vs-bd & e2e & verify & rfc-key-sorting & schubfach-rs & json-canon-rs & schubfach-rs & 1.02$\times$ & [0.947, 1.100] & 0.5989 & 0.6944 & -0.20 & 0.0552 & 9.9 \\
rs:schubfach-vs-ryu & e2e & verify & rfc-key-sorting & schubfach-rs & ryu-rs & ryu-rs & 1.08$\times$ & [0.987, 1.169] & 0.0982 & 0.1576 & 0.62 & 0.0742 & 14.3 \\
rs:schubfach-vs-dragonbox & e2e & verify & rfc-key-sorting & schubfach-rs & dragonbox-rs & dragonbox-rs & 1.02$\times$ & [0.953, 1.072] & 0.4801 & 0.5822 & 0.27 & 0.0348 & 6.4 \\
rs:ryu-vs-dragonbox & e2e & verify & rfc-key-sorting & ryu-rs & dragonbox-rs & ryu-rs & 1.05$\times$ & [0.978, 1.137] & 0.2054 & 0.2956 & -0.48 & 0.0742 & 14.3 \\
rs:ryu-vs-bd & e2e & verify & rfc-key-sorting & ryu-rs & json-canon-rs & ryu-rs & 1.10$\times$ & [1.005, 1.204] & 0.0494 & 0.0849 & -0.74 & 0.0742 & 14.3 \\
rs:dragonbox-vs-bd & e2e & verify & rfc-key-sorting & dragonbox-rs & json-canon-rs & dragonbox-rs & 1.04$\times$ & [0.975, 1.107] & 0.2080 & 0.2987 & -0.47 & 0.0348 & 6.4 \\
schubfach:go-vs-rs & e2e & verify & rfc-key-sorting & schubfach & schubfach-rs & schubfach-rs$^*$ & 2.57$\times$ & [2.271, 2.890] & 0.0002 & 0.0004 & 3.86 & 0.0552 & 9.9 \\
bd:go-vs-rs & e2e & verify & rfc-key-sorting & json-canon & json-canon-rs & json-canon-rs$^*$ & 2.74$\times$ & [2.450, 3.064] & 0.0002 & 0.0004 & 4.42 & 0.0657 & 11.5 \\
ryu:go-vs-rs & e2e & verify & rfc-key-sorting & ryu & ryu-rs & ryu-rs$^*$ & 2.85$\times$ & [2.516, 3.219] & 0.0002 & 0.0004 & 4.27 & 0.0742 & 14.3 \\
dragonbox:go-vs-rs & e2e & verify & rfc-key-sorting & dragonbox & dragonbox-rs & dragonbox-rs$^*$ & 2.81$\times$ & [2.540, 3.052] & 0.0002 & 0.0004 & 5.08 & 0.0348 & 6.4 \\
go:schubfach-vs-bd & e2e & verify & small & schubfach & json-canon & schubfach$^*$ & 1.22$\times$ & [1.069, 1.407] & 0.0122 & 0.0231 & -1.00 & 0.2995 & 21.5 \\
go:schubfach-vs-ryu & e2e & verify & small & schubfach & ryu & schubfach & 1.02$\times$ & [0.896, 1.172] & 0.7626 & 0.8258 & -0.11 & 0.2995 & 21.5 \\
go:schubfach-vs-dragonbox & e2e & verify & small & schubfach & dragonbox & dragonbox & 1.00$\times$ & [0.862, 1.154] & 0.9988 & 0.9994 & 0.00 & 0.2984 & 21.4 \\
go:ryu-vs-dragonbox & e2e & verify & small & ryu & dragonbox & dragonbox & 1.02$\times$ & [0.893, 1.158] & 0.7674 & 0.8282 & 0.11 & 0.2984 & 21.4 \\
go:ryu-vs-bd & e2e & verify & small & ryu & json-canon & ryu$^*$ & 1.20$\times$ & [1.062, 1.361] & 0.0114 & 0.0217 & -0.98 & 0.2454 & 17.3 \\
go:dragonbox-vs-bd & e2e & verify & small & dragonbox & json-canon & dragonbox$^*$ & 1.22$\times$ & [1.063, 1.404] & 0.0110 & 0.0210 & -1.00 & 0.2984 & 21.4 \\
rs:schubfach-vs-bd & e2e & verify & small & schubfach-rs & json-canon-rs & schubfach-rs$^*$ & 1.25$\times$ & [1.184, 1.331] & 0.0002 & 0.0004 & -2.53 & 0.0385 & 7.0 \\
rs:schubfach-vs-ryu & e2e & verify & small & schubfach-rs & ryu-rs & ryu-rs & 1.01$\times$ & [0.953, 1.083] & 0.8572 & 0.8967 & 0.07 & 0.0612 & 11.2 \\
rs:schubfach-vs-dragonbox & e2e & verify & small & schubfach-rs & dragonbox-rs & schubfach-rs & 1.02$\times$ & [0.975, 1.081] & 0.4721 & 0.5766 & -0.27 & 0.0385 & 7.0 \\
rs:ryu-vs-dragonbox & e2e & verify & small & ryu-rs & dragonbox-rs & ryu-rs & 1.03$\times$ & [0.967, 1.109] & 0.4487 & 0.5578 & -0.28 & 0.0612 & 11.2 \\
rs:ryu-vs-bd & e2e & verify & small & ryu-rs & json-canon-rs & ryu-rs$^*$ & 1.26$\times$ & [1.180, 1.366] & 0.0002 & 0.0004 & -2.22 & 0.0612 & 11.2 \\
rs:dragonbox-vs-bd & e2e & verify & small & dragonbox-rs & json-canon-rs & dragonbox-rs$^*$ & 1.22$\times$ & [1.151, 1.301] & 0.0002 & 0.0004 & -2.21 & 0.0461 & 8.2 \\
schubfach:go-vs-rs & e2e & verify & small & schubfach & schubfach-rs & schubfach-rs$^*$ & 2.53$\times$ & [2.254, 2.797] & 0.0002 & 0.0004 & 4.03 & 0.0385 & 7.0 \\
bd:go-vs-rs & e2e & verify & small & json-canon & json-canon-rs & json-canon-rs$^*$ & 2.48$\times$ & [2.224, 2.741] & 0.0002 & 0.0004 & 4.33 & 0.0680 & 9.9 \\
ryu:go-vs-rs & e2e & verify & small & ryu & ryu-rs & ryu-rs$^*$ & 2.60$\times$ & [2.371, 2.891] & 0.0002 & 0.0004 & 5.00 & 0.0612 & 11.2 \\
dragonbox:go-vs-rs & e2e & verify & small & dragonbox & dragonbox-rs & dragonbox-rs$^*$ & 2.48$\times$ & [2.217, 2.766] & 0.0002 & 0.0004 & 3.97 & 0.0461 & 8.2 \\
go:schubfach-vs-bd & e2e & verify & surrogate-pair & schubfach & json-canon & schubfach & 1.09$\times$ & [0.948, 1.240] & 0.2515 & 0.3508 & -0.43 & 0.2810 & 18.9 \\
go:schubfach-vs-ryu & e2e & verify & surrogate-pair & schubfach & ryu & schubfach & 1.05$\times$ & [0.945, 1.179] & 0.3573 & 0.4624 & -0.34 & 0.2810 & 18.9 \\
go:schubfach-vs-dragonbox & e2e & verify & surrogate-pair & schubfach & dragonbox & schubfach & 1.06$\times$ & [0.960, 1.185] & 0.2779 & 0.3778 & -0.40 & 0.2810 & 18.9 \\
go:ryu-vs-dragonbox & e2e & verify & surrogate-pair & ryu & dragonbox & ryu & 1.01$\times$ & [0.933, 1.099] & 0.8854 & 0.9173 & -0.05 & 0.2001 & 12.8 \\
go:ryu-vs-bd & e2e & verify & surrogate-pair & ryu & json-canon & ryu & 1.03$\times$ & [0.919, 1.161] & 0.6271 & 0.7156 & -0.18 & 0.2001 & 12.8 \\
go:dragonbox-vs-bd & e2e & verify & surrogate-pair & dragonbox & json-canon & dragonbox & 1.02$\times$ & [0.911, 1.145] & 0.6883 & 0.7623 & -0.15 & 0.1711 & 10.9 \\
rs:schubfach-vs-bd & e2e & verify & surrogate-pair & schubfach-rs & json-canon-rs & schubfach-rs$^*$ & 1.09$\times$ & [1.023, 1.152] & 0.0088 & 0.0169 & -1.01 & 0.0395 & 7.1 \\
rs:schubfach-vs-ryu & e2e & verify & surrogate-pair & schubfach-rs & ryu-rs & ryu-rs & 1.03$\times$ & [0.946, 1.124] & 0.5363 & 0.6349 & 0.23 & 0.0926 & 17.1 \\
rs:schubfach-vs-dragonbox & e2e & verify & surrogate-pair & schubfach-rs & dragonbox-rs & dragonbox-rs & 1.02$\times$ & [0.934, 1.089] & 0.6805 & 0.7584 & 0.16 & 0.0776 & 14.2 \\
rs:ryu-vs-dragonbox & e2e & verify & surrogate-pair & ryu-rs & dragonbox-rs & ryu-rs & 1.01$\times$ & [0.911, 1.130] & 0.8296 & 0.8772 & -0.07 & 0.0926 & 17.1 \\
rs:ryu-vs-bd & e2e & verify & surrogate-pair & ryu-rs & json-canon-rs & ryu-rs$^*$ & 1.12$\times$ & [1.021, 1.237] & 0.0264 & 0.0481 & -0.87 & 0.0926 & 17.1 \\
rs:dragonbox-vs-bd & e2e & verify & surrogate-pair & dragonbox-rs & json-canon-rs & dragonbox-rs$^*$ & 1.11$\times$ & [1.012, 1.197] & 0.0194 & 0.0360 & -0.89 & 0.0776 & 14.2 \\
schubfach:go-vs-rs & e2e & verify & surrogate-pair & schubfach & schubfach-rs & schubfach-rs$^*$ & 2.67$\times$ & [2.398, 2.925] & 0.0002 & 0.0004 & 4.73 & 0.0395 & 7.1 \\
bd:go-vs-rs & e2e & verify & surrogate-pair & json-canon & json-canon-rs & json-canon-rs$^*$ & 2.66$\times$ & [2.367, 2.947] & 0.0002 & 0.0004 & 4.21 & 0.0615 & 10.1 \\
ryu:go-vs-rs & e2e & verify & surrogate-pair & ryu & ryu-rs & ryu-rs$^*$ & 2.90$\times$ & [2.585, 3.188] & 0.0002 & 0.0004 & 6.72 & 0.0926 & 17.1 \\
dragonbox:go-vs-rs & e2e & verify & surrogate-pair & dragonbox & dragonbox-rs & dragonbox-rs$^*$ & 2.88$\times$ & [2.610, 3.115] & 0.0002 & 0.0004 & 7.91 & 0.0776 & 14.2 \\
go:schubfach-vs-bd & e2e & verify & unicode & schubfach & json-canon & schubfach & 1.06$\times$ & [0.958, 1.179] & 0.3065 & 0.4079 & -0.38 & 0.2361 & 16.0 \\
go:schubfach-vs-ryu & e2e & verify & unicode & schubfach & ryu & ryu & 1.05$\times$ & [0.930, 1.169] & 0.4685 & 0.5746 & 0.27 & 0.2514 & 17.8 \\
go:schubfach-vs-dragonbox & e2e & verify & unicode & schubfach & dragonbox & schubfach & 1.26$\times$ & [1.055, 1.899] & 0.0468 & 0.0811 & -0.53 & 0.2361 & 16.0 \\
go:ryu-vs-dragonbox & e2e & verify & unicode & ryu & dragonbox & ryu$^*$ & 1.32$\times$ & [1.101, 1.990] & 0.0056 & 0.0109 & -0.62 & 0.2514 & 17.8 \\
go:ryu-vs-bd & e2e & verify & unicode & ryu & json-canon & ryu & 1.11$\times$ & [0.995, 1.243] & 0.0928 & 0.1505 & -0.64 & 0.2514 & 17.8 \\
go:dragonbox-vs-bd & e2e & verify & unicode & dragonbox & json-canon & json-canon & 1.19$\times$ & [0.999, 1.760] & 0.2737 & 0.3748 & 0.41 & 0.2293 & 14.7 \\
rs:schubfach-vs-bd & e2e & verify & unicode & schubfach-rs & json-canon-rs & schubfach-rs$^*$ & 1.13$\times$ & [1.033, 1.229] & 0.0132 & 0.0248 & -0.98 & 0.0489 & 9.0 \\
rs:schubfach-vs-ryu & e2e & verify & unicode & schubfach-rs & ryu-rs & ryu-rs & 1.02$\times$ & [0.964, 1.097] & 0.4895 & 0.5925 & 0.25 & 0.0531 & 10.0 \\
rs:schubfach-vs-dragonbox & e2e & verify & unicode & schubfach-rs & dragonbox-rs & schubfach-rs & 1.02$\times$ & [0.960, 1.083] & 0.4477 & 0.5571 & -0.29 & 0.0489 & 9.0 \\
rs:ryu-vs-dragonbox & e2e & verify & unicode & ryu-rs & dragonbox-rs & ryu-rs & 1.05$\times$ & [0.982, 1.115] & 0.1670 & 0.2509 & -0.53 & 0.0531 & 10.0 \\
rs:ryu-vs-bd & e2e & verify & unicode & ryu-rs & json-canon-rs & ryu-rs$^*$ & 1.16$\times$ & [1.057, 1.260] & 0.0040 & 0.0079 & -1.13 & 0.0531 & 10.0 \\
rs:dragonbox-vs-bd & e2e & verify & unicode & dragonbox-rs & json-canon-rs & dragonbox-rs & 1.10$\times$ & [1.013, 1.196] & 0.0380 & 0.0673 & -0.80 & 0.0469 & 8.4 \\
schubfach:go-vs-rs & e2e & verify & unicode & schubfach & schubfach-rs & schubfach-rs$^*$ & 2.70$\times$ & [2.459, 2.937] & 0.0002 & 0.0004 & 5.57 & 0.0489 & 9.0 \\
bd:go-vs-rs & e2e & verify & unicode & json-canon & json-canon-rs & json-canon-rs$^*$ & 2.53$\times$ & [2.295, 2.805] & 0.0002 & 0.0004 & 5.51 & 0.0930 & 15.1 \\
ryu:go-vs-rs & e2e & verify & unicode & ryu & ryu-rs & ryu-rs$^*$ & 2.65$\times$ & [2.400, 2.915] & 0.0002 & 0.0004 & 4.93 & 0.0531 & 10.0 \\
dragonbox:go-vs-rs & e2e & verify & unicode & dragonbox & dragonbox-rs & dragonbox-rs$^*$ & 3.33$\times$ & [2.818, 5.147] & 0.0002 & 0.0004 & 1.85 & 0.0469 & 8.4 \\
go:schubfach-vs-bd & e2e & verify & verify-whitespace & schubfach & json-canon & schubfach & 1.02$\times$ & [0.891, 1.180] & 0.7598 & 0.8242 & -0.11 & 0.3596 & 22.4 \\
go:schubfach-vs-ryu & e2e & verify & verify-whitespace & schubfach & ryu & ryu & 1.15$\times$ & [0.988, 1.314] & 0.0918 & 0.1494 & 0.64 & 0.2944 & 21.1 \\
go:schubfach-vs-dragonbox & e2e & verify & verify-whitespace & schubfach & dragonbox & dragonbox & 1.09$\times$ & [0.941, 1.246] & 0.2831 & 0.3825 & 0.41 & 0.2708 & 18.3 \\
go:ryu-vs-dragonbox & e2e & verify & verify-whitespace & ryu & dragonbox & ryu & 1.06$\times$ & [0.919, 1.199] & 0.4279 & 0.5360 & -0.28 & 0.2944 & 21.1 \\
go:ryu-vs-bd & e2e & verify & verify-whitespace & ryu & json-canon & ryu & 1.17$\times$ & [1.028, 1.340] & 0.0370 & 0.0656 & -0.82 & 0.2944 & 21.1 \\
go:dragonbox-vs-bd & e2e & verify & verify-whitespace & dragonbox & json-canon & dragonbox & 1.11$\times$ & [0.986, 1.261] & 0.1340 & 0.2086 & -0.57 & 0.2708 & 18.3 \\
rs:schubfach-vs-bd & e2e & verify & verify-whitespace & schubfach-rs & json-canon-rs & schubfach-rs$^*$ & 1.26$\times$ & [1.161, 1.361] & 0.0004 & 0.0008 & -1.95 & 0.0638 & 11.8 \\
rs:schubfach-vs-ryu & e2e & verify & verify-whitespace & schubfach-rs & ryu-rs & schubfach-rs & 1.00$\times$ & [0.910, 1.090] & 0.9404 & 0.9601 & -0.03 & 0.0638 & 11.8 \\
rs:schubfach-vs-dragonbox & e2e & verify & verify-whitespace & schubfach-rs & dragonbox-rs & schubfach-rs & 1.02$\times$ & [0.940, 1.107] & 0.6777 & 0.7559 & -0.16 & 0.0638 & 11.8 \\
rs:ryu-vs-dragonbox & e2e & verify & verify-whitespace & ryu-rs & dragonbox-rs & ryu-rs & 1.01$\times$ & [0.932, 1.121] & 0.7680 & 0.8282 & -0.11 & 0.0795 & 14.7 \\
rs:ryu-vs-bd & e2e & verify & verify-whitespace & ryu-rs & json-canon-rs & ryu-rs$^*$ & 1.25$\times$ & [1.146, 1.379] & 0.0002 & 0.0004 & -1.75 & 0.0795 & 14.7 \\
rs:dragonbox-vs-bd & e2e & verify & verify-whitespace & dragonbox-rs & json-canon-rs & dragonbox-rs$^*$ & 1.23$\times$ & [1.137, 1.341] & 0.0002 & 0.0004 & -1.78 & 0.0669 & 12.2 \\
schubfach:go-vs-rs & e2e & verify & verify-whitespace & schubfach & schubfach-rs & schubfach-rs$^*$ & 2.97$\times$ & [2.612, 3.340] & 0.0002 & 0.0004 & 4.22 & 0.0638 & 11.8 \\
bd:go-vs-rs & e2e & verify & verify-whitespace & json-canon & json-canon-rs & json-canon-rs$^*$ & 2.42$\times$ & [2.172, 2.692] & 0.0002 & 0.0004 & 4.36 & 0.0797 & 11.8 \\
ryu:go-vs-rs & e2e & verify & verify-whitespace & ryu & ryu-rs & ryu-rs$^*$ & 2.58$\times$ & [2.287, 2.927] & 0.0002 & 0.0004 & 4.06 & 0.0795 & 14.7 \\
dragonbox:go-vs-rs & e2e & verify & verify-whitespace & dragonbox & dragonbox-rs & dragonbox-rs$^*$ & 2.69$\times$ & [2.382, 2.968] & 0.0002 & 0.0004 & 4.81 & 0.0669 & 12.2 \\
go:schubfach-vs-bd & api & canonicalize & array-2048 & schubfach & json-canon & schubfach$^*$ & 1.64$\times$ & [1.631, 1.648] & 0.0002 & 0.0004 & -81.57 & 0.0177 & 0.9 \\
go:schubfach-vs-ryu & api & canonicalize & array-2048 & schubfach & ryu & ryu$^*$ & 1.01$\times$ & [1.007, 1.018] & 0.0006 & 0.0013 & 1.73 & 0.0155 & 0.8 \\
go:schubfach-vs-dragonbox & api & canonicalize & array-2048 & schubfach & dragonbox & dragonbox & 1.00$\times$ & [0.996, 1.007] & 0.8756 & 0.9101 & 0.08 & 0.0155 & 0.8 \\
go:ryu-vs-dragonbox & api & canonicalize & array-2048 & ryu & dragonbox & ryu$^*$ & 1.01$\times$ & [1.006, 1.016] & 0.0022 & 0.0045 & -1.78 & 0.0155 & 0.8 \\
go:ryu-vs-bd & api & canonicalize & array-2048 & ryu & json-canon & ryu$^*$ & 1.66$\times$ & [1.652, 1.668] & 0.0002 & 0.0004 & -87.22 & 0.0155 & 0.8 \\
go:dragonbox-vs-bd & api & canonicalize & array-2048 & dragonbox & json-canon & dragonbox$^*$ & 1.64$\times$ & [1.635, 1.650] & 0.0002 & 0.0004 & -85.65 & 0.0155 & 0.8 \\
rs:schubfach-vs-bd & api & canonicalize & array-2048 & schubfach-rs & json-canon-rs & schubfach-rs$^*$ & 2.27$\times$ & [1.910, 2.545] & 0.0002 & 0.0004 & -7.67 & 0.3523 & 29.1 \\
rs:schubfach-vs-ryu & api & canonicalize & array-2048 & schubfach-rs & ryu-rs & ryu-rs & 1.08$\times$ & [0.958, 1.310] & 0.3667 & 0.4712 & 0.40 & 0.1717 & 15.3 \\
rs:schubfach-vs-dragonbox & api & canonicalize & array-2048 & schubfach-rs & dragonbox-rs & dragonbox-rs$^*$ & 1.13$\times$ & [1.008, 1.340] & 0.0040 & 0.0079 & 0.71 & 0.0070 & 0.7 \\
rs:ryu-vs-dragonbox & api & canonicalize & array-2048 & ryu-rs & dragonbox-rs & dragonbox-rs$^*$ & 1.05$\times$ & [1.006, 1.174] & 0.0060 & 0.0117 & 0.54 & 0.0070 & 0.7 \\
rs:ryu-vs-bd & api & canonicalize & array-2048 & ryu-rs & json-canon-rs & ryu-rs$^*$ & 2.44$\times$ & [2.117, 2.548] & 0.0002 & 0.0004 & -16.56 & 0.1717 & 15.3 \\
rs:dragonbox-vs-bd & api & canonicalize & array-2048 & dragonbox-rs & json-canon-rs & dragonbox-rs$^*$ & 2.56$\times$ & [2.551, 2.576] & 0.0002 & 0.0004 & -130.15 & 0.0070 & 0.7 \\
schubfach:go-vs-rs & api & canonicalize & array-2048 & schubfach & schubfach-rs & schubfach-rs$^*$ & 1.67$\times$ & [1.418, 1.880] & 0.0002 & 0.0004 & 4.08 & 0.3523 & 29.1 \\
bd:go-vs-rs & api & canonicalize & array-2048 & json-canon & json-canon-rs & json-canon-rs$^*$ & 1.21$\times$ & [1.206, 1.218] & 0.0002 & 0.0004 & 33.39 & 0.0216 & 0.8 \\
ryu:go-vs-rs & api & canonicalize & array-2048 & ryu & ryu-rs & ryu-rs$^*$ & 1.78$\times$ & [1.546, 1.860] & 0.0002 & 0.0004 & 9.03 & 0.1717 & 15.3 \\
dragonbox:go-vs-rs & api & canonicalize & array-2048 & dragonbox & dragonbox-rs & dragonbox-rs$^*$ & 1.89$\times$ & [1.883, 1.901] & 0.0002 & 0.0004 & 98.98 & 0.0070 & 0.7 \\
go:schubfach-vs-bd & api & canonicalize & array-256 & schubfach & json-canon & schubfach$^*$ & 1.57$\times$ & [1.558, 1.581] & 0.0002 & 0.0004 & -54.05 & 0.0032 & 1.3 \\
go:schubfach-vs-ryu & api & canonicalize & array-256 & schubfach & ryu & ryu$^*$ & 1.01$\times$ & [1.002, 1.016] & 0.0268 & 0.0488 & 1.04 & 0.0017 & 0.7 \\
go:schubfach-vs-dragonbox & api & canonicalize & array-256 & schubfach & dragonbox & dragonbox$^*$ & 1.01$\times$ & [1.003, 1.017] & 0.0142 & 0.0267 & 1.14 & 0.0013 & 0.5 \\
go:ryu-vs-dragonbox & api & canonicalize & array-256 & ryu & dragonbox & dragonbox & 1.00$\times$ & [0.996, 1.004] & 0.8696 & 0.9075 & 0.07 & 0.0013 & 0.5 \\
go:ryu-vs-bd & api & canonicalize & array-256 & ryu & json-canon & ryu$^*$ & 1.58$\times$ & [1.576, 1.592] & 0.0002 & 0.0004 & -69.20 & 0.0017 & 0.7 \\
go:dragonbox-vs-bd & api & canonicalize & array-256 & dragonbox & json-canon & dragonbox$^*$ & 1.58$\times$ & [1.577, 1.592] & 0.0002 & 0.0004 & -72.77 & 0.0013 & 0.5 \\
rs:schubfach-vs-bd & api & canonicalize & array-256 & schubfach-rs & json-canon-rs & schubfach-rs$^*$ & 2.30$\times$ & [2.252, 2.353] & 0.0002 & 0.0004 & -22.56 & 0.0025 & 1.9 \\
rs:schubfach-vs-ryu & api & canonicalize & array-256 & schubfach-rs & ryu-rs & schubfach-rs & 1.02$\times$ & [1.001, 1.083] & 0.2306 & 0.3268 & -0.54 & 0.0025 & 1.9 \\
rs:schubfach-vs-dragonbox & api & canonicalize & array-256 & schubfach-rs & dragonbox-rs & dragonbox-rs & 1.00$\times$ & [0.994, 1.014] & 0.8124 & 0.8631 & 0.13 & 0.0009 & 0.7 \\
rs:ryu-vs-dragonbox & api & canonicalize & array-256 & ryu-rs & dragonbox-rs & dragonbox-rs & 1.02$\times$ & [1.004, 1.078] & 0.0394 & 0.0694 & 0.59 & 0.0009 & 0.7 \\
rs:ryu-vs-bd & api & canonicalize & array-256 & ryu-rs & json-canon-rs & ryu-rs$^*$ & 2.25$\times$ & [2.136, 2.314] & 0.0002 & 0.0004 & -18.57 & 0.0093 & 6.8 \\
rs:dragonbox-vs-bd & api & canonicalize & array-256 & dragonbox-rs & json-canon-rs & dragonbox-rs$^*$ & 2.30$\times$ & [2.260, 2.356] & 0.0002 & 0.0004 & -22.92 & 0.0009 & 0.7 \\
schubfach:go-vs-rs & api & canonicalize & array-256 & schubfach & schubfach-rs & schubfach-rs$^*$ & 1.81$\times$ & [1.789, 1.829] & 0.0002 & 0.0004 & 47.16 & 0.0025 & 1.9 \\
bd:go-vs-rs & api & canonicalize & array-256 & json-canon & json-canon-rs & json-canon-rs$^*$ & 1.24$\times$ & [1.209, 1.261] & 0.0002 & 0.0004 & 9.41 & 0.0135 & 4.4 \\
ryu:go-vs-rs & api & canonicalize & array-256 & ryu & ryu-rs & ryu-rs$^*$ & 1.76$\times$ & [1.646, 1.792] & 0.0002 & 0.0004 & 19.51 & 0.0093 & 6.8 \\
dragonbox:go-vs-rs & api & canonicalize & array-256 & dragonbox & dragonbox-rs & dragonbox-rs$^*$ & 1.80$\times$ & [1.789, 1.804] & 0.0002 & 0.0004 & 121.35 & 0.0009 & 0.7 \\
go:schubfach-vs-bd & api & canonicalize & canonical-minimal & schubfach & json-canon & schubfach$^*$ & 1.49$\times$ & [1.482, 1.510] & 0.0002 & 0.0004 & -38.06 & 0.0000 & 1.8 \\
go:schubfach-vs-ryu & api & canonicalize & canonical-minimal & schubfach & ryu & ryu$^*$ & 1.03$\times$ & [1.018, 1.038] & 0.0002 & 0.0004 & 2.40 & 0.0000 & 1.1 \\
go:schubfach-vs-dragonbox & api & canonicalize & canonical-minimal & schubfach & dragonbox & dragonbox$^*$ & 1.02$\times$ & [1.014, 1.032] & 0.0006 & 0.0013 & 2.17 & 0.0000 & 0.7 \\
go:ryu-vs-dragonbox & api & canonicalize & canonical-minimal & ryu & dragonbox & ryu & 1.00$\times$ & [0.999, 1.011] & 0.1680 & 0.2516 & -0.63 & 0.0000 & 1.1 \\
go:ryu-vs-bd & api & canonicalize & canonical-minimal & ryu & json-canon & ryu$^*$ & 1.54$\times$ & [1.527, 1.549] & 0.0002 & 0.0004 & -51.22 & 0.0000 & 1.1 \\
go:dragonbox-vs-bd & api & canonicalize & canonical-minimal & dragonbox & json-canon & dragonbox$^*$ & 1.53$\times$ & [1.522, 1.539] & 0.0002 & 0.0004 & -56.22 & 0.0000 & 0.7 \\
rs:schubfach-vs-bd & api & canonicalize & canonical-minimal & schubfach-rs & json-canon-rs & schubfach-rs$^*$ & 2.10$\times$ & [2.023, 2.191] & 0.0002 & 0.0004 & -14.78 & 0.0000 & 6.5 \\
rs:schubfach-vs-ryu & api & canonicalize & canonical-minimal & schubfach-rs & ryu-rs & ryu-rs & 1.04$\times$ & [0.989, 1.079] & 0.1028 & 0.1646 & 0.76 & 0.0000 & 6.4 \\
rs:schubfach-vs-dragonbox & api & canonicalize & canonical-minimal & schubfach-rs & dragonbox-rs & dragonbox-rs & 1.01$\times$ & [0.942, 1.061] & 0.6945 & 0.7670 & 0.19 & 0.0000 & 10.6 \\
rs:ryu-vs-dragonbox & api & canonicalize & canonical-minimal & ryu-rs & dragonbox-rs & ryu-rs & 1.03$\times$ & [0.978, 1.102] & 0.4983 & 0.5998 & -0.36 & 0.0000 & 6.4 \\
rs:ryu-vs-bd & api & canonicalize & canonical-minimal & ryu-rs & json-canon-rs & ryu-rs$^*$ & 2.19$\times$ & [2.085, 2.259] & 0.0002 & 0.0004 & -15.49 & 0.0000 & 6.4 \\
rs:dragonbox-vs-bd & api & canonicalize & canonical-minimal & dragonbox-rs & json-canon-rs & dragonbox-rs$^*$ & 2.13$\times$ & [1.981, 2.226] & 0.0002 & 0.0004 & -12.68 & 0.0000 & 10.6 \\
schubfach:go-vs-rs & api & canonicalize & canonical-minimal & schubfach & schubfach-rs & schubfach-rs$^*$ & 1.68$\times$ & [1.624, 1.729] & 0.0002 & 0.0004 & 16.87 & 0.0000 & 6.5 \\
bd:go-vs-rs & api & canonicalize & canonical-minimal & json-canon & json-canon-rs & json-canon-rs$^*$ & 1.19$\times$ & [1.158, 1.220] & 0.0002 & 0.0004 & 6.09 & 0.0000 & 5.5 \\
ryu:go-vs-rs & api & canonicalize & canonical-minimal & ryu & ryu-rs & ryu-rs$^*$ & 1.70$\times$ & [1.612, 1.730] & 0.0002 & 0.0004 & 18.61 & 0.0000 & 6.4 \\
dragonbox:go-vs-rs & api & canonicalize & canonical-minimal & dragonbox & dragonbox-rs & dragonbox-rs$^*$ & 1.66$\times$ & [1.541, 1.722] & 0.0002 & 0.0004 & 10.92 & 0.0000 & 10.6 \\
go:schubfach-vs-bd & api & canonicalize & control-escapes & schubfach & json-canon & schubfach$^*$ & 1.01$\times$ & [1.009, 1.019] & 0.0004 & 0.0008 & -2.25 & 0.0000 & 1.0 \\
go:schubfach-vs-ryu & api & canonicalize & control-escapes & schubfach & ryu & ryu & 1.01$\times$ & [1.001, 1.016] & 0.0504 & 0.0865 & 0.91 & 0.0000 & 1.3 \\
go:schubfach-vs-dragonbox & api & canonicalize & control-escapes & schubfach & dragonbox & dragonbox & 1.00$\times$ & [0.996, 1.012] & 0.4271 & 0.5355 & 0.36 & 0.0000 & 1.4 \\
go:ryu-vs-dragonbox & api & canonicalize & control-escapes & ryu & dragonbox & ryu & 1.00$\times$ & [0.996, 1.013] & 0.2959 & 0.3960 & -0.47 & 0.0000 & 1.3 \\
go:ryu-vs-bd & api & canonicalize & control-escapes & ryu & json-canon & ryu$^*$ & 1.02$\times$ & [1.016, 1.029] & 0.0004 & 0.0008 & -2.90 & 0.0000 & 1.3 \\
go:dragonbox-vs-bd & api & canonicalize & control-escapes & dragonbox & json-canon & dragonbox$^*$ & 1.02$\times$ & [1.011, 1.026] & 0.0002 & 0.0004 & -2.16 & 0.0000 & 1.4 \\
rs:schubfach-vs-bd & api & canonicalize & control-escapes & schubfach-rs & json-canon-rs & json-canon-rs & 1.02$\times$ & [0.981, 1.052] & 0.3429 & 0.4464 & 0.43 & 0.0000 & 6.9 \\
rs:schubfach-vs-ryu & api & canonicalize & control-escapes & schubfach-rs & ryu-rs & schubfach-rs & 1.03$\times$ & [0.989, 1.097] & 0.2747 & 0.3750 & -0.52 & 0.0000 & 3.0 \\
rs:schubfach-vs-dragonbox & api & canonicalize & control-escapes & schubfach-rs & dragonbox-rs & dragonbox-rs & 1.04$\times$ & [0.995, 1.071] & 0.0524 & 0.0898 & 0.95 & 0.0000 & 6.8 \\
rs:ryu-vs-dragonbox & api & canonicalize & control-escapes & ryu-rs & dragonbox-rs & dragonbox-rs & 1.08$\times$ & [1.017, 1.147] & 0.0356 & 0.0635 & 1.00 & 0.0000 & 6.8 \\
rs:ryu-vs-bd & api & canonicalize & control-escapes & ryu-rs & json-canon-rs & json-canon-rs & 1.05$\times$ & [0.998, 1.126] & 0.1336 & 0.2082 & 0.71 & 0.0000 & 6.9 \\
rs:dragonbox-vs-bd & api & canonicalize & control-escapes & dragonbox-rs & json-canon-rs & dragonbox-rs & 1.02$\times$ & [0.974, 1.066] & 0.3929 & 0.5004 & -0.39 & 0.0000 & 6.8 \\
schubfach:go-vs-rs & api & canonicalize & control-escapes & schubfach & schubfach-rs & schubfach-rs$^*$ & 2.00$\times$ & [1.971, 2.028] & 0.0002 & 0.0004 & 49.04 & 0.0000 & 3.0 \\
bd:go-vs-rs & api & canonicalize & control-escapes & json-canon & json-canon-rs & json-canon-rs$^*$ & 2.06$\times$ & [1.996, 2.126] & 0.0002 & 0.0004 & 27.13 & 0.0000 & 6.9 \\
ryu:go-vs-rs & api & canonicalize & control-escapes & ryu & ryu-rs & ryu-rs$^*$ & 1.92$\times$ & [1.799, 1.998] & 0.0002 & 0.0004 & 14.82 & 0.0001 & 10.7 \\
dragonbox:go-vs-rs & api & canonicalize & control-escapes & dragonbox & dragonbox-rs & dragonbox-rs$^*$ & 2.07$\times$ & [1.985, 2.127] & 0.0002 & 0.0004 & 25.74 & 0.0000 & 6.8 \\
go:schubfach-vs-bd & api & canonicalize & deep-64 & schubfach & json-canon & json-canon & 1.00$\times$ & [0.999, 1.007] & 0.1824 & 0.2669 & 0.62 & 0.0001 & 0.6 \\
go:schubfach-vs-ryu & api & canonicalize & deep-64 & schubfach & ryu & schubfach$^*$ & 1.01$\times$ & [1.008, 1.016] & 0.0002 & 0.0004 & -2.44 & 0.0001 & 0.6 \\
go:schubfach-vs-dragonbox & api & canonicalize & deep-64 & schubfach & dragonbox & schubfach$^*$ & 1.02$\times$ & [1.013, 1.025] & 0.0004 & 0.0008 & -2.67 & 0.0001 & 0.6 \\
go:ryu-vs-dragonbox & api & canonicalize & deep-64 & ryu & dragonbox & ryu & 1.01$\times$ & [1.000, 1.012] & 0.0646 & 0.1097 & -0.89 & 0.0001 & 0.7 \\
go:ryu-vs-bd & api & canonicalize & deep-64 & ryu & json-canon & json-canon$^*$ & 1.02$\times$ & [1.011, 1.020] & 0.0002 & 0.0004 & 3.06 & 0.0001 & 0.6 \\
go:dragonbox-vs-bd & api & canonicalize & deep-64 & dragonbox & json-canon & json-canon$^*$ & 1.02$\times$ & [1.016, 1.028] & 0.0002 & 0.0004 & 3.11 & 0.0001 & 0.6 \\
rs:schubfach-vs-bd & api & canonicalize & deep-64 & schubfach-rs & json-canon-rs & json-canon-rs & 1.05$\times$ & [1.004, 1.154] & 0.1570 & 0.2387 & 0.59 & 0.0004 & 4.2 \\
rs:schubfach-vs-ryu & api & canonicalize & deep-64 & schubfach-rs & ryu-rs & ryu-rs$^*$ & 1.06$\times$ & [1.021, 1.186] & 0.0002 & 0.0004 & 0.77 & 0.0000 & 0.6 \\
rs:schubfach-vs-dragonbox & api & canonicalize & deep-64 & schubfach-rs & dragonbox-rs & dragonbox-rs & 1.04$\times$ & [1.003, 1.139] & 0.0918 & 0.1494 & 0.54 & 0.0001 & 1.1 \\
rs:ryu-vs-dragonbox & api & canonicalize & deep-64 & ryu-rs & dragonbox-rs & ryu-rs$^*$ & 1.02$\times$ & [1.011, 1.023] & 0.0002 & 0.0004 & -2.31 & 0.0000 & 0.6 \\
rs:ryu-vs-bd & api & canonicalize & deep-64 & ryu-rs & json-canon-rs & ryu-rs & 1.01$\times$ & [0.994, 1.036] & 0.3901 & 0.4988 & -0.43 & 0.0000 & 0.6 \\
rs:dragonbox-vs-bd & api & canonicalize & deep-64 & dragonbox-rs & json-canon-rs & json-canon-rs & 1.01$\times$ & [0.981, 1.024] & 0.5913 & 0.6877 & 0.25 & 0.0004 & 4.2 \\
schubfach:go-vs-rs & api & canonicalize & deep-64 & schubfach & schubfach-rs & schubfach-rs$^*$ & 1.39$\times$ & [1.241, 1.439] & 0.0002 & 0.0004 & 5.65 & 0.0011 & 12.2 \\
bd:go-vs-rs & api & canonicalize & deep-64 & json-canon & json-canon-rs & json-canon-rs$^*$ & 1.45$\times$ & [1.415, 1.474] & 0.0002 & 0.0004 & 18.48 & 0.0004 & 4.2 \\
ryu:go-vs-rs & api & canonicalize & deep-64 & ryu & ryu-rs & ryu-rs$^*$ & 1.49$\times$ & [1.481, 1.493] & 0.0002 & 0.0004 & 76.96 & 0.0000 & 0.6 \\
dragonbox:go-vs-rs & api & canonicalize & deep-64 & dragonbox & dragonbox-rs & dragonbox-rs$^*$ & 1.47$\times$ & [1.461, 1.483] & 0.0002 & 0.0004 & 43.09 & 0.0001 & 1.1 \\
go:schubfach-vs-bd & api & canonicalize & deep & schubfach & json-canon & json-canon & 1.00$\times$ & [0.996, 1.008] & 0.6075 & 0.6994 & 0.24 & 0.0002 & 0.6 \\
go:schubfach-vs-ryu & api & canonicalize & deep & schubfach & ryu & schubfach$^*$ & 1.02$\times$ & [1.011, 1.025] & 0.0004 & 0.0008 & -2.21 & 0.0004 & 1.2 \\
go:schubfach-vs-dragonbox & api & canonicalize & deep & schubfach & dragonbox & schubfach$^*$ & 1.02$\times$ & [1.014, 1.028] & 0.0004 & 0.0008 & -2.58 & 0.0004 & 1.2 \\
go:ryu-vs-dragonbox & api & canonicalize & deep & ryu & dragonbox & ryu & 1.00$\times$ & [0.998, 1.009] & 0.3521 & 0.4566 & -0.44 & 0.0003 & 0.9 \\
go:ryu-vs-bd & api & canonicalize & deep & ryu & json-canon & json-canon$^*$ & 1.02$\times$ & [1.015, 1.025] & 0.0002 & 0.0004 & 3.31 & 0.0002 & 0.6 \\
go:dragonbox-vs-bd & api & canonicalize & deep & dragonbox & json-canon & json-canon$^*$ & 1.02$\times$ & [1.019, 1.029] & 0.0002 & 0.0004 & 3.85 & 0.0002 & 0.6 \\
rs:schubfach-vs-bd & api & canonicalize & deep & schubfach-rs & json-canon-rs & schubfach-rs$^*$ & 1.04$\times$ & [1.026, 1.054] & 0.0004 & 0.0008 & -2.39 & 0.0005 & 2.0 \\
rs:schubfach-vs-ryu & api & canonicalize & deep & schubfach-rs & ryu-rs & schubfach-rs & 1.01$\times$ & [0.993, 1.059] & 0.5063 & 0.6077 & -0.39 & 0.0005 & 2.0 \\
rs:schubfach-vs-dragonbox & api & canonicalize & deep & schubfach-rs & dragonbox-rs & schubfach-rs$^*$ & 1.06$\times$ & [1.044, 1.069] & 0.0002 & 0.0004 & -3.94 & 0.0005 & 2.0 \\
rs:ryu-vs-dragonbox & api & canonicalize & deep & ryu-rs & dragonbox-rs & ryu-rs$^*$ & 1.04$\times$ & [0.996, 1.063] & 0.0130 & 0.0245 & -1.18 & 0.0014 & 6.1 \\
rs:ryu-vs-bd & api & canonicalize & deep & ryu-rs & json-canon-rs & ryu-rs & 1.03$\times$ & [0.979, 1.048] & 0.1194 & 0.1888 & -0.71 & 0.0014 & 6.1 \\
rs:dragonbox-vs-bd & api & canonicalize & deep & dragonbox-rs & json-canon-rs & json-canon-rs & 1.02$\times$ & [1.003, 1.028] & 0.0316 & 0.0571 & 1.05 & 0.0005 & 2.2 \\
schubfach:go-vs-rs & api & canonicalize & deep & schubfach & schubfach-rs & schubfach-rs$^*$ & 1.27$\times$ & [1.252, 1.281] & 0.0002 & 0.0004 & 18.76 & 0.0005 & 2.0 \\
bd:go-vs-rs & api & canonicalize & deep & json-canon & json-canon-rs & json-canon-rs$^*$ & 1.22$\times$ & [1.203, 1.228] & 0.0002 & 0.0004 & 16.66 & 0.0005 & 2.2 \\
ryu:go-vs-rs & api & canonicalize & deep & ryu & ryu-rs & ryu-rs$^*$ & 1.27$\times$ & [1.215, 1.298] & 0.0002 & 0.0004 & 7.81 & 0.0014 & 6.1 \\
dragonbox:go-vs-rs & api & canonicalize & deep & dragonbox & dragonbox-rs & dragonbox-rs$^*$ & 1.22$\times$ & [1.215, 1.235] & 0.0002 & 0.0004 & 21.92 & 0.0004 & 1.5 \\
go:schubfach-vs-bd & api & canonicalize & escaped-key-order & schubfach & json-canon & schubfach$^*$ & 1.59$\times$ & [1.571, 1.609] & 0.0002 & 0.0004 & -33.61 & 0.0000 & 2.0 \\
go:schubfach-vs-ryu & api & canonicalize & escaped-key-order & schubfach & ryu & ryu & 1.01$\times$ & [0.998, 1.019] & 0.2178 & 0.3096 & 0.58 & 0.0000 & 0.8 \\
go:schubfach-vs-dragonbox & api & canonicalize & escaped-key-order & schubfach & dragonbox & schubfach & 1.00$\times$ & [0.990, 1.033] & 0.7473 & 0.8126 & -0.20 & 0.0000 & 2.0 \\
go:ryu-vs-dragonbox & api & canonicalize & escaped-key-order & ryu & dragonbox & ryu & 1.01$\times$ & [1.000, 1.040] & 0.2006 & 0.2893 & -0.56 & 0.0000 & 0.8 \\
go:ryu-vs-bd & api & canonicalize & escaped-key-order & ryu & json-canon & ryu$^*$ & 1.60$\times$ & [1.591, 1.617] & 0.0002 & 0.0004 & -41.45 & 0.0000 & 0.8 \\
go:dragonbox-vs-bd & api & canonicalize & escaped-key-order & dragonbox & json-canon & dragonbox$^*$ & 1.58$\times$ & [1.541, 1.605] & 0.0002 & 0.0004 & -24.11 & 0.0000 & 3.6 \\
rs:schubfach-vs-bd & api & canonicalize & escaped-key-order & schubfach-rs & json-canon-rs & schubfach-rs$^*$ & 2.55$\times$ & [2.505, 2.601] & 0.0002 & 0.0004 & -38.00 & 0.0000 & 3.0 \\
rs:schubfach-vs-ryu & api & canonicalize & escaped-key-order & schubfach-rs & ryu-rs & schubfach-rs$^*$ & 1.03$\times$ & [1.008, 1.055] & 0.0248 & 0.0455 & -1.12 & 0.0000 & 3.0 \\
rs:schubfach-vs-dragonbox & api & canonicalize & escaped-key-order & schubfach-rs & dragonbox-rs & schubfach-rs$^*$ & 1.04$\times$ & [1.013, 1.064] & 0.0118 & 0.0224 & -1.22 & 0.0000 & 3.0 \\
rs:ryu-vs-dragonbox & api & canonicalize & escaped-key-order & ryu-rs & dragonbox-rs & ryu-rs & 1.00$\times$ & [0.979, 1.033] & 0.7680 & 0.8282 & -0.14 & 0.0000 & 3.9 \\
rs:ryu-vs-bd & api & canonicalize & escaped-key-order & ryu-rs & json-canon-rs & ryu-rs$^*$ & 2.47$\times$ & [2.421, 2.527] & 0.0002 & 0.0004 & -35.04 & 0.0000 & 3.9 \\
rs:dragonbox-vs-bd & api & canonicalize & escaped-key-order & dragonbox-rs & json-canon-rs & dragonbox-rs$^*$ & 2.46$\times$ & [2.401, 2.516] & 0.0002 & 0.0004 & -34.23 & 0.0000 & 4.1 \\
schubfach:go-vs-rs & api & canonicalize & escaped-key-order & schubfach & schubfach-rs & schubfach-rs$^*$ & 2.20$\times$ & [2.162, 2.234] & 0.0002 & 0.0004 & 40.23 & 0.0000 & 3.0 \\
bd:go-vs-rs & api & canonicalize & escaped-key-order & json-canon & json-canon-rs & json-canon-rs$^*$ & 1.37$\times$ & [1.351, 1.390] & 0.0002 & 0.0004 & 19.73 & 0.0000 & 2.6 \\
ryu:go-vs-rs & api & canonicalize & escaped-key-order & ryu & ryu-rs & ryu-rs$^*$ & 2.11$\times$ & [2.076, 2.154] & 0.0002 & 0.0004 & 46.50 & 0.0000 & 3.9 \\
dragonbox:go-vs-rs & api & canonicalize & escaped-key-order & dragonbox & dragonbox-rs & dragonbox-rs$^*$ & 2.13$\times$ & [2.080, 2.189] & 0.0002 & 0.0004 & 23.20 & 0.0000 & 4.1 \\
go:schubfach-vs-bd & api & canonicalize & large & schubfach & json-canon & schubfach$^*$ & 1.77$\times$ & [1.761, 1.774] & 0.0002 & 0.0004 & -130.43 & 0.0595 & 0.6 \\
go:schubfach-vs-ryu & api & canonicalize & large & schubfach & ryu & ryu & 1.01$\times$ & [1.000, 1.009] & 0.0404 & 0.0709 & 1.00 & 0.0710 & 0.8 \\
go:schubfach-vs-dragonbox & api & canonicalize & large & schubfach & dragonbox & dragonbox & 1.01$\times$ & [1.000, 1.011] & 0.0954 & 0.1541 & 0.78 & 0.0998 & 1.1 \\
go:ryu-vs-dragonbox & api & canonicalize & large & ryu & dragonbox & ryu & 1.00$\times$ & [0.993, 1.006] & 0.9638 & 0.9751 & -0.02 & 0.0710 & 0.8 \\
go:ryu-vs-bd & api & canonicalize & large & ryu & json-canon & ryu$^*$ & 1.78$\times$ & [1.767, 1.783] & 0.0002 & 0.0004 & -122.25 & 0.0710 & 0.8 \\
go:dragonbox-vs-bd & api & canonicalize & large & dragonbox & json-canon & dragonbox$^*$ & 1.78$\times$ & [1.767, 1.786] & 0.0002 & 0.0004 & -101.90 & 0.0998 & 1.1 \\
rs:schubfach-vs-bd & api & canonicalize & large & schubfach-rs & json-canon-rs & schubfach-rs$^*$ & 3.28$\times$ & [3.181, 3.344] & 0.0002 & 0.0004 & -48.92 & 0.2203 & 4.7 \\
rs:schubfach-vs-ryu & api & canonicalize & large & schubfach-rs & ryu-rs & schubfach-rs & 1.01$\times$ & [0.974, 1.048] & 0.6511 & 0.7394 & -0.16 & 0.2203 & 4.7 \\
rs:schubfach-vs-dragonbox & api & canonicalize & large & schubfach-rs & dragonbox-rs & dragonbox-rs & 1.01$\times$ & [0.985, 1.042] & 0.3257 & 0.4271 & 0.40 & 0.1710 & 3.7 \\
rs:ryu-vs-dragonbox & api & canonicalize & large & ryu-rs & dragonbox-rs & dragonbox-rs & 1.02$\times$ & [0.990, 1.062] & 0.2484 & 0.3485 & 0.50 & 0.1710 & 3.7 \\
rs:ryu-vs-bd & api & canonicalize & large & ryu-rs & json-canon-rs & ryu-rs$^*$ & 3.26$\times$ & [3.135, 3.338] & 0.0002 & 0.0004 & -43.69 & 0.2911 & 6.2 \\
rs:dragonbox-vs-bd & api & canonicalize & large & dragonbox-rs & json-canon-rs & dragonbox-rs$^*$ & 3.33$\times$ & [3.228, 3.373] & 0.0002 & 0.0004 & -52.80 & 0.1710 & 3.7 \\
schubfach:go-vs-rs & api & canonicalize & large & schubfach & schubfach-rs & schubfach-rs$^*$ & 2.04$\times$ & [1.979, 2.074] & 0.0002 & 0.0004 & 37.40 & 0.2203 & 4.7 \\
bd:go-vs-rs & api & canonicalize & large & json-canon & json-canon-rs & json-canon-rs$^*$ & 1.10$\times$ & [1.087, 1.108] & 0.0002 & 0.0004 & 8.12 & 0.3138 & 2.1 \\
ryu:go-vs-rs & api & canonicalize & large & ryu & ryu-rs & ryu-rs$^*$ & 2.01$\times$ & [1.937, 2.059] & 0.0002 & 0.0004 & 27.98 & 0.2911 & 6.2 \\
dragonbox:go-vs-rs & api & canonicalize & large & dragonbox & dragonbox-rs & dragonbox-rs$^*$ & 2.05$\times$ & [1.993, 2.079] & 0.0002 & 0.0004 & 43.21 & 0.1710 & 3.7 \\
go:schubfach-vs-bd & api & canonicalize & long-string & schubfach & json-canon & schubfach & 1.01$\times$ & [0.996, 1.021] & 0.2975 & 0.3977 & -0.49 & 0.0011 & 1.8 \\
go:schubfach-vs-ryu & api & canonicalize & long-string & schubfach & ryu & ryu & 1.00$\times$ & [0.987, 1.014] & 0.9742 & 0.9795 & 0.01 & 0.0013 & 2.2 \\
go:schubfach-vs-dragonbox & api & canonicalize & long-string & schubfach & dragonbox & schubfach & 1.00$\times$ & [0.992, 1.015] & 0.6451 & 0.7335 & -0.22 & 0.0011 & 1.8 \\
go:ryu-vs-dragonbox & api & canonicalize & long-string & ryu & dragonbox & ryu & 1.00$\times$ & [0.991, 1.017] & 0.6589 & 0.7438 & -0.21 & 0.0013 & 2.2 \\
go:ryu-vs-bd & api & canonicalize & long-string & ryu & json-canon & ryu & 1.01$\times$ & [0.995, 1.022] & 0.3197 & 0.4214 & -0.46 & 0.0013 & 2.2 \\
go:dragonbox-vs-bd & api & canonicalize & long-string & dragonbox & json-canon & dragonbox & 1.00$\times$ & [0.993, 1.016] & 0.5297 & 0.6294 & -0.28 & 0.0011 & 1.7 \\
rs:schubfach-vs-bd & api & canonicalize & long-string & schubfach-rs & json-canon-rs & schubfach-rs$^*$ & 1.05$\times$ & [1.039, 1.051] & 0.0002 & 0.0004 & -6.92 & 0.0000 & 0.2 \\
rs:schubfach-vs-ryu & api & canonicalize & long-string & schubfach-rs & ryu-rs & schubfach-rs$^*$ & 1.02$\times$ & [1.020, 1.022] & 0.0002 & 0.0004 & -13.18 & 0.0000 & 0.2 \\
rs:schubfach-vs-dragonbox & api & canonicalize & long-string & schubfach-rs & dragonbox-rs & schubfach-rs$^*$ & 1.02$\times$ & [1.023, 1.026] & 0.0002 & 0.0004 & -15.28 & 0.0000 & 0.2 \\
rs:ryu-vs-dragonbox & api & canonicalize & long-string & ryu-rs & dragonbox-rs & ryu-rs$^*$ & 1.00$\times$ & [1.002, 1.005] & 0.0004 & 0.0008 & -2.11 & 0.0000 & 0.2 \\
rs:ryu-vs-bd & api & canonicalize & long-string & ryu-rs & json-canon-rs & ryu-rs$^*$ & 1.03$\times$ & [1.018, 1.030] & 0.0002 & 0.0004 & -3.79 & 0.0000 & 0.2 \\
rs:dragonbox-vs-bd & api & canonicalize & long-string & dragonbox-rs & json-canon-rs & dragonbox-rs$^*$ & 1.02$\times$ & [1.014, 1.026] & 0.0002 & 0.0004 & -3.28 & 0.0000 & 0.2 \\
schubfach:go-vs-rs & api & canonicalize & long-string & schubfach & schubfach-rs & schubfach-rs$^*$ & 4.50$\times$ & [4.456, 4.536] & 0.0002 & 0.0004 & 74.61 & 0.0000 & 0.2 \\
bd:go-vs-rs & api & canonicalize & long-string & json-canon & json-canon-rs & json-canon-rs$^*$ & 4.33$\times$ & [4.293, 4.378] & 0.0002 & 0.0004 & 77.40 & 0.0002 & 1.1 \\
ryu:go-vs-rs & api & canonicalize & long-string & ryu & ryu-rs & ryu-rs$^*$ & 4.41$\times$ & [4.356, 4.448] & 0.0002 & 0.0004 & 62.47 & 0.0000 & 0.2 \\
dragonbox:go-vs-rs & api & canonicalize & long-string & dragonbox & dragonbox-rs & dragonbox-rs$^*$ & 4.41$\times$ & [4.372, 4.445] & 0.0002 & 0.0004 & 78.27 & 0.0000 & 0.2 \\
go:schubfach-vs-bd & api & canonicalize & medium & schubfach & json-canon & schubfach$^*$ & 1.75$\times$ & [1.743, 1.766] & 0.0002 & 0.0004 & -75.66 & 0.0034 & 1.2 \\
go:schubfach-vs-ryu & api & canonicalize & medium & schubfach & ryu & ryu & 1.01$\times$ & [1.000, 1.013] & 0.0642 & 0.1093 & 0.90 & 0.0012 & 0.4 \\
go:schubfach-vs-dragonbox & api & canonicalize & medium & schubfach & dragonbox & dragonbox$^*$ & 1.01$\times$ & [1.003, 1.016] & 0.0148 & 0.0277 & 1.19 & 0.0017 & 0.6 \\
go:ryu-vs-dragonbox & api & canonicalize & medium & ryu & dragonbox & dragonbox & 1.00$\times$ & [0.999, 1.006] & 0.1766 & 0.2602 & 0.64 & 0.0017 & 0.6 \\
go:ryu-vs-bd & api & canonicalize & medium & ryu & json-canon & ryu$^*$ & 1.77$\times$ & [1.759, 1.773] & 0.0002 & 0.0004 & -101.53 & 0.0012 & 0.4 \\
go:dragonbox-vs-bd & api & canonicalize & medium & dragonbox & json-canon & dragonbox$^*$ & 1.77$\times$ & [1.763, 1.779] & 0.0002 & 0.0004 & -95.92 & 0.0017 & 0.6 \\
rs:schubfach-vs-bd & api & canonicalize & medium & schubfach-rs & json-canon-rs & schubfach-rs$^*$ & 2.91$\times$ & [2.484, 3.068] & 0.0002 & 0.0004 & -17.82 & 0.0265 & 18.1 \\
rs:schubfach-vs-ryu & api & canonicalize & medium & schubfach-rs & ryu-rs & ryu-rs & 1.03$\times$ & [0.974, 1.176] & 0.9924 & 0.9949 & 0.24 & 0.0027 & 1.9 \\
rs:schubfach-vs-dragonbox & api & canonicalize & medium & schubfach-rs & dragonbox-rs & schubfach-rs & 1.02$\times$ & [0.916, 1.178] & 0.5131 & 0.6142 & -0.11 & 0.0265 & 18.1 \\
rs:ryu-vs-dragonbox & api & canonicalize & medium & ryu-rs & dragonbox-rs & ryu-rs & 1.04$\times$ & [0.986, 1.216] & 0.9884 & 0.9922 & -0.36 & 0.0027 & 1.9 \\
rs:ryu-vs-bd & api & canonicalize & medium & ryu-rs & json-canon-rs & ryu-rs$^*$ & 2.99$\times$ & [2.943, 3.020] & 0.0002 & 0.0004 & -57.35 & 0.0027 & 1.9 \\
rs:dragonbox-vs-bd & api & canonicalize & medium & dragonbox-rs & json-canon-rs & dragonbox-rs$^*$ & 2.86$\times$ & [2.447, 3.032] & 0.0002 & 0.0004 & -15.29 & 0.0309 & 20.8 \\
schubfach:go-vs-rs & api & canonicalize & medium & schubfach & schubfach-rs & schubfach-rs$^*$ & 1.86$\times$ & [1.577, 1.959] & 0.0002 & 0.0004 & 8.34 & 0.0265 & 18.1 \\
bd:go-vs-rs & api & canonicalize & medium & json-canon & json-canon-rs & json-canon-rs$^*$ & 1.12$\times$ & [1.110, 1.132] & 0.0002 & 0.0004 & 10.14 & 0.0083 & 2.0 \\
ryu:go-vs-rs & api & canonicalize & medium & ryu & ryu-rs & ryu-rs$^*$ & 1.90$\times$ & [1.867, 1.908] & 0.0002 & 0.0004 & 76.72 & 0.0027 & 1.9 \\
dragonbox:go-vs-rs & api & canonicalize & medium & dragonbox & dragonbox-rs & dragonbox-rs$^*$ & 1.81$\times$ & [1.497, 1.915] & 0.0002 & 0.0004 & 6.89 & 0.0309 & 20.8 \\
go:schubfach-vs-bd & api & canonicalize & mixed-prod & schubfach & json-canon & schubfach$^*$ & 1.36$\times$ & [1.351, 1.370] & 0.0002 & 0.0004 & -39.25 & 0.0013 & 1.2 \\
go:schubfach-vs-ryu & api & canonicalize & mixed-prod & schubfach & ryu & schubfach & 1.00$\times$ & [0.997, 1.011] & 0.2725 & 0.3739 & -0.51 & 0.0013 & 1.2 \\
go:schubfach-vs-dragonbox & api & canonicalize & mixed-prod & schubfach & dragonbox & schubfach & 1.00$\times$ & [0.995, 1.008] & 0.5385 & 0.6369 & -0.27 & 0.0013 & 1.2 \\
go:ryu-vs-dragonbox & api & canonicalize & mixed-prod & ryu & dragonbox & dragonbox & 1.00$\times$ & [0.996, 1.009] & 0.5061 & 0.6077 & 0.31 & 0.0009 & 0.8 \\
go:ryu-vs-bd & api & canonicalize & mixed-prod & ryu & json-canon & ryu$^*$ & 1.35$\times$ & [1.347, 1.364] & 0.0002 & 0.0004 & -42.30 & 0.0011 & 1.0 \\
go:dragonbox-vs-bd & api & canonicalize & mixed-prod & dragonbox & json-canon & dragonbox$^*$ & 1.36$\times$ & [1.351, 1.366] & 0.0002 & 0.0004 & -46.68 & 0.0009 & 0.8 \\
rs:schubfach-vs-bd & api & canonicalize & mixed-prod & schubfach-rs & json-canon-rs & schubfach-rs$^*$ & 1.78$\times$ & [1.738, 1.802] & 0.0002 & 0.0004 & -33.60 & 0.0018 & 3.3 \\
rs:schubfach-vs-ryu & api & canonicalize & mixed-prod & schubfach-rs & ryu-rs & schubfach-rs & 1.02$\times$ & [0.997, 1.041] & 0.1238 & 0.1937 & -0.74 & 0.0018 & 3.3 \\
rs:schubfach-vs-dragonbox & api & canonicalize & mixed-prod & schubfach-rs & dragonbox-rs & dragonbox-rs$^*$ & 1.03$\times$ & [1.013, 1.051] & 0.0048 & 0.0094 & 1.24 & 0.0010 & 1.8 \\
rs:ryu-vs-dragonbox & api & canonicalize & mixed-prod & ryu-rs & dragonbox-rs & dragonbox-rs$^*$ & 1.05$\times$ & [1.032, 1.068] & 0.0002 & 0.0004 & 2.18 & 0.0010 & 1.8 \\
rs:ryu-vs-bd & api & canonicalize & mixed-prod & ryu-rs & json-canon-rs & ryu-rs$^*$ & 1.75$\times$ & [1.710, 1.768] & 0.0002 & 0.0004 & -32.91 & 0.0018 & 3.2 \\
rs:dragonbox-vs-bd & api & canonicalize & mixed-prod & dragonbox-rs & json-canon-rs & dragonbox-rs$^*$ & 1.83$\times$ & [1.807, 1.845] & 0.0002 & 0.0004 & -47.89 & 0.0010 & 1.8 \\
schubfach:go-vs-rs & api & canonicalize & mixed-prod & schubfach & schubfach-rs & schubfach-rs$^*$ & 1.94$\times$ & [1.900, 1.970] & 0.0002 & 0.0004 & 41.05 & 0.0018 & 3.3 \\
bd:go-vs-rs & api & canonicalize & mixed-prod & json-canon & json-canon-rs & json-canon-rs$^*$ & 1.49$\times$ & [1.476, 1.498] & 0.0002 & 0.0004 & 47.86 & 0.0013 & 1.4 \\
ryu:go-vs-rs & api & canonicalize & mixed-prod & ryu & ryu-rs & ryu-rs$^*$ & 1.92$\times$ & [1.877, 1.941] & 0.0002 & 0.0004 & 42.94 & 0.0018 & 3.2 \\
dragonbox:go-vs-rs & api & canonicalize & mixed-prod & dragonbox & dragonbox-rs & dragonbox-rs$^*$ & 2.00$\times$ & [1.980, 2.018] & 0.0002 & 0.0004 & 73.12 & 0.0010 & 1.8 \\
go:schubfach-vs-bd & api & canonicalize & nested-mixed & schubfach & json-canon & schubfach$^*$ & 1.42$\times$ & [1.408, 1.424] & 0.0002 & 0.0004 & -48.55 & 0.0000 & 0.8 \\
go:schubfach-vs-ryu & api & canonicalize & nested-mixed & schubfach & ryu & schubfach & 1.01$\times$ & [1.002, 1.013] & 0.0328 & 0.0589 & -1.05 & 0.0000 & 0.8 \\
go:schubfach-vs-dragonbox & api & canonicalize & nested-mixed & schubfach & dragonbox & schubfach & 1.00$\times$ & [0.996, 1.007] & 0.5915 & 0.6877 & -0.25 & 0.0000 & 0.8 \\
go:ryu-vs-dragonbox & api & canonicalize & nested-mixed & ryu & dragonbox & dragonbox & 1.01$\times$ & [1.000, 1.012] & 0.1224 & 0.1922 & 0.74 & 0.0000 & 0.9 \\
go:ryu-vs-bd & api & canonicalize & nested-mixed & ryu & json-canon & ryu$^*$ & 1.41$\times$ & [1.398, 1.414] & 0.0002 & 0.0004 & -46.65 & 0.0000 & 0.9 \\
go:dragonbox-vs-bd & api & canonicalize & nested-mixed & dragonbox & json-canon & dragonbox$^*$ & 1.41$\times$ & [1.405, 1.423] & 0.0002 & 0.0004 & -46.24 & 0.0000 & 0.9 \\
rs:schubfach-vs-bd & api & canonicalize & nested-mixed & schubfach-rs & json-canon-rs & schubfach-rs$^*$ & 1.97$\times$ & [1.923, 2.010] & 0.0002 & 0.0004 & -33.02 & 0.0000 & 4.4 \\
rs:schubfach-vs-ryu & api & canonicalize & nested-mixed & schubfach-rs & ryu-rs & schubfach-rs & 1.02$\times$ & [0.988, 1.052] & 0.2519 & 0.3509 & -0.54 & 0.0000 & 4.4 \\
rs:schubfach-vs-dragonbox & api & canonicalize & nested-mixed & schubfach-rs & dragonbox-rs & schubfach-rs & 1.03$\times$ & [0.999, 1.052] & 0.0466 & 0.0808 & -0.95 & 0.0000 & 4.4 \\
rs:ryu-vs-dragonbox & api & canonicalize & nested-mixed & ryu-rs & dragonbox-rs & ryu-rs & 1.01$\times$ & [0.979, 1.035] & 0.5805 & 0.6773 & -0.25 & 0.0000 & 5.1 \\
rs:ryu-vs-bd & api & canonicalize & nested-mixed & ryu-rs & json-canon-rs & ryu-rs$^*$ & 1.93$\times$ & [1.880, 1.977] & 0.0002 & 0.0004 & -28.49 & 0.0000 & 5.1 \\
rs:dragonbox-vs-bd & api & canonicalize & nested-mixed & dragonbox-rs & json-canon-rs & dragonbox-rs$^*$ & 1.92$\times$ & [1.889, 1.950] & 0.0002 & 0.0004 & -39.12 & 0.0000 & 3.1 \\
schubfach:go-vs-rs & api & canonicalize & nested-mixed & schubfach & schubfach-rs & schubfach-rs$^*$ & 2.18$\times$ & [2.128, 2.221] & 0.0002 & 0.0004 & 44.21 & 0.0000 & 4.4 \\
bd:go-vs-rs & api & canonicalize & nested-mixed & json-canon & json-canon-rs & json-canon-rs$^*$ & 1.57$\times$ & [1.556, 1.580] & 0.0002 & 0.0004 & 49.86 & 0.0000 & 1.4 \\
ryu:go-vs-rs & api & canonicalize & nested-mixed & ryu & ryu-rs & ryu-rs$^*$ & 2.15$\times$ & [2.097, 2.202] & 0.0002 & 0.0004 & 37.60 & 0.0000 & 5.1 \\
dragonbox:go-vs-rs & api & canonicalize & nested-mixed & dragonbox & dragonbox-rs & dragonbox-rs$^*$ & 2.12$\times$ & [2.096, 2.160] & 0.0002 & 0.0004 & 53.86 & 0.0000 & 3.1 \\
go:schubfach-vs-bd & api & canonicalize & number-heavy & schubfach & json-canon & schubfach$^*$ & 1.57$\times$ & [1.552, 1.579] & 0.0002 & 0.0004 & -43.36 & 0.0002 & 1.5 \\
go:schubfach-vs-ryu & api & canonicalize & number-heavy & schubfach & ryu & schubfach & 1.00$\times$ & [0.993, 1.011] & 0.6015 & 0.6962 & -0.25 & 0.0002 & 1.5 \\
go:schubfach-vs-dragonbox & api & canonicalize & number-heavy & schubfach & dragonbox & dragonbox & 1.01$\times$ & [0.999, 1.016] & 0.1844 & 0.2696 & 0.62 & 0.0001 & 1.1 \\
go:ryu-vs-dragonbox & api & canonicalize & number-heavy & ryu & dragonbox & dragonbox & 1.01$\times$ & [1.002, 1.017] & 0.0472 & 0.0815 & 0.98 & 0.0001 & 1.1 \\
go:ryu-vs-bd & api & canonicalize & number-heavy & ryu & json-canon & ryu$^*$ & 1.56$\times$ & [1.550, 1.575] & 0.0002 & 0.0004 & -45.51 & 0.0001 & 1.3 \\
go:dragonbox-vs-bd & api & canonicalize & number-heavy & dragonbox & json-canon & dragonbox$^*$ & 1.58$\times$ & [1.565, 1.588] & 0.0002 & 0.0004 & -48.70 & 0.0001 & 1.1 \\
rs:schubfach-vs-bd & api & canonicalize & number-heavy & schubfach-rs & json-canon-rs & schubfach-rs$^*$ & 8.23$\times$ & [8.106, 8.340] & 0.0004 & 0.0008 & -106.99 & 0.0000 & 2.7 \\
rs:schubfach-vs-ryu & api & canonicalize & number-heavy & schubfach-rs & ryu-rs & schubfach-rs & 1.03$\times$ & [1.004, 1.054] & 0.0338 & 0.0604 & -1.01 & 0.0000 & 2.7 \\
rs:schubfach-vs-dragonbox & api & canonicalize & number-heavy & schubfach-rs & dragonbox-rs & dragonbox-rs & 1.00$\times$ & [0.986, 1.021] & 0.7640 & 0.8266 & 0.13 & 0.0000 & 2.8 \\
rs:ryu-vs-dragonbox & api & canonicalize & number-heavy & ryu-rs & dragonbox-rs & dragonbox-rs$^*$ & 1.03$\times$ & [1.009, 1.059] & 0.0234 & 0.0430 & 1.10 & 0.0000 & 2.8 \\
rs:ryu-vs-bd & api & canonicalize & number-heavy & ryu-rs & json-canon-rs & ryu-rs$^*$ & 7.99$\times$ & [7.815, 8.161] & 0.0002 & 0.0004 & -101.80 & 0.0001 & 4.4 \\
rs:dragonbox-vs-bd & api & canonicalize & number-heavy & dragonbox-rs & json-canon-rs & dragonbox-rs$^*$ & 8.25$\times$ & [8.129, 8.371] & 0.0002 & 0.0004 & -106.87 & 0.0000 & 2.8 \\
schubfach:go-vs-rs & api & canonicalize & number-heavy & schubfach & schubfach-rs & schubfach-rs$^*$ & 9.65$\times$ & [9.505, 9.784] & 0.0002 & 0.0004 & 106.41 & 0.0000 & 2.7 \\
bd:go-vs-rs & api & canonicalize & number-heavy & json-canon & json-canon-rs & json-canon-rs$^*$ & 1.84$\times$ & [1.821, 1.852] & 0.0002 & 0.0004 & 58.48 & 0.0001 & 1.4 \\
ryu:go-vs-rs & api & canonicalize & number-heavy & ryu & ryu-rs & ryu-rs$^*$ & 9.39$\times$ & [9.190, 9.605] & 0.0002 & 0.0004 & 117.25 & 0.0001 & 4.4 \\
dragonbox:go-vs-rs & api & canonicalize & number-heavy & dragonbox & dragonbox-rs & dragonbox-rs$^*$ & 9.61$\times$ & [9.472, 9.743] & 0.0002 & 0.0004 & 142.99 & 0.0000 & 2.8 \\
go:schubfach-vs-bd & api & canonicalize & numeric-boundary & schubfach & json-canon & schubfach$^*$ & 1.50$\times$ & [1.487, 1.510] & 0.0002 & 0.0004 & -41.30 & 0.0001 & 1.0 \\
go:schubfach-vs-ryu & api & canonicalize & numeric-boundary & schubfach & ryu & ryu & 1.00$\times$ & [0.991, 1.009] & 0.8524 & 0.8939 & 0.09 & 0.0002 & 1.6 \\
go:schubfach-vs-dragonbox & api & canonicalize & numeric-boundary & schubfach & dragonbox & dragonbox & 1.00$\times$ & [0.998, 1.010] & 0.2561 & 0.3555 & 0.54 & 0.0001 & 0.8 \\
go:ryu-vs-dragonbox & api & canonicalize & numeric-boundary & ryu & dragonbox & dragonbox & 1.00$\times$ & [0.996, 1.013] & 0.5421 & 0.6400 & 0.30 & 0.0001 & 0.8 \\
go:ryu-vs-bd & api & canonicalize & numeric-boundary & ryu & json-canon & ryu$^*$ & 1.50$\times$ & [1.485, 1.513] & 0.0002 & 0.0004 & -36.11 & 0.0002 & 1.6 \\
go:dragonbox-vs-bd & api & canonicalize & numeric-boundary & dragonbox & json-canon & dragonbox$^*$ & 1.50$\times$ & [1.495, 1.515] & 0.0002 & 0.0004 & -43.81 & 0.0001 & 0.8 \\
rs:schubfach-vs-bd & api & canonicalize & numeric-boundary & schubfach-rs & json-canon-rs & schubfach-rs$^*$ & 9.33$\times$ & [9.163, 9.542] & 0.0002 & 0.0004 & -178.05 & 0.0000 & 4.1 \\
rs:schubfach-vs-ryu & api & canonicalize & numeric-boundary & schubfach-rs & ryu-rs & schubfach-rs & 1.02$\times$ & [0.992, 1.058] & 0.2767 & 0.3770 & -0.52 & 0.0000 & 4.1 \\
rs:schubfach-vs-dragonbox & api & canonicalize & numeric-boundary & schubfach-rs & dragonbox-rs & dragonbox-rs & 1.01$\times$ & [0.982, 1.038] & 0.3353 & 0.4379 & 0.44 & 0.0000 & 4.0 \\
rs:ryu-vs-dragonbox & api & canonicalize & numeric-boundary & ryu-rs & dragonbox-rs & dragonbox-rs & 1.03$\times$ & [1.004, 1.070] & 0.0556 & 0.0949 & 0.89 & 0.0000 & 4.0 \\
rs:ryu-vs-bd & api & canonicalize & numeric-boundary & ryu-rs & json-canon-rs & ryu-rs$^*$ & 9.15$\times$ & [8.861, 9.348] & 0.0002 & 0.0004 & -161.50 & 0.0000 & 5.5 \\
rs:dragonbox-vs-bd & api & canonicalize & numeric-boundary & dragonbox-rs & json-canon-rs & dragonbox-rs$^*$ & 9.46$\times$ & [9.239, 9.614] & 0.0002 & 0.0004 & -179.88 & 0.0000 & 4.0 \\
schubfach:go-vs-rs & api & canonicalize & numeric-boundary & schubfach & schubfach-rs & schubfach-rs$^*$ & 12.98$\times$ & [12.744, 13.263] & 0.0002 & 0.0004 & 154.21 & 0.0000 & 4.1 \\
bd:go-vs-rs & api & canonicalize & numeric-boundary & json-canon & json-canon-rs & json-canon-rs$^*$ & 2.08$\times$ & [2.071, 2.099] & 0.0002 & 0.0004 & 70.32 & 0.0001 & 0.8 \\
ryu:go-vs-rs & api & canonicalize & numeric-boundary & ryu & ryu-rs & ryu-rs$^*$ & 12.71$\times$ & [12.314, 13.008] & 0.0002 & 0.0004 & 100.35 & 0.0000 & 5.5 \\
dragonbox:go-vs-rs & api & canonicalize & numeric-boundary & dragonbox & dragonbox-rs & dragonbox-rs$^*$ & 13.11$\times$ & [12.813, 13.320] & 0.0002 & 0.0004 & 199.14 & 0.0000 & 4.0 \\
go:schubfach-vs-bd & api & canonicalize & rfc-key-sorting & schubfach & json-canon & schubfach$^*$ & 1.03$\times$ & [1.022, 1.044] & 0.0002 & 0.0004 & -2.42 & 0.0000 & 0.8 \\
go:schubfach-vs-ryu & api & canonicalize & rfc-key-sorting & schubfach & ryu & schubfach & 1.01$\times$ & [0.998, 1.012] & 0.1656 & 0.2491 & -0.64 & 0.0000 & 0.8 \\
go:schubfach-vs-dragonbox & api & canonicalize & rfc-key-sorting & schubfach & dragonbox & schubfach & 1.00$\times$ & [0.997, 1.008] & 0.3959 & 0.5033 & -0.39 & 0.0000 & 0.8 \\
go:ryu-vs-dragonbox & api & canonicalize & rfc-key-sorting & ryu & dragonbox & dragonbox & 1.00$\times$ & [0.995, 1.010] & 0.5119 & 0.6133 & 0.30 & 0.0000 & 0.9 \\
go:ryu-vs-bd & api & canonicalize & rfc-key-sorting & ryu & json-canon & ryu$^*$ & 1.03$\times$ & [1.015, 1.039] & 0.0002 & 0.0004 & -1.84 & 0.0000 & 1.2 \\
go:dragonbox-vs-bd & api & canonicalize & rfc-key-sorting & dragonbox & json-canon & dragonbox$^*$ & 1.03$\times$ & [1.019, 1.042] & 0.0002 & 0.0004 & -2.15 & 0.0000 & 0.9 \\
rs:schubfach-vs-bd & api & canonicalize & rfc-key-sorting & schubfach-rs & json-canon-rs & json-canon-rs & 1.02$\times$ & [0.985, 1.060] & 0.2727 & 0.3739 & 0.50 & 0.0000 & 3.7 \\
rs:schubfach-vs-ryu & api & canonicalize & rfc-key-sorting & schubfach-rs & ryu-rs & schubfach-rs & 1.03$\times$ & [0.990, 1.071] & 0.1932 & 0.2805 & -0.62 & 0.0001 & 7.0 \\
rs:schubfach-vs-dragonbox & api & canonicalize & rfc-key-sorting & schubfach-rs & dragonbox-rs & schubfach-rs & 1.00$\times$ & [0.964, 1.044] & 0.9438 & 0.9624 & -0.03 & 0.0001 & 7.0 \\
rs:ryu-vs-dragonbox & api & canonicalize & rfc-key-sorting & ryu-rs & dragonbox-rs & dragonbox-rs & 1.03$\times$ & [0.997, 1.059] & 0.1088 & 0.1736 & 0.75 & 0.0000 & 4.5 \\
rs:ryu-vs-bd & api & canonicalize & rfc-key-sorting & ryu-rs & json-canon-rs & json-canon-rs$^*$ & 1.05$\times$ & [1.025, 1.083] & 0.0032 & 0.0064 & 1.53 & 0.0000 & 3.7 \\
rs:dragonbox-vs-bd & api & canonicalize & rfc-key-sorting & dragonbox-rs & json-canon-rs & json-canon-rs & 1.02$\times$ & [0.999, 1.055] & 0.1342 & 0.2087 & 0.72 & 0.0000 & 3.7 \\
schubfach:go-vs-rs & api & canonicalize & rfc-key-sorting & schubfach & schubfach-rs & schubfach-rs$^*$ & 2.18$\times$ & [2.113, 2.255] & 0.0002 & 0.0004 & 29.28 & 0.0001 & 7.0 \\
bd:go-vs-rs & api & canonicalize & rfc-key-sorting & json-canon & json-canon-rs & json-canon-rs$^*$ & 2.30$\times$ & [2.254, 2.348] & 0.0002 & 0.0004 & 38.30 & 0.0000 & 3.7 \\
ryu:go-vs-rs & api & canonicalize & rfc-key-sorting & ryu & ryu-rs & ryu-rs$^*$ & 2.13$\times$ & [2.079, 2.175] & 0.0002 & 0.0004 & 38.16 & 0.0001 & 4.5 \\
dragonbox:go-vs-rs & api & canonicalize & rfc-key-sorting & dragonbox & dragonbox-rs & dragonbox-rs$^*$ & 2.18$\times$ & [2.130, 2.226] & 0.0002 & 0.0004 & 42.16 & 0.0000 & 4.5 \\
go:schubfach-vs-bd & api & canonicalize & small & schubfach & json-canon & schubfach$^*$ & 1.36$\times$ & [1.354, 1.365] & 0.0002 & 0.0004 & -58.36 & 0.0000 & 0.5 \\
go:schubfach-vs-ryu & api & canonicalize & small & schubfach & ryu & schubfach & 1.01$\times$ & [1.000, 1.012] & 0.1032 & 0.1650 & -0.77 & 0.0000 & 0.5 \\
go:schubfach-vs-dragonbox & api & canonicalize & small & schubfach & dragonbox & schubfach & 1.00$\times$ & [0.997, 1.009] & 0.5211 & 0.6215 & -0.29 & 0.0000 & 0.5 \\
go:ryu-vs-dragonbox & api & canonicalize & small & ryu & dragonbox & dragonbox & 1.00$\times$ & [0.995, 1.011] & 0.4711 & 0.5766 & 0.34 & 0.0000 & 1.3 \\
go:ryu-vs-bd & api & canonicalize & small & ryu & json-canon & ryu$^*$ & 1.35$\times$ & [1.343, 1.360] & 0.0002 & 0.0004 & -41.66 & 0.0000 & 1.1 \\
go:dragonbox-vs-bd & api & canonicalize & small & dragonbox & json-canon & dragonbox$^*$ & 1.36$\times$ & [1.346, 1.365] & 0.0002 & 0.0004 & -39.66 & 0.0000 & 1.3 \\
rs:schubfach-vs-bd & api & canonicalize & small & schubfach-rs & json-canon-rs & schubfach-rs$^*$ & 1.75$\times$ & [1.675, 1.807] & 0.0002 & 0.0004 & -13.36 & 0.0000 & 6.7 \\
rs:schubfach-vs-ryu & api & canonicalize & small & schubfach-rs & ryu-rs & schubfach-rs$^*$ & 1.06$\times$ & [1.020, 1.107] & 0.0102 & 0.0195 & -1.25 & 0.0000 & 6.7 \\
rs:schubfach-vs-dragonbox & api & canonicalize & small & schubfach-rs & dragonbox-rs & schubfach-rs & 1.04$\times$ & [0.983, 1.118] & 0.3127 & 0.4146 & -0.49 & 0.0000 & 6.7 \\
rs:ryu-vs-dragonbox & api & canonicalize & small & ryu-rs & dragonbox-rs & dragonbox-rs & 1.02$\times$ & [0.953, 1.079] & 0.4721 & 0.5766 & 0.33 & 0.0001 & 11.7 \\
rs:ryu-vs-bd & api & canonicalize & small & ryu-rs & json-canon-rs & ryu-rs$^*$ & 1.64$\times$ & [1.588, 1.696] & 0.0002 & 0.0004 & -12.65 & 0.0000 & 5.8 \\
rs:dragonbox-vs-bd & api & canonicalize & small & dragonbox-rs & json-canon-rs & dragonbox-rs$^*$ & 1.68$\times$ & [1.560, 1.761] & 0.0002 & 0.0004 & -8.85 & 0.0001 & 11.7 \\
schubfach:go-vs-rs & api & canonicalize & small & schubfach & schubfach-rs & schubfach-rs$^*$ & 1.89$\times$ & [1.821, 1.945] & 0.0002 & 0.0004 & 23.42 & 0.0000 & 6.7 \\
bd:go-vs-rs & api & canonicalize & small & json-canon & json-canon-rs & json-canon-rs$^*$ & 1.47$\times$ & [1.445, 1.504] & 0.0002 & 0.0004 & 19.44 & 0.0000 & 4.2 \\
ryu:go-vs-rs & api & canonicalize & small & ryu & ryu-rs & ryu-rs$^*$ & 1.79$\times$ & [1.735, 1.835] & 0.0002 & 0.0004 & 22.70 & 0.0000 & 5.8 \\
dragonbox:go-vs-rs & api & canonicalize & small & dragonbox & dragonbox-rs & dragonbox-rs$^*$ & 1.83$\times$ & [1.687, 1.905] & 0.0002 & 0.0004 & 12.28 & 0.0001 & 11.7 \\
go:schubfach-vs-bd & api & canonicalize & surrogate-pair & schubfach & json-canon & schubfach$^*$ & 1.02$\times$ & [1.011, 1.021] & 0.0002 & 0.0004 & -2.58 & 0.0000 & 0.5 \\
go:schubfach-vs-ryu & api & canonicalize & surrogate-pair & schubfach & ryu & schubfach & 1.00$\times$ & [0.998, 1.005] & 0.3067 & 0.4079 & -0.49 & 0.0000 & 0.5 \\
go:schubfach-vs-dragonbox & api & canonicalize & surrogate-pair & schubfach & dragonbox & schubfach & 1.01$\times$ & [1.002, 1.040] & 0.1610 & 0.2439 & -0.68 & 0.0000 & 0.5 \\
go:ryu-vs-dragonbox & api & canonicalize & surrogate-pair & ryu & dragonbox & ryu & 1.01$\times$ & [1.000, 1.038] & 0.2783 & 0.3779 & -0.59 & 0.0000 & 0.5 \\
go:ryu-vs-bd & api & canonicalize & surrogate-pair & ryu & json-canon & ryu$^*$ & 1.01$\times$ & [1.009, 1.019] & 0.0002 & 0.0004 & -2.21 & 0.0000 & 0.5 \\
go:dragonbox-vs-bd & api & canonicalize & surrogate-pair & dragonbox & json-canon & dragonbox & 1.00$\times$ & [0.977, 1.014] & 0.9568 & 0.9707 & -0.04 & 0.0000 & 3.6 \\
rs:schubfach-vs-bd & api & canonicalize & surrogate-pair & schubfach-rs & json-canon-rs & schubfach-rs & 1.01$\times$ & [0.933, 1.073] & 0.7303 & 0.7981 & -0.17 & 0.0000 & 11.7 \\
rs:schubfach-vs-ryu & api & canonicalize & surrogate-pair & schubfach-rs & ryu-rs & schubfach-rs & 1.00$\times$ & [0.924, 1.056] & 0.9508 & 0.9673 & -0.03 & 0.0000 & 11.7 \\
rs:schubfach-vs-dragonbox & api & canonicalize & surrogate-pair & schubfach-rs & dragonbox-rs & schubfach-rs & 1.07$\times$ & [0.985, 1.158] & 0.1510 & 0.2315 & -0.69 & 0.0000 & 11.7 \\
rs:ryu-vs-dragonbox & api & canonicalize & surrogate-pair & ryu-rs & dragonbox-rs & ryu-rs & 1.07$\times$ & [1.007, 1.157] & 0.0968 & 0.1560 & -0.77 & 0.0000 & 7.0 \\
rs:ryu-vs-bd & api & canonicalize & surrogate-pair & ryu-rs & json-canon-rs & ryu-rs & 1.01$\times$ & [0.962, 1.062] & 0.6887 & 0.7623 & -0.18 & 0.0000 & 7.0 \\
rs:dragonbox-vs-bd & api & canonicalize & surrogate-pair & dragonbox-rs & json-canon-rs & json-canon-rs & 1.06$\times$ & [0.990, 1.148] & 0.1908 & 0.2777 & 0.62 & 0.0000 & 8.3 \\
schubfach:go-vs-rs & api & canonicalize & surrogate-pair & schubfach & schubfach-rs & schubfach-rs$^*$ & 2.05$\times$ & [1.873, 2.129] & 0.0002 & 0.0004 & 15.80 & 0.0000 & 11.7 \\
bd:go-vs-rs & api & canonicalize & surrogate-pair & json-canon & json-canon-rs & json-canon-rs$^*$ & 2.05$\times$ & [1.964, 2.126] & 0.0002 & 0.0004 & 22.02 & 0.0000 & 8.3 \\
ryu:go-vs-rs & api & canonicalize & surrogate-pair & ryu & ryu-rs & ryu-rs$^*$ & 2.05$\times$ & [1.973, 2.109] & 0.0002 & 0.0004 & 26.37 & 0.0000 & 7.0 \\
dragonbox:go-vs-rs & api & canonicalize & surrogate-pair & dragonbox & dragonbox-rs & dragonbox-rs$^*$ & 1.94$\times$ & [1.801, 2.048] & 0.0002 & 0.0004 & 11.37 & 0.0000 & 12.9 \\
go:schubfach-vs-bd & api & canonicalize & unicode & schubfach & json-canon & json-canon & 1.00$\times$ & [0.998, 1.006] & 0.4655 & 0.5726 & 0.35 & 0.0000 & 0.7 \\
go:schubfach-vs-ryu & api & canonicalize & unicode & schubfach & ryu & schubfach$^*$ & 1.01$\times$ & [1.003, 1.015] & 0.0096 & 0.0184 & -1.24 & 0.0000 & 0.5 \\
go:schubfach-vs-dragonbox & api & canonicalize & unicode & schubfach & dragonbox & schubfach & 1.00$\times$ & [0.999, 1.008] & 0.1476 & 0.2273 & -0.69 & 0.0000 & 0.5 \\
go:ryu-vs-dragonbox & api & canonicalize & unicode & ryu & dragonbox & dragonbox & 1.00$\times$ & [0.999, 1.012] & 0.1764 & 0.2602 & 0.62 & 0.0000 & 0.8 \\
go:ryu-vs-bd & api & canonicalize & unicode & ryu & json-canon & json-canon$^*$ & 1.01$\times$ & [1.004, 1.017] & 0.0036 & 0.0072 & 1.37 & 0.0000 & 0.7 \\
go:dragonbox-vs-bd & api & canonicalize & unicode & dragonbox & json-canon & json-canon & 1.01$\times$ & [1.000, 1.010] & 0.0710 & 0.1190 & 0.89 & 0.0000 & 0.7 \\
rs:schubfach-vs-bd & api & canonicalize & unicode & schubfach-rs & json-canon-rs & schubfach-rs$^*$ & 1.04$\times$ & [1.021, 1.076] & 0.0040 & 0.0079 & -1.38 & 0.0000 & 3.1 \\
rs:schubfach-vs-ryu & api & canonicalize & unicode & schubfach-rs & ryu-rs & schubfach-rs & 1.02$\times$ & [0.996, 1.045] & 0.2110 & 0.3023 & -0.60 & 0.0000 & 3.1 \\
rs:schubfach-vs-dragonbox & api & canonicalize & unicode & schubfach-rs & dragonbox-rs & schubfach-rs$^*$ & 1.06$\times$ & [1.040, 1.079] & 0.0004 & 0.0008 & -2.69 & 0.0000 & 3.1 \\
rs:ryu-vs-dragonbox & api & canonicalize & unicode & ryu-rs & dragonbox-rs & ryu-rs$^*$ & 1.04$\times$ & [1.017, 1.063] & 0.0022 & 0.0045 & -1.56 & 0.0001 & 4.1 \\
rs:ryu-vs-bd & api & canonicalize & unicode & ryu-rs & json-canon-rs & ryu-rs & 1.03$\times$ & [0.998, 1.058] & 0.1096 & 0.1746 & -0.76 & 0.0001 & 4.1 \\
rs:dragonbox-vs-bd & api & canonicalize & unicode & dragonbox-rs & json-canon-rs & json-canon-rs & 1.01$\times$ & [0.984, 1.035] & 0.2931 & 0.3935 & 0.48 & 0.0001 & 4.7 \\
schubfach:go-vs-rs & api & canonicalize & unicode & schubfach & schubfach-rs & schubfach-rs$^*$ & 1.04$\times$ & [1.020, 1.049] & 0.0006 & 0.0013 & 2.00 & 0.0000 & 3.1 \\
bd:go-vs-rs & api & canonicalize & unicode & json-canon & json-canon-rs & json-canon & 1.01$\times$ & [0.993, 1.039] & 0.4135 & 0.5210 & -0.41 & 0.0000 & 0.7 \\
ryu:go-vs-rs & api & canonicalize & unicode & ryu & ryu-rs & ryu-rs$^*$ & 1.03$\times$ & [1.002, 1.044] & 0.0256 & 0.0469 & 1.09 & 0.0001 & 4.1 \\
dragonbox:go-vs-rs & api & canonicalize & unicode & dragonbox & dragonbox-rs & dragonbox$^*$ & 1.02$\times$ & [1.010, 1.033] & 0.0026 & 0.0053 & -1.42 & 0.0000 & 0.8 \\
go:schubfach-vs-bd & api & canonicalize & verify-whitespace & schubfach & json-canon & schubfach$^*$ & 1.73$\times$ & [1.716, 1.734] & 0.0002 & 0.0004 & -76.68 & 0.0000 & 0.7 \\
go:schubfach-vs-ryu & api & canonicalize & verify-whitespace & schubfach & ryu & ryu & 1.00$\times$ & [0.998, 1.006] & 0.4207 & 0.5280 & 0.38 & 0.0000 & 0.6 \\
go:schubfach-vs-dragonbox & api & canonicalize & verify-whitespace & schubfach & dragonbox & schubfach & 1.00$\times$ & [0.995, 1.005] & 0.9304 & 0.9548 & -0.05 & 0.0000 & 0.7 \\
go:ryu-vs-dragonbox & api & canonicalize & verify-whitespace & ryu & dragonbox & ryu & 1.00$\times$ & [0.997, 1.007] & 0.3991 & 0.5063 & -0.40 & 0.0000 & 0.6 \\
go:ryu-vs-bd & api & canonicalize & verify-whitespace & ryu & json-canon & ryu$^*$ & 1.73$\times$ & [1.720, 1.737] & 0.0002 & 0.0004 & -78.82 & 0.0000 & 0.6 \\
go:dragonbox-vs-bd & api & canonicalize & verify-whitespace & dragonbox & json-canon & dragonbox$^*$ & 1.72$\times$ & [1.715, 1.734] & 0.0002 & 0.0004 & -74.10 & 0.0000 & 0.8 \\
rs:schubfach-vs-bd & api & canonicalize & verify-whitespace & schubfach-rs & json-canon-rs & schubfach-rs$^*$ & 2.96$\times$ & [2.866, 3.039] & 0.0002 & 0.0004 & -30.83 & 0.0000 & 5.2 \\
rs:schubfach-vs-ryu & api & canonicalize & verify-whitespace & schubfach-rs & ryu-rs & schubfach-rs & 1.02$\times$ & [0.985, 1.051] & 0.2625 & 0.3626 & -0.51 & 0.0000 & 5.2 \\
rs:schubfach-vs-dragonbox & api & canonicalize & verify-whitespace & schubfach-rs & dragonbox-rs & schubfach-rs & 1.03$\times$ & [0.995, 1.058] & 0.0980 & 0.1576 & -0.77 & 0.0000 & 5.2 \\
rs:ryu-vs-dragonbox & api & canonicalize & verify-whitespace & ryu-rs & dragonbox-rs & ryu-rs & 1.01$\times$ & [0.982, 1.036] & 0.5683 & 0.6649 & -0.25 & 0.0000 & 4.4 \\
rs:ryu-vs-bd & api & canonicalize & verify-whitespace & ryu-rs & json-canon-rs & ryu-rs$^*$ & 2.90$\times$ & [2.830, 2.980] & 0.0002 & 0.0004 & -31.41 & 0.0000 & 4.4 \\
rs:dragonbox-vs-bd & api & canonicalize & verify-whitespace & dragonbox-rs & json-canon-rs & dragonbox-rs$^*$ & 2.87$\times$ & [2.813, 2.947] & 0.0002 & 0.0004 & -32.05 & 0.0000 & 3.7 \\
schubfach:go-vs-rs & api & canonicalize & verify-whitespace & schubfach & schubfach-rs & schubfach-rs$^*$ & 2.28$\times$ & [2.222, 2.337] & 0.0002 & 0.0004 & 41.56 & 0.0000 & 5.2 \\
bd:go-vs-rs & api & canonicalize & verify-whitespace & json-canon & json-canon-rs & json-canon-rs$^*$ & 1.33$\times$ & [1.307, 1.352] & 0.0002 & 0.0004 & 16.58 & 0.0000 & 3.4 \\
ryu:go-vs-rs & api & canonicalize & verify-whitespace & ryu & ryu-rs & ryu-rs$^*$ & 2.24$\times$ & [2.190, 2.283] & 0.0002 & 0.0004 & 47.42 & 0.0000 & 4.4 \\
dragonbox:go-vs-rs & api & canonicalize & verify-whitespace & dragonbox & dragonbox-rs & dragonbox-rs$^*$ & 2.22$\times$ & [2.185, 2.263] & 0.0002 & 0.0004 & 52.24 & 0.0000 & 3.7 \\
go:schubfach-vs-bd & api & verify & canonical/array-2048 & schubfach & json-canon & schubfach$^*$ & 1.66$\times$ & [1.652, 1.670] & 0.0002 & 0.0004 & -62.27 & 0.0068 & 0.3 \\
go:schubfach-vs-ryu & api & verify & canonical/array-2048 & schubfach & ryu & schubfach & 1.01$\times$ & [1.000, 1.009] & 0.0360 & 0.0642 & -1.03 & 0.0068 & 0.3 \\
go:schubfach-vs-dragonbox & api & verify & canonical/array-2048 & schubfach & dragonbox & schubfach$^*$ & 1.01$\times$ & [1.011, 1.017] & 0.0002 & 0.0004 & -3.32 & 0.0068 & 0.3 \\
go:ryu-vs-dragonbox & api & verify & canonical/array-2048 & ryu & dragonbox & ryu$^*$ & 1.01$\times$ & [1.004, 1.014] & 0.0036 & 0.0072 & -1.43 & 0.0160 & 0.8 \\
go:ryu-vs-bd & api & verify & canonical/array-2048 & ryu & json-canon & ryu$^*$ & 1.65$\times$ & [1.643, 1.664] & 0.0002 & 0.0004 & -57.59 & 0.0160 & 0.8 \\
go:dragonbox-vs-bd & api & verify & canonical/array-2048 & dragonbox & json-canon & dragonbox$^*$ & 1.64$\times$ & [1.630, 1.649] & 0.0002 & 0.0004 & -58.83 & 0.0123 & 0.6 \\
rs:schubfach-vs-bd & api & verify & canonical/array-2048 & schubfach-rs & json-canon-rs & schubfach-rs$^*$ & 2.50$\times$ & [2.472, 2.517] & 0.0002 & 0.0004 & -85.54 & 0.0167 & 1.5 \\
rs:schubfach-vs-ryu & api & verify & canonical/array-2048 & schubfach-rs & ryu-rs & schubfach-rs & 1.02$\times$ & [1.001, 1.060] & 0.1818 & 0.2664 & -0.57 & 0.0167 & 1.5 \\
rs:schubfach-vs-dragonbox & api & verify & canonical/array-2048 & schubfach-rs & dragonbox-rs & schubfach-rs & 1.00$\times$ & [0.991, 1.009] & 0.9152 & 0.9429 & -0.05 & 0.0167 & 1.5 \\
rs:ryu-vs-dragonbox & api & verify & canonical/array-2048 & ryu-rs & dragonbox-rs & dragonbox-rs & 1.02$\times$ & [1.001, 1.063] & 0.1582 & 0.2402 & 0.57 & 0.0114 & 1.1 \\
rs:ryu-vs-bd & api & verify & canonical/array-2048 & ryu-rs & json-canon-rs & ryu-rs$^*$ & 2.45$\times$ & [2.349, 2.491] & 0.0002 & 0.0004 & -44.51 & 0.0562 & 5.1 \\
rs:dragonbox-vs-bd & api & verify & canonical/array-2048 & dragonbox-rs & json-canon-rs & dragonbox-rs$^*$ & 2.49$\times$ & [2.478, 2.512] & 0.0002 & 0.0004 & -91.89 & 0.0114 & 1.1 \\
schubfach:go-vs-rs & api & verify & canonical/array-2048 & schubfach & schubfach-rs & schubfach-rs$^*$ & 1.84$\times$ & [1.824, 1.851] & 0.0002 & 0.0004 & 88.74 & 0.0167 & 1.5 \\
bd:go-vs-rs & api & verify & canonical/array-2048 & json-canon & json-canon-rs & json-canon-rs$^*$ & 1.23$\times$ & [1.215, 1.233] & 0.0002 & 0.0004 & 22.97 & 0.0288 & 1.1 \\
ryu:go-vs-rs & api & verify & canonical/array-2048 & ryu & ryu-rs & ryu-rs$^*$ & 1.82$\times$ & [1.728, 1.845] & 0.0002 & 0.0004 & 27.06 & 0.0562 & 5.1 \\
dragonbox:go-vs-rs & api & verify & canonical/array-2048 & dragonbox & dragonbox-rs & dragonbox-rs$^*$ & 1.86$\times$ & [1.853, 1.874] & 0.0002 & 0.0004 & 98.12 & 0.0114 & 1.1 \\
go:schubfach-vs-bd & api & verify & canonical/array-256 & schubfach & json-canon & schubfach$^*$ & 1.54$\times$ & [1.534, 1.554] & 0.0002 & 0.0004 & -57.10 & 0.0026 & 1.0 \\
go:schubfach-vs-ryu & api & verify & canonical/array-256 & schubfach & ryu & ryu & 1.01$\times$ & [1.001, 1.013] & 0.0432 & 0.0750 & 0.95 & 0.0017 & 0.7 \\
go:schubfach-vs-dragonbox & api & verify & canonical/array-256 & schubfach & dragonbox & dragonbox & 1.01$\times$ & [1.001, 1.014] & 0.0674 & 0.1138 & 0.86 & 0.0023 & 0.9 \\
go:ryu-vs-dragonbox & api & verify & canonical/array-256 & ryu & dragonbox & dragonbox & 1.00$\times$ & [0.995, 1.006] & 0.9630 & 0.9751 & 0.02 & 0.0023 & 0.9 \\
go:ryu-vs-bd & api & verify & canonical/array-256 & ryu & json-canon & ryu$^*$ & 1.56$\times$ & [1.548, 1.564] & 0.0002 & 0.0004 & -65.44 & 0.0017 & 0.7 \\
go:dragonbox-vs-bd & api & verify & canonical/array-256 & dragonbox & json-canon & dragonbox$^*$ & 1.56$\times$ & [1.547, 1.565] & 0.0002 & 0.0004 & -60.53 & 0.0023 & 0.9 \\
rs:schubfach-vs-bd & api & verify & canonical/array-256 & schubfach-rs & json-canon-rs & schubfach-rs$^*$ & 2.26$\times$ & [2.237, 2.274] & 0.0002 & 0.0004 & -71.98 & 0.0016 & 1.2 \\
rs:schubfach-vs-ryu & api & verify & canonical/array-256 & schubfach-rs & ryu-rs & schubfach-rs & 1.01$\times$ & [0.987, 1.025] & 0.4521 & 0.5609 & -0.34 & 0.0016 & 1.2 \\
rs:schubfach-vs-dragonbox & api & verify & canonical/array-256 & schubfach-rs & dragonbox-rs & dragonbox-rs & 1.01$\times$ & [0.998, 1.012] & 0.1986 & 0.2871 & 0.59 & 0.0012 & 0.9 \\
rs:ryu-vs-dragonbox & api & verify & canonical/array-256 & ryu-rs & dragonbox-rs & dragonbox-rs & 1.01$\times$ & [0.993, 1.030] & 0.2304 & 0.3268 & 0.57 & 0.0012 & 0.9 \\
rs:ryu-vs-bd & api & verify & canonical/array-256 & ryu-rs & json-canon-rs & ryu-rs$^*$ & 2.24$\times$ & [2.199, 2.283] & 0.0002 & 0.0004 & -46.44 & 0.0051 & 3.8 \\
rs:dragonbox-vs-bd & api & verify & canonical/array-256 & dragonbox-rs & json-canon-rs & dragonbox-rs$^*$ & 2.27$\times$ & [2.250, 2.283] & 0.0002 & 0.0004 & -74.98 & 0.0012 & 0.9 \\
schubfach:go-vs-rs & api & verify & canonical/array-256 & schubfach & schubfach-rs & schubfach-rs$^*$ & 1.88$\times$ & [1.869, 1.898] & 0.0002 & 0.0004 & 67.88 & 0.0016 & 1.2 \\
bd:go-vs-rs & api & verify & canonical/array-256 & json-canon & json-canon-rs & json-canon-rs$^*$ & 1.29$\times$ & [1.280, 1.299] & 0.0002 & 0.0004 & 30.58 & 0.0038 & 1.3 \\
ryu:go-vs-rs & api & verify & canonical/array-256 & ryu & ryu-rs & ryu-rs$^*$ & 1.86$\times$ & [1.826, 1.894] & 0.0002 & 0.0004 & 38.00 & 0.0051 & 3.8 \\
dragonbox:go-vs-rs & api & verify & canonical/array-256 & dragonbox & dragonbox-rs & dragonbox-rs$^*$ & 1.88$\times$ & [1.868, 1.891] & 0.0002 & 0.0004 & 80.13 & 0.0012 & 0.9 \\
go:schubfach-vs-bd & api & verify & canonical/canonical-minimal & schubfach & json-canon & schubfach$^*$ & 1.51$\times$ & [1.504, 1.518] & 0.0002 & 0.0004 & -70.33 & 0.0000 & 0.5 \\
go:schubfach-vs-ryu & api & verify & canonical/canonical-minimal & schubfach & ryu & ryu & 1.00$\times$ & [0.996, 1.008] & 0.6033 & 0.6976 & 0.26 & 0.0000 & 1.2 \\
go:schubfach-vs-dragonbox & api & verify & canonical/canonical-minimal & schubfach & dragonbox & dragonbox & 1.00$\times$ & [0.999, 1.009] & 0.1696 & 0.2531 & 0.68 & 0.0000 & 0.9 \\
go:ryu-vs-dragonbox & api & verify & canonical/canonical-minimal & ryu & dragonbox & dragonbox & 1.00$\times$ & [0.995, 1.010] & 0.5559 & 0.6533 & 0.26 & 0.0000 & 0.9 \\
go:ryu-vs-bd & api & verify & canonical/canonical-minimal & ryu & json-canon & ryu$^*$ & 1.51$\times$ & [1.503, 1.524] & 0.0002 & 0.0004 & -54.09 & 0.0000 & 1.2 \\
go:dragonbox-vs-bd & api & verify & canonical/canonical-minimal & dragonbox & json-canon & dragonbox$^*$ & 1.52$\times$ & [1.509, 1.526] & 0.0002 & 0.0004 & -61.09 & 0.0000 & 0.9 \\
rs:schubfach-vs-bd & api & verify & canonical/canonical-minimal & schubfach-rs & json-canon-rs & schubfach-rs$^*$ & 2.08$\times$ & [2.021, 2.147] & 0.0002 & 0.0004 & -18.07 & 0.0000 & 4.6 \\
rs:schubfach-vs-ryu & api & verify & canonical/canonical-minimal & schubfach-rs & ryu-rs & schubfach-rs & 1.03$\times$ & [0.988, 1.080] & 0.2541 & 0.3532 & -0.53 & 0.0000 & 4.6 \\
rs:schubfach-vs-dragonbox & api & verify & canonical/canonical-minimal & schubfach-rs & dragonbox-rs & dragonbox-rs & 1.01$\times$ & [0.979, 1.045] & 0.4795 & 0.5821 & 0.33 & 0.0000 & 5.2 \\
rs:ryu-vs-dragonbox & api & verify & canonical/canonical-minimal & ryu-rs & dragonbox-rs & dragonbox-rs & 1.04$\times$ & [1.000, 1.093] & 0.1158 & 0.1838 & 0.74 & 0.0000 & 5.2 \\
rs:ryu-vs-bd & api & verify & canonical/canonical-minimal & ryu-rs & json-canon-rs & ryu-rs$^*$ & 2.03$\times$ & [1.932, 2.107] & 0.0002 & 0.0004 & -14.57 & 0.0000 & 8.3 \\
rs:dragonbox-vs-bd & api & verify & canonical/canonical-minimal & dragonbox-rs & json-canon-rs & dragonbox-rs$^*$ & 2.11$\times$ & [2.043, 2.176] & 0.0002 & 0.0004 & -17.84 & 0.0000 & 5.2 \\
schubfach:go-vs-rs & api & verify & canonical/canonical-minimal & schubfach & schubfach-rs & schubfach-rs$^*$ & 1.73$\times$ & [1.678, 1.756] & 0.0002 & 0.0004 & 27.62 & 0.0000 & 4.6 \\
bd:go-vs-rs & api & verify & canonical/canonical-minimal & json-canon & json-canon-rs & json-canon-rs$^*$ & 1.25$\times$ & [1.221, 1.276] & 0.0002 & 0.0004 & 9.52 & 0.0000 & 4.6 \\
ryu:go-vs-rs & api & verify & canonical/canonical-minimal & ryu & ryu-rs & ryu-rs$^*$ & 1.68$\times$ & [1.596, 1.732] & 0.0002 & 0.0004 & 14.04 & 0.0000 & 8.3 \\
dragonbox:go-vs-rs & api & verify & canonical/canonical-minimal & dragonbox & dragonbox-rs & dragonbox-rs$^*$ & 1.74$\times$ & [1.690, 1.778] & 0.0002 & 0.0004 & 24.01 & 0.0000 & 5.2 \\
go:schubfach-vs-bd & api & verify & canonical/control-escapes & schubfach & json-canon & schubfach$^*$ & 1.02$\times$ & [1.017, 1.029] & 0.0002 & 0.0004 & -3.42 & 0.0000 & 1.1 \\
go:schubfach-vs-ryu & api & verify & canonical/control-escapes & schubfach & ryu & ryu & 1.00$\times$ & [0.999, 1.011] & 0.1552 & 0.2368 & 0.68 & 0.0000 & 0.5 \\
go:schubfach-vs-dragonbox & api & verify & canonical/control-escapes & schubfach & dragonbox & schubfach & 1.00$\times$ & [0.994, 1.007] & 0.9200 & 0.9471 & -0.05 & 0.0000 & 1.1 \\
go:ryu-vs-dragonbox & api & verify & canonical/control-escapes & ryu & dragonbox & ryu & 1.01$\times$ & [1.001, 1.010] & 0.0404 & 0.0709 & -0.96 & 0.0000 & 0.5 \\
go:ryu-vs-bd & api & verify & canonical/control-escapes & ryu & json-canon & ryu$^*$ & 1.03$\times$ & [1.025, 1.032] & 0.0002 & 0.0004 & -7.08 & 0.0000 & 0.5 \\
go:dragonbox-vs-bd & api & verify & canonical/control-escapes & dragonbox & json-canon & dragonbox$^*$ & 1.02$\times$ & [1.018, 1.027] & 0.0002 & 0.0004 & -4.38 & 0.0000 & 0.8 \\
rs:schubfach-vs-bd & api & verify & canonical/control-escapes & schubfach-rs & json-canon-rs & json-canon-rs & 1.04$\times$ & [0.998, 1.083] & 0.0784 & 0.1299 & 0.86 & 0.0000 & 6.0 \\
rs:schubfach-vs-ryu & api & verify & canonical/control-escapes & schubfach-rs & ryu-rs & schubfach-rs & 1.02$\times$ & [0.972, 1.076] & 0.5995 & 0.6945 & -0.26 & 0.0000 & 6.4 \\
rs:schubfach-vs-dragonbox & api & verify & canonical/control-escapes & schubfach-rs & dragonbox-rs & dragonbox-rs & 1.02$\times$ & [0.988, 1.053] & 0.2851 & 0.3847 & 0.51 & 0.0000 & 2.6 \\
rs:ryu-vs-dragonbox & api & verify & canonical/control-escapes & ryu-rs & dragonbox-rs & dragonbox-rs & 1.04$\times$ & [1.003, 1.093] & 0.1210 & 0.1907 & 0.70 & 0.0000 & 2.6 \\
rs:ryu-vs-bd & api & verify & canonical/control-escapes & ryu-rs & json-canon-rs & json-canon-rs & 1.06$\times$ & [1.014, 1.117] & 0.0334 & 0.0598 & 0.99 & 0.0000 & 6.0 \\
rs:dragonbox-vs-bd & api & verify & canonical/control-escapes & dragonbox-rs & json-canon-rs & json-canon-rs & 1.02$\times$ & [0.981, 1.049] & 0.1886 & 0.2748 & 0.62 & 0.0000 & 6.0 \\
schubfach:go-vs-rs & api & verify & canonical/control-escapes & schubfach & schubfach-rs & schubfach-rs$^*$ & 2.04$\times$ & [1.976, 2.102] & 0.0002 & 0.0004 & 27.11 & 0.0000 & 6.4 \\
bd:go-vs-rs & api & verify & canonical/control-escapes & json-canon & json-canon-rs & json-canon-rs$^*$ & 2.18$\times$ & [2.096, 2.228] & 0.0002 & 0.0004 & 33.96 & 0.0000 & 6.0 \\
ryu:go-vs-rs & api & verify & canonical/control-escapes & ryu & ryu-rs & ryu-rs$^*$ & 2.00$\times$ & [1.882, 2.054] & 0.0002 & 0.0004 & 21.06 & 0.0000 & 8.3 \\
dragonbox:go-vs-rs & api & verify & canonical/control-escapes & dragonbox & dragonbox-rs & dragonbox-rs$^*$ & 2.08$\times$ & [2.051, 2.105] & 0.0002 & 0.0004 & 62.33 & 0.0000 & 2.6 \\
go:schubfach-vs-bd & api & verify & canonical/deep-64 & schubfach & json-canon & json-canon$^*$ & 1.01$\times$ & [1.009, 1.017] & 0.0002 & 0.0004 & 2.80 & 0.0001 & 0.6 \\
go:schubfach-vs-ryu & api & verify & canonical/deep-64 & schubfach & ryu & schubfach & 1.00$\times$ & [0.996, 1.006] & 0.8342 & 0.8790 & -0.09 & 0.0001 & 0.6 \\
go:schubfach-vs-dragonbox & api & verify & canonical/deep-64 & schubfach & dragonbox & schubfach & 1.00$\times$ & [1.001, 1.008] & 0.0318 & 0.0573 & -1.02 & 0.0001 & 0.6 \\
go:ryu-vs-dragonbox & api & verify & canonical/deep-64 & ryu & dragonbox & ryu & 1.00$\times$ & [0.998, 1.008] & 0.1760 & 0.2602 & -0.64 & 0.0001 & 0.9 \\
go:ryu-vs-bd & api & verify & canonical/deep-64 & ryu & json-canon & json-canon$^*$ & 1.01$\times$ & [1.009, 1.019] & 0.0004 & 0.0008 & 2.17 & 0.0001 & 0.6 \\
go:dragonbox-vs-bd & api & verify & canonical/deep-64 & dragonbox & json-canon & json-canon$^*$ & 1.02$\times$ & [1.014, 1.021] & 0.0002 & 0.0004 & 3.95 & 0.0001 & 0.6 \\
rs:schubfach-vs-bd & api & verify & canonical/deep-64 & schubfach-rs & json-canon-rs & json-canon-rs$^*$ & 1.02$\times$ & [1.019, 1.031] & 0.0002 & 0.0004 & 3.43 & 0.0001 & 0.9 \\
rs:schubfach-vs-ryu & api & verify & canonical/deep-64 & schubfach-rs & ryu-rs & ryu-rs & 1.00$\times$ & [0.996, 1.010] & 0.4519 & 0.5609 & 0.35 & 0.0001 & 1.2 \\
rs:schubfach-vs-dragonbox & api & verify & canonical/deep-64 & schubfach-rs & dragonbox-rs & schubfach-rs$^*$ & 1.03$\times$ & [1.010, 1.067] & 0.0022 & 0.0045 & -0.88 & 0.0001 & 0.9 \\
rs:ryu-vs-dragonbox & api & verify & canonical/deep-64 & ryu-rs & dragonbox-rs & ryu-rs$^*$ & 1.03$\times$ & [1.013, 1.071] & 0.0010 & 0.0021 & -0.97 & 0.0001 & 1.2 \\
rs:ryu-vs-bd & api & verify & canonical/deep-64 & ryu-rs & json-canon-rs & json-canon-rs$^*$ & 1.02$\times$ & [1.015, 1.029] & 0.0002 & 0.0004 & 2.59 & 0.0001 & 0.9 \\
rs:dragonbox-vs-bd & api & verify & canonical/deep-64 & dragonbox-rs & json-canon-rs & json-canon-rs$^*$ & 1.05$\times$ & [1.035, 1.094] & 0.0002 & 0.0004 & 1.71 & 0.0001 & 0.9 \\
schubfach:go-vs-rs & api & verify & canonical/deep-64 & schubfach & schubfach-rs & schubfach-rs$^*$ & 1.52$\times$ & [1.510, 1.525] & 0.0004 & 0.0008 & 74.84 & 0.0001 & 0.9 \\
bd:go-vs-rs & api & verify & canonical/deep-64 & json-canon & json-canon-rs & json-canon-rs$^*$ & 1.54$\times$ & [1.528, 1.544] & 0.0002 & 0.0004 & 74.60 & 0.0001 & 0.9 \\
ryu:go-vs-rs & api & verify & canonical/deep-64 & ryu & ryu-rs & ryu-rs$^*$ & 1.52$\times$ & [1.513, 1.535] & 0.0002 & 0.0004 & 51.05 & 0.0001 & 1.2 \\
dragonbox:go-vs-rs & api & verify & canonical/deep-64 & dragonbox & dragonbox-rs & dragonbox-rs$^*$ & 1.49$\times$ & [1.428, 1.509] & 0.0002 & 0.0004 & 17.24 & 0.0004 & 4.9 \\
go:schubfach-vs-bd & api & verify & canonical/deep & schubfach & json-canon & json-canon$^*$ & 1.01$\times$ & [1.005, 1.017] & 0.0014 & 0.0029 & 1.49 & 0.0002 & 0.8 \\
go:schubfach-vs-ryu & api & verify & canonical/deep & schubfach & ryu & ryu & 1.00$\times$ & [0.995, 1.008] & 0.7580 & 0.8229 & 0.14 & 0.0004 & 1.2 \\
go:schubfach-vs-dragonbox & api & verify & canonical/deep & schubfach & dragonbox & schubfach & 1.01$\times$ & [1.000, 1.015] & 0.0950 & 0.1536 & -0.80 & 0.0003 & 0.9 \\
go:ryu-vs-dragonbox & api & verify & canonical/deep & ryu & dragonbox & ryu & 1.01$\times$ & [1.000, 1.017] & 0.0760 & 0.1267 & -0.84 & 0.0004 & 1.2 \\
go:ryu-vs-bd & api & verify & canonical/deep & ryu & json-canon & json-canon$^*$ & 1.01$\times$ & [1.002, 1.015] & 0.0200 & 0.0370 & 1.12 & 0.0002 & 0.8 \\
go:dragonbox-vs-bd & api & verify & canonical/deep & dragonbox & json-canon & json-canon$^*$ & 1.02$\times$ & [1.010, 1.025] & 0.0008 & 0.0017 & 1.97 & 0.0002 & 0.8 \\
rs:schubfach-vs-bd & api & verify & canonical/deep & schubfach-rs & json-canon-rs & schubfach-rs$^*$ & 1.02$\times$ & [1.014, 1.033] & 0.0004 & 0.0008 & -2.09 & 0.0004 & 1.8 \\
rs:schubfach-vs-ryu & api & verify & canonical/deep & schubfach-rs & ryu-rs & schubfach-rs$^*$ & 1.07$\times$ & [1.061, 1.083] & 0.0004 & 0.0008 & -5.58 & 0.0004 & 1.8 \\
rs:schubfach-vs-dragonbox & api & verify & canonical/deep & schubfach-rs & dragonbox-rs & schubfach-rs$^*$ & 1.06$\times$ & [1.047, 1.068] & 0.0002 & 0.0004 & -4.80 & 0.0004 & 1.8 \\
rs:ryu-vs-dragonbox & api & verify & canonical/deep & ryu-rs & dragonbox-rs & dragonbox-rs$^*$ & 1.01$\times$ & [1.006, 1.022] & 0.0038 & 0.0075 & 1.41 & 0.0003 & 1.1 \\
rs:ryu-vs-bd & api & verify & canonical/deep & ryu-rs & json-canon-rs & json-canon-rs$^*$ & 1.05$\times$ & [1.041, 1.056] & 0.0002 & 0.0004 & 5.31 & 0.0002 & 0.8 \\
rs:dragonbox-vs-bd & api & verify & canonical/deep & dragonbox-rs & json-canon-rs & json-canon-rs$^*$ & 1.03$\times$ & [1.027, 1.040] & 0.0004 & 0.0008 & 4.36 & 0.0002 & 0.8 \\
schubfach:go-vs-rs & api & verify & canonical/deep & schubfach & schubfach-rs & schubfach-rs$^*$ & 1.33$\times$ & [1.321, 1.346] & 0.0002 & 0.0004 & 26.94 & 0.0004 & 1.8 \\
bd:go-vs-rs & api & verify & canonical/deep & json-canon & json-canon-rs & json-canon-rs$^*$ & 1.29$\times$ & [1.284, 1.297] & 0.0002 & 0.0004 & 40.45 & 0.0002 & 0.8 \\
ryu:go-vs-rs & api & verify & canonical/deep & ryu & ryu-rs & ryu-rs$^*$ & 1.24$\times$ & [1.232, 1.252] & 0.0002 & 0.0004 & 22.13 & 0.0003 & 1.3 \\
dragonbox:go-vs-rs & api & verify & canonical/deep & dragonbox & dragonbox-rs & dragonbox-rs$^*$ & 1.27$\times$ & [1.260, 1.280] & 0.0002 & 0.0004 & 23.71 & 0.0003 & 1.1 \\
go:schubfach-vs-bd & api & verify & canonical/escaped-key-order & schubfach & json-canon & schubfach$^*$ & 1.61$\times$ & [1.605, 1.622] & 0.0002 & 0.0004 & -90.06 & 0.0000 & 1.0 \\
go:schubfach-vs-ryu & api & verify & canonical/escaped-key-order & schubfach & ryu & ryu & 1.00$\times$ & [0.998, 1.010] & 0.1648 & 0.2482 & 0.63 & 0.0000 & 0.8 \\
go:schubfach-vs-dragonbox & api & verify & canonical/escaped-key-order & schubfach & dragonbox & dragonbox & 1.00$\times$ & [0.996, 1.008] & 0.6399 & 0.7282 & 0.21 & 0.0000 & 0.9 \\
go:ryu-vs-dragonbox & api & verify & canonical/escaped-key-order & ryu & dragonbox & ryu & 1.00$\times$ & [0.997, 1.008] & 0.3105 & 0.4126 & -0.45 & 0.0000 & 0.8 \\
go:ryu-vs-bd & api & verify & canonical/escaped-key-order & ryu & json-canon & ryu$^*$ & 1.62$\times$ & [1.613, 1.627] & 0.0002 & 0.0004 & -108.72 & 0.0000 & 0.8 \\
go:dragonbox-vs-bd & api & verify & canonical/escaped-key-order & dragonbox & json-canon & dragonbox$^*$ & 1.62$\times$ & [1.609, 1.624] & 0.0002 & 0.0004 & -100.56 & 0.0000 & 0.9 \\
rs:schubfach-vs-bd & api & verify & canonical/escaped-key-order & schubfach-rs & json-canon-rs & schubfach-rs$^*$ & 2.53$\times$ & [2.335, 2.622] & 0.0002 & 0.0004 & -19.20 & 0.0000 & 10.3 \\
rs:schubfach-vs-ryu & api & verify & canonical/escaped-key-order & schubfach-rs & ryu-rs & schubfach-rs & 1.02$\times$ & [0.943, 1.058] & 0.6165 & 0.7083 & -0.27 & 0.0000 & 10.3 \\
rs:schubfach-vs-dragonbox & api & verify & canonical/escaped-key-order & schubfach-rs & dragonbox-rs & schubfach-rs & 1.02$\times$ & [0.937, 1.055] & 0.6181 & 0.7084 & -0.26 & 0.0000 & 10.3 \\
rs:ryu-vs-dragonbox & api & verify & canonical/escaped-key-order & ryu-rs & dragonbox-rs & dragonbox-rs & 1.00$\times$ & [0.968, 1.032] & 0.9696 & 0.9787 & 0.02 & 0.0000 & 4.6 \\
rs:ryu-vs-bd & api & verify & canonical/escaped-key-order & ryu-rs & json-canon-rs & ryu-rs$^*$ & 2.48$\times$ & [2.407, 2.552] & 0.0002 & 0.0004 & -24.73 & 0.0000 & 5.0 \\
rs:dragonbox-vs-bd & api & verify & canonical/escaped-key-order & dragonbox-rs & json-canon-rs & dragonbox-rs$^*$ & 2.48$\times$ & [2.414, 2.553] & 0.0002 & 0.0004 & -25.13 & 0.0000 & 4.6 \\
schubfach:go-vs-rs & api & verify & canonical/escaped-key-order & schubfach & schubfach-rs & schubfach-rs$^*$ & 2.11$\times$ & [1.947, 2.180] & 0.0002 & 0.0004 & 18.63 & 0.0000 & 10.3 \\
bd:go-vs-rs & api & verify & canonical/escaped-key-order & json-canon & json-canon-rs & json-canon-rs$^*$ & 1.35$\times$ & [1.320, 1.369] & 0.0002 & 0.0004 & 16.21 & 0.0000 & 3.8 \\
ryu:go-vs-rs & api & verify & canonical/escaped-key-order & ryu & ryu-rs & ryu-rs$^*$ & 2.07$\times$ & [2.009, 2.109] & 0.0002 & 0.0004 & 36.01 & 0.0000 & 5.0 \\
dragonbox:go-vs-rs & api & verify & canonical/escaped-key-order & dragonbox & dragonbox-rs & dragonbox-rs$^*$ & 2.07$\times$ & [2.022, 2.115] & 0.0002 & 0.0004 & 38.22 & 0.0000 & 4.6 \\
go:schubfach-vs-bd & api & verify & canonical/large & schubfach & json-canon & schubfach$^*$ & 1.79$\times$ & [1.779, 1.799] & 0.0002 & 0.0004 & -79.41 & 0.0820 & 0.9 \\
go:schubfach-vs-ryu & api & verify & canonical/large & schubfach & ryu & ryu & 1.00$\times$ & [0.998, 1.008] & 0.2797 & 0.3790 & 0.49 & 0.0515 & 0.6 \\
go:schubfach-vs-dragonbox & api & verify & canonical/large & schubfach & dragonbox & schubfach & 1.01$\times$ & [1.001, 1.014] & 0.0430 & 0.0748 & -0.97 & 0.0820 & 0.9 \\
go:ryu-vs-dragonbox & api & verify & canonical/large & ryu & dragonbox & ryu$^*$ & 1.01$\times$ & [1.004, 1.016] & 0.0038 & 0.0075 & -1.53 & 0.0515 & 0.6 \\
go:ryu-vs-bd & api & verify & canonical/large & ryu & json-canon & ryu$^*$ & 1.79$\times$ & [1.785, 1.802] & 0.0002 & 0.0004 & -86.83 & 0.0515 & 0.6 \\
go:dragonbox-vs-bd & api & verify & canonical/large & dragonbox & json-canon & dragonbox$^*$ & 1.78$\times$ & [1.765, 1.787] & 0.0002 & 0.0004 & -73.80 & 0.1010 & 1.1 \\
rs:schubfach-vs-bd & api & verify & canonical/large & schubfach-rs & json-canon-rs & schubfach-rs$^*$ & 3.30$\times$ & [3.219, 3.357] & 0.0002 & 0.0004 & -66.39 & 0.1928 & 4.1 \\
rs:schubfach-vs-ryu & api & verify & canonical/large & schubfach-rs & ryu-rs & schubfach-rs & 1.00$\times$ & [0.977, 1.042] & 0.8978 & 0.9279 & -0.06 & 0.1928 & 4.1 \\
rs:schubfach-vs-dragonbox & api & verify & canonical/large & schubfach-rs & dragonbox-rs & dragonbox-rs & 1.01$\times$ & [0.977, 1.032] & 0.6949 & 0.7670 & 0.23 & 0.1912 & 4.1 \\
rs:ryu-vs-dragonbox & api & verify & canonical/large & ryu-rs & dragonbox-rs & dragonbox-rs & 1.01$\times$ & [0.983, 1.047] & 0.5075 & 0.6086 & 0.26 & 0.1912 & 4.1 \\
rs:ryu-vs-bd & api & verify & canonical/large & ryu-rs & json-canon-rs & ryu-rs$^*$ & 3.30$\times$ & [3.156, 3.349] & 0.0002 & 0.0004 & -58.83 & 0.2448 & 5.2 \\
rs:dragonbox-vs-bd & api & verify & canonical/large & dragonbox-rs & json-canon-rs & dragonbox-rs$^*$ & 3.33$\times$ & [3.223, 3.372] & 0.0002 & 0.0004 & -66.83 & 0.1912 & 4.1 \\
schubfach:go-vs-rs & api & verify & canonical/large & schubfach & schubfach-rs & schubfach-rs$^*$ & 1.93$\times$ & [1.885, 1.964] & 0.0002 & 0.0004 & 37.28 & 0.1928 & 4.1 \\
bd:go-vs-rs & api & verify & canonical/large & json-canon & json-canon-rs & json-canon-rs$^*$ & 1.05$\times$ & [1.038, 1.055] & 0.0002 & 0.0004 & 5.09 & 0.2168 & 1.4 \\
ryu:go-vs-rs & api & verify & canonical/large & ryu & ryu-rs & ryu-rs$^*$ & 1.92$\times$ & [1.844, 1.954] & 0.0002 & 0.0004 & 30.94 & 0.2448 & 5.2 \\
dragonbox:go-vs-rs & api & verify & canonical/large & dragonbox & dragonbox-rs & dragonbox-rs$^*$ & 1.96$\times$ & [1.892, 1.988] & 0.0002 & 0.0004 & 36.97 & 0.1912 & 4.1 \\
go:schubfach-vs-bd & api & verify & canonical/long-string & schubfach & json-canon & schubfach & 1.00$\times$ & [0.991, 1.021] & 0.5505 & 0.6476 & -0.27 & 0.0011 & 1.8 \\
go:schubfach-vs-ryu & api & verify & canonical/long-string & schubfach & ryu & schubfach & 1.01$\times$ & [0.998, 1.020] & 0.1684 & 0.2519 & -0.65 & 0.0011 & 1.8 \\
go:schubfach-vs-dragonbox & api & verify & canonical/long-string & schubfach & dragonbox & schubfach & 1.00$\times$ & [0.990, 1.013] & 0.8222 & 0.8707 & -0.10 & 0.0011 & 1.8 \\
go:ryu-vs-dragonbox & api & verify & canonical/long-string & ryu & dragonbox & dragonbox & 1.01$\times$ & [0.998, 1.018] & 0.2026 & 0.2919 & 0.59 & 0.0010 & 1.6 \\
go:ryu-vs-bd & api & verify & canonical/long-string & ryu & json-canon & json-canon & 1.00$\times$ & [0.989, 1.017] & 0.6283 & 0.7163 & 0.22 & 0.0016 & 2.6 \\
go:dragonbox-vs-bd & api & verify & canonical/long-string & dragonbox & json-canon & dragonbox & 1.00$\times$ & [0.990, 1.018] & 0.6531 & 0.7394 & -0.20 & 0.0010 & 1.6 \\
rs:schubfach-vs-bd & api & verify & canonical/long-string & schubfach-rs & json-canon-rs & json-canon-rs & 1.01$\times$ & [0.992, 1.028] & 0.5499 & 0.6476 & 0.25 & 0.0002 & 1.3 \\
rs:schubfach-vs-ryu & api & verify & canonical/long-string & schubfach-rs & ryu-rs & schubfach-rs & 1.01$\times$ & [0.987, 1.021] & 0.3533 & 0.4577 & -0.45 & 0.0005 & 3.3 \\
rs:schubfach-vs-dragonbox & api & verify & canonical/long-string & schubfach-rs & dragonbox-rs & schubfach-rs & 1.00$\times$ & [0.980, 1.018] & 0.6039 & 0.6977 & -0.22 & 0.0005 & 3.3 \\
rs:ryu-vs-dragonbox & api & verify & canonical/long-string & ryu-rs & dragonbox-rs & dragonbox-rs & 1.00$\times$ & [0.999, 1.015] & 0.3619 & 0.4665 & 0.48 & 0.0002 & 1.4 \\
rs:ryu-vs-bd & api & verify & canonical/long-string & ryu-rs & json-canon-rs & json-canon-rs$^*$ & 1.01$\times$ & [1.009, 1.022] & 0.0002 & 0.0004 & 1.76 & 0.0002 & 1.3 \\
rs:dragonbox-vs-bd & api & verify & canonical/long-string & dragonbox-rs & json-canon-rs & json-canon-rs & 1.01$\times$ & [1.000, 1.017] & 0.0646 & 0.1097 & 0.90 & 0.0002 & 1.3 \\
schubfach:go-vs-rs & api & verify & canonical/long-string & schubfach & schubfach-rs & schubfach-rs$^*$ & 4.45$\times$ & [4.350, 4.515] & 0.0002 & 0.0004 & 69.72 & 0.0005 & 3.3 \\
bd:go-vs-rs & api & verify & canonical/long-string & json-canon & json-canon-rs & json-canon-rs$^*$ & 4.50$\times$ & [4.440, 4.560] & 0.0002 & 0.0004 & 52.91 & 0.0002 & 1.3 \\
ryu:go-vs-rs & api & verify & canonical/long-string & ryu & ryu-rs & ryu-rs$^*$ & 4.45$\times$ & [4.422, 4.483] & 0.0002 & 0.0004 & 96.37 & 0.0001 & 0.4 \\
dragonbox:go-vs-rs & api & verify & canonical/long-string & dragonbox & dragonbox-rs & dragonbox-rs$^*$ & 4.44$\times$ & [4.400, 4.489] & 0.0002 & 0.0004 & 86.34 & 0.0002 & 1.4 \\
go:schubfach-vs-bd & api & verify & canonical/medium & schubfach & json-canon & schubfach$^*$ & 1.74$\times$ & [1.727, 1.742] & 0.0002 & 0.0004 & -101.97 & 0.0020 & 0.7 \\
go:schubfach-vs-ryu & api & verify & canonical/medium & schubfach & ryu & schubfach & 1.01$\times$ & [1.001, 1.013] & 0.0424 & 0.0739 & -0.97 & 0.0020 & 0.7 \\
go:schubfach-vs-dragonbox & api & verify & canonical/medium & schubfach & dragonbox & schubfach & 1.00$\times$ & [0.997, 1.008] & 0.3855 & 0.4934 & -0.40 & 0.0020 & 0.7 \\
go:ryu-vs-dragonbox & api & verify & canonical/medium & ryu & dragonbox & dragonbox & 1.00$\times$ & [0.998, 1.011] & 0.2490 & 0.3486 & 0.55 & 0.0026 & 0.9 \\
go:ryu-vs-bd & api & verify & canonical/medium & ryu & json-canon & ryu$^*$ & 1.72$\times$ & [1.712, 1.732] & 0.0002 & 0.0004 & -84.79 & 0.0030 & 1.1 \\
go:dragonbox-vs-bd & api & verify & canonical/medium & dragonbox & json-canon & dragonbox$^*$ & 1.73$\times$ & [1.722, 1.740] & 0.0002 & 0.0004 & -92.10 & 0.0026 & 0.9 \\
rs:schubfach-vs-bd & api & verify & canonical/medium & schubfach-rs & json-canon-rs & schubfach-rs$^*$ & 3.02$\times$ & [2.954, 3.088] & 0.0002 & 0.0004 & -49.14 & 0.0060 & 4.2 \\
rs:schubfach-vs-ryu & api & verify & canonical/medium & schubfach-rs & ryu-rs & schubfach-rs & 1.02$\times$ & [0.990, 1.060] & 0.2168 & 0.3086 & -0.58 & 0.0060 & 4.2 \\
rs:schubfach-vs-dragonbox & api & verify & canonical/medium & schubfach-rs & dragonbox-rs & schubfach-rs & 1.01$\times$ & [0.993, 1.033] & 0.2743 & 0.3749 & -0.51 & 0.0060 & 4.2 \\
rs:ryu-vs-dragonbox & api & verify & canonical/medium & ryu-rs & dragonbox-rs & dragonbox-rs & 1.01$\times$ & [0.982, 1.038] & 0.4433 & 0.5542 & 0.34 & 0.0013 & 0.9 \\
rs:ryu-vs-bd & api & verify & canonical/medium & ryu-rs & json-canon-rs & ryu-rs$^*$ & 2.95$\times$ & [2.869, 3.040] & 0.0002 & 0.0004 & -41.88 & 0.0085 & 5.9 \\
rs:dragonbox-vs-bd & api & verify & canonical/medium & dragonbox-rs & json-canon-rs & dragonbox-rs$^*$ & 2.98$\times$ & [2.959, 3.024] & 0.0002 & 0.0004 & -59.17 & 0.0013 & 0.9 \\
schubfach:go-vs-rs & api & verify & canonical/medium & schubfach & schubfach-rs & schubfach-rs$^*$ & 1.91$\times$ & [1.873, 1.951] & 0.0002 & 0.0004 & 36.34 & 0.0060 & 4.2 \\
bd:go-vs-rs & api & verify & canonical/medium & json-canon & json-canon-rs & json-canon-rs$^*$ & 1.10$\times$ & [1.082, 1.106] & 0.0002 & 0.0004 & 8.46 & 0.0084 & 2.0 \\
ryu:go-vs-rs & api & verify & canonical/medium & ryu & ryu-rs & ryu-rs$^*$ & 1.88$\times$ & [1.831, 1.935] & 0.0002 & 0.0004 & 25.11 & 0.0085 & 5.9 \\
dragonbox:go-vs-rs & api & verify & canonical/medium & dragonbox & dragonbox-rs & dragonbox-rs$^*$ & 1.89$\times$ & [1.881, 1.904] & 0.0002 & 0.0004 & 78.59 & 0.0013 & 0.9 \\
go:schubfach-vs-bd & api & verify & canonical/mixed-prod & schubfach & json-canon & schubfach$^*$ & 1.34$\times$ & [1.336, 1.353] & 0.0002 & 0.0004 & -36.31 & 0.0009 & 0.9 \\
go:schubfach-vs-ryu & api & verify & canonical/mixed-prod & schubfach & ryu & schubfach & 1.01$\times$ & [1.000, 1.017] & 0.1504 & 0.2311 & -0.67 & 0.0009 & 0.9 \\
go:schubfach-vs-dragonbox & api & verify & canonical/mixed-prod & schubfach & dragonbox & schubfach & 1.01$\times$ & [1.001, 1.012] & 0.0556 & 0.0949 & -0.94 & 0.0009 & 0.9 \\
go:ryu-vs-dragonbox & api & verify & canonical/mixed-prod & ryu & dragonbox & dragonbox & 1.00$\times$ & [0.994, 1.011] & 0.8646 & 0.9030 & 0.08 & 0.0008 & 0.7 \\
go:ryu-vs-bd & api & verify & canonical/mixed-prod & ryu & json-canon & ryu$^*$ & 1.34$\times$ & [1.322, 1.346] & 0.0002 & 0.0004 & -28.22 & 0.0017 & 1.6 \\
go:dragonbox-vs-bd & api & verify & canonical/mixed-prod & dragonbox & json-canon & dragonbox$^*$ & 1.34$\times$ & [1.329, 1.345] & 0.0002 & 0.0004 & -37.67 & 0.0008 & 0.7 \\
rs:schubfach-vs-bd & api & verify & canonical/mixed-prod & schubfach-rs & json-canon-rs & schubfach-rs$^*$ & 1.83$\times$ & [1.807, 1.848] & 0.0002 & 0.0004 & -46.73 & 0.0011 & 2.0 \\
rs:schubfach-vs-ryu & api & verify & canonical/mixed-prod & schubfach-rs & ryu-rs & schubfach-rs & 1.01$\times$ & [0.997, 1.032] & 0.2166 & 0.3086 & -0.58 & 0.0011 & 2.0 \\
rs:schubfach-vs-dragonbox & api & verify & canonical/mixed-prod & schubfach-rs & dragonbox-rs & schubfach-rs & 1.01$\times$ & [0.994, 1.028] & 0.5185 & 0.6195 & -0.34 & 0.0011 & 2.0 \\
rs:ryu-vs-dragonbox & api & verify & canonical/mixed-prod & ryu-rs & dragonbox-rs & dragonbox-rs & 1.01$\times$ & [0.986, 1.026] & 0.5785 & 0.6756 & 0.26 & 0.0014 & 2.7 \\
rs:ryu-vs-bd & api & verify & canonical/mixed-prod & ryu-rs & json-canon-rs & ryu-rs$^*$ & 1.81$\times$ & [1.772, 1.832] & 0.0002 & 0.0004 & -35.58 & 0.0017 & 3.2 \\
rs:dragonbox-vs-bd & api & verify & canonical/mixed-prod & dragonbox-rs & json-canon-rs & dragonbox-rs$^*$ & 1.82$\times$ & [1.781, 1.834] & 0.0002 & 0.0004 & -40.37 & 0.0014 & 2.7 \\
schubfach:go-vs-rs & api & verify & canonical/mixed-prod & schubfach & schubfach-rs & schubfach-rs$^*$ & 2.00$\times$ & [1.981, 2.023] & 0.0002 & 0.0004 & 65.93 & 0.0011 & 2.0 \\
bd:go-vs-rs & api & verify & canonical/mixed-prod & json-canon & json-canon-rs & json-canon-rs$^*$ & 1.47$\times$ & [1.463, 1.486] & 0.0002 & 0.0004 & 41.14 & 0.0013 & 1.3 \\
ryu:go-vs-rs & api & verify & canonical/mixed-prod & ryu & ryu-rs & ryu-rs$^*$ & 1.99$\times$ & [1.954, 2.022] & 0.0002 & 0.0004 & 39.21 & 0.0017 & 3.2 \\
dragonbox:go-vs-rs & api & verify & canonical/mixed-prod & dragonbox & dragonbox-rs & dragonbox-rs$^*$ & 2.00$\times$ & [1.961, 2.021] & 0.0002 & 0.0004 & 58.81 & 0.0014 & 2.7 \\
go:schubfach-vs-bd & api & verify & canonical/nested-mixed & schubfach & json-canon & schubfach$^*$ & 1.41$\times$ & [1.398, 1.413] & 0.0002 & 0.0004 & -50.92 & 0.0000 & 0.9 \\
go:schubfach-vs-ryu & api & verify & canonical/nested-mixed & schubfach & ryu & schubfach$^*$ & 1.01$\times$ & [1.004, 1.016] & 0.0030 & 0.0061 & -1.58 & 0.0000 & 0.9 \\
go:schubfach-vs-dragonbox & api & verify & canonical/nested-mixed & schubfach & dragonbox & schubfach & 1.01$\times$ & [1.000, 1.010] & 0.0658 & 0.1116 & -0.91 & 0.0000 & 0.9 \\
go:ryu-vs-dragonbox & api & verify & canonical/nested-mixed & ryu & dragonbox & dragonbox & 1.01$\times$ & [1.000, 1.010] & 0.0696 & 0.1169 & 0.86 & 0.0000 & 0.7 \\
go:ryu-vs-bd & api & verify & canonical/nested-mixed & ryu & json-canon & ryu$^*$ & 1.39$\times$ & [1.384, 1.400] & 0.0002 & 0.0004 & -49.92 & 0.0000 & 0.8 \\
go:dragonbox-vs-bd & api & verify & canonical/nested-mixed & dragonbox & json-canon & dragonbox$^*$ & 1.40$\times$ & [1.392, 1.405] & 0.0002 & 0.0004 & -54.75 & 0.0000 & 0.7 \\
rs:schubfach-vs-bd & api & verify & canonical/nested-mixed & schubfach-rs & json-canon-rs & schubfach-rs$^*$ & 1.97$\times$ & [1.913, 2.010] & 0.0002 & 0.0004 & -26.02 & 0.0000 & 4.8 \\
rs:schubfach-vs-ryu & api & verify & canonical/nested-mixed & schubfach-rs & ryu-rs & schubfach-rs & 1.00$\times$ & [0.966, 1.039] & 0.9322 & 0.9557 & -0.04 & 0.0000 & 4.8 \\
rs:schubfach-vs-dragonbox & api & verify & canonical/nested-mixed & schubfach-rs & dragonbox-rs & schubfach-rs & 1.00$\times$ & [0.971, 1.026] & 0.9282 & 0.9540 & -0.04 & 0.0000 & 4.8 \\
rs:ryu-vs-dragonbox & api & verify & canonical/nested-mixed & ryu-rs & dragonbox-rs & dragonbox-rs & 1.00$\times$ & [0.971, 1.039] & 0.9722 & 0.9790 & 0.02 & 0.0000 & 3.6 \\
rs:ryu-vs-bd & api & verify & canonical/nested-mixed & ryu-rs & json-canon-rs & ryu-rs$^*$ & 1.96$\times$ & [1.895, 2.016] & 0.0002 & 0.0004 & -22.30 & 0.0000 & 6.2 \\
rs:dragonbox-vs-bd & api & verify & canonical/nested-mixed & dragonbox-rs & json-canon-rs & dragonbox-rs$^*$ & 1.96$\times$ & [1.931, 2.014] & 0.0002 & 0.0004 & -29.67 & 0.0000 & 3.6 \\
schubfach:go-vs-rs & api & verify & canonical/nested-mixed & schubfach & schubfach-rs & schubfach-rs$^*$ & 2.19$\times$ & [2.135, 2.234] & 0.0002 & 0.0004 & 40.79 & 0.0000 & 4.8 \\
bd:go-vs-rs & api & verify & canonical/nested-mixed & json-canon & json-canon-rs & json-canon-rs$^*$ & 1.57$\times$ & [1.548, 1.583] & 0.0002 & 0.0004 & 38.67 & 0.0000 & 2.3 \\
ryu:go-vs-rs & api & verify & canonical/nested-mixed & ryu & ryu-rs & ryu-rs$^*$ & 2.21$\times$ & [2.141, 2.267] & 0.0002 & 0.0004 & 33.14 & 0.0000 & 6.2 \\
dragonbox:go-vs-rs & api & verify & canonical/nested-mixed & dragonbox & dragonbox-rs & dragonbox-rs$^*$ & 2.20$\times$ & [2.172, 2.253] & 0.0002 & 0.0004 & 54.73 & 0.0000 & 3.6 \\
go:schubfach-vs-bd & api & verify & canonical/number-heavy & schubfach & json-canon & schubfach$^*$ & 1.56$\times$ & [1.546, 1.569] & 0.0002 & 0.0004 & -55.30 & 0.0002 & 1.4 \\
go:schubfach-vs-ryu & api & verify & canonical/number-heavy & schubfach & ryu & schubfach & 1.01$\times$ & [0.996, 1.016] & 0.3435 & 0.4468 & -0.44 & 0.0002 & 1.4 \\
go:schubfach-vs-dragonbox & api & verify & canonical/number-heavy & schubfach & dragonbox & dragonbox & 1.00$\times$ & [0.991, 1.008] & 0.8806 & 0.9145 & 0.07 & 0.0001 & 1.1 \\
go:ryu-vs-dragonbox & api & verify & canonical/number-heavy & ryu & dragonbox & dragonbox & 1.01$\times$ & [0.997, 1.016] & 0.2498 & 0.3494 & 0.54 & 0.0001 & 1.1 \\
go:ryu-vs-bd & api & verify & canonical/number-heavy & ryu & json-canon & ryu$^*$ & 1.55$\times$ & [1.534, 1.560] & 0.0002 & 0.0004 & -49.02 & 0.0002 & 1.6 \\
go:dragonbox-vs-bd & api & verify & canonical/number-heavy & dragonbox & json-canon & dragonbox$^*$ & 1.56$\times$ & [1.546, 1.566] & 0.0002 & 0.0004 & -64.32 & 0.0001 & 1.1 \\
rs:schubfach-vs-bd & api & verify & canonical/number-heavy & schubfach-rs & json-canon-rs & schubfach-rs$^*$ & 8.23$\times$ & [8.115, 8.394] & 0.0002 & 0.0004 & -136.63 & 0.0000 & 3.4 \\
rs:schubfach-vs-ryu & api & verify & canonical/number-heavy & schubfach-rs & ryu-rs & schubfach-rs & 1.02$\times$ & [0.998, 1.043] & 0.0838 & 0.1383 & -0.80 & 0.0000 & 3.4 \\
rs:schubfach-vs-dragonbox & api & verify & canonical/number-heavy & schubfach-rs & dragonbox-rs & dragonbox-rs & 1.00$\times$ & [0.974, 1.027] & 0.7910 & 0.8446 & 0.13 & 0.0000 & 4.3 \\
rs:ryu-vs-dragonbox & api & verify & canonical/number-heavy & ryu-rs & dragonbox-rs & dragonbox-rs & 1.02$\times$ & [0.997, 1.048] & 0.0816 & 0.1348 & 0.82 & 0.0000 & 4.3 \\
rs:ryu-vs-bd & api & verify & canonical/number-heavy & ryu-rs & json-canon-rs & ryu-rs$^*$ & 8.06$\times$ & [7.965, 8.225] & 0.0002 & 0.0004 & -137.57 & 0.0000 & 3.1 \\
rs:dragonbox-vs-bd & api & verify & canonical/number-heavy & dragonbox-rs & json-canon-rs & dragonbox-rs$^*$ & 8.26$\times$ & [8.069, 8.414] & 0.0002 & 0.0004 & -131.48 & 0.0000 & 4.3 \\
schubfach:go-vs-rs & api & verify & canonical/number-heavy & schubfach & schubfach-rs & schubfach-rs$^*$ & 9.81$\times$ & [9.666, 10.010] & 0.0002 & 0.0004 & 110.36 & 0.0000 & 3.4 \\
bd:go-vs-rs & api & verify & canonical/number-heavy & json-canon & json-canon-rs & json-canon-rs$^*$ & 1.86$\times$ & [1.844, 1.867] & 0.0002 & 0.0004 & 89.58 & 0.0001 & 1.1 \\
ryu:go-vs-rs & api & verify & canonical/number-heavy & ryu & ryu-rs & ryu-rs$^*$ & 9.66$\times$ & [9.531, 9.858] & 0.0002 & 0.0004 & 94.58 & 0.0000 & 3.1 \\
dragonbox:go-vs-rs & api & verify & canonical/number-heavy & dragonbox & dragonbox-rs & dragonbox-rs$^*$ & 9.84$\times$ & [9.601, 10.022] & 0.0002 & 0.0004 & 137.99 & 0.0000 & 4.3 \\
go:schubfach-vs-bd & api & verify & canonical/numeric-boundary & schubfach & json-canon & schubfach$^*$ & 1.49$\times$ & [1.475, 1.495] & 0.0002 & 0.0004 & -45.59 & 0.0001 & 1.0 \\
go:schubfach-vs-ryu & api & verify & canonical/numeric-boundary & schubfach & ryu & schubfach & 1.01$\times$ & [1.000, 1.019] & 0.0856 & 0.1409 & -0.84 & 0.0001 & 1.0 \\
go:schubfach-vs-dragonbox & api & verify & canonical/numeric-boundary & schubfach & dragonbox & schubfach & 1.01$\times$ & [0.998, 1.013] & 0.1872 & 0.2730 & -0.60 & 0.0001 & 1.0 \\
go:ryu-vs-dragonbox & api & verify & canonical/numeric-boundary & ryu & dragonbox & dragonbox & 1.00$\times$ & [0.994, 1.014] & 0.4471 & 0.5568 & 0.34 & 0.0001 & 1.2 \\
go:ryu-vs-bd & api & verify & canonical/numeric-boundary & ryu & json-canon & ryu$^*$ & 1.47$\times$ & [1.457, 1.485] & 0.0002 & 0.0004 & -35.63 & 0.0002 & 1.7 \\
go:dragonbox-vs-bd & api & verify & canonical/numeric-boundary & dragonbox & json-canon & dragonbox$^*$ & 1.48$\times$ & [1.465, 1.488] & 0.0002 & 0.0004 & -41.95 & 0.0001 & 1.2 \\
rs:schubfach-vs-bd & api & verify & canonical/numeric-boundary & schubfach-rs & json-canon-rs & schubfach-rs$^*$ & 9.46$\times$ & [9.328, 9.773] & 0.0002 & 0.0004 & -154.74 & 0.0000 & 4.2 \\
rs:schubfach-vs-ryu & api & verify & canonical/numeric-boundary & schubfach-rs & ryu-rs & schubfach-rs & 1.02$\times$ & [0.990, 1.051] & 0.1784 & 0.2623 & -0.64 & 0.0000 & 4.2 \\
rs:schubfach-vs-dragonbox & api & verify & canonical/numeric-boundary & schubfach-rs & dragonbox-rs & schubfach-rs & 1.01$\times$ & [0.988, 1.037] & 0.6219 & 0.7109 & -0.25 & 0.0000 & 4.2 \\
rs:ryu-vs-dragonbox & api & verify & canonical/numeric-boundary & ryu-rs & dragonbox-rs & dragonbox-rs & 1.02$\times$ & [0.979, 1.036] & 0.2879 & 0.3873 & 0.51 & 0.0000 & 2.7 \\
rs:ryu-vs-bd & api & verify & canonical/numeric-boundary & ryu-rs & json-canon-rs & ryu-rs$^*$ & 9.25$\times$ & [9.107, 9.587] & 0.0002 & 0.0004 & -149.96 & 0.0000 & 4.7 \\
rs:dragonbox-vs-bd & api & verify & canonical/numeric-boundary & dragonbox-rs & json-canon-rs & dragonbox-rs$^*$ & 9.39$\times$ & [9.261, 9.521] & 0.0002 & 0.0004 & -163.23 & 0.0000 & 2.7 \\
schubfach:go-vs-rs & api & verify & canonical/numeric-boundary & schubfach & schubfach-rs & schubfach-rs$^*$ & 13.33$\times$ & [13.142, 13.737] & 0.0002 & 0.0004 & 162.35 & 0.0000 & 4.2 \\
bd:go-vs-rs & api & verify & canonical/numeric-boundary & json-canon & json-canon-rs & json-canon-rs$^*$ & 2.09$\times$ & [2.077, 2.106] & 0.0002 & 0.0004 & 78.51 & 0.0001 & 0.9 \\
ryu:go-vs-rs & api & verify & canonical/numeric-boundary & ryu & ryu-rs & ryu-rs$^*$ & 13.15$\times$ & [12.934, 13.603] & 0.0002 & 0.0004 & 93.92 & 0.0000 & 4.7 \\
dragonbox:go-vs-rs & api & verify & canonical/numeric-boundary & dragonbox & dragonbox-rs & dragonbox-rs$^*$ & 13.31$\times$ & [13.112, 13.488] & 0.0002 & 0.0004 & 133.71 & 0.0000 & 2.7 \\
go:schubfach-vs-bd & api & verify & canonical/rfc-key-sorting & schubfach & json-canon & schubfach$^*$ & 1.01$\times$ & [1.002, 1.012] & 0.0230 & 0.0423 & -1.17 & 0.0000 & 0.8 \\
go:schubfach-vs-ryu & api & verify & canonical/rfc-key-sorting & schubfach & ryu & schubfach$^*$ & 1.01$\times$ & [1.002, 1.012] & 0.0166 & 0.0310 & -1.21 & 0.0000 & 0.8 \\
go:schubfach-vs-dragonbox & api & verify & canonical/rfc-key-sorting & schubfach & dragonbox & schubfach$^*$ & 1.01$\times$ & [1.003, 1.015] & 0.0096 & 0.0184 & -1.30 & 0.0000 & 0.8 \\
go:ryu-vs-dragonbox & api & verify & canonical/rfc-key-sorting & ryu & dragonbox & ryu & 1.00$\times$ & [0.996, 1.008] & 0.4919 & 0.5949 & -0.32 & 0.0000 & 0.7 \\
go:ryu-vs-bd & api & verify & canonical/rfc-key-sorting & ryu & json-canon & json-canon & 1.00$\times$ & [0.995, 1.005] & 0.9706 & 0.9789 & 0.01 & 0.0000 & 0.7 \\
go:dragonbox-vs-bd & api & verify & canonical/rfc-key-sorting & dragonbox & json-canon & json-canon & 1.00$\times$ & [0.996, 1.008] & 0.4775 & 0.5807 & 0.32 & 0.0000 & 0.7 \\
rs:schubfach-vs-bd & api & verify & canonical/rfc-key-sorting & schubfach-rs & json-canon-rs & json-canon-rs & 1.01$\times$ & [0.988, 1.043] & 0.3265 & 0.4277 & 0.45 & 0.0000 & 4.0 \\
rs:schubfach-vs-ryu & api & verify & canonical/rfc-key-sorting & schubfach-rs & ryu-rs & schubfach-rs$^*$ & 1.05$\times$ & [1.022, 1.076] & 0.0036 & 0.0072 & -1.56 & 0.0000 & 4.2 \\
rs:schubfach-vs-dragonbox & api & verify & canonical/rfc-key-sorting & schubfach-rs & dragonbox-rs & schubfach-rs$^*$ & 1.06$\times$ & [1.025, 1.107] & 0.0110 & 0.0210 & -1.25 & 0.0000 & 4.2 \\
rs:ryu-vs-dragonbox & api & verify & canonical/rfc-key-sorting & ryu-rs & dragonbox-rs & ryu-rs & 1.01$\times$ & [0.977, 1.051] & 0.5831 & 0.6791 & -0.25 & 0.0000 & 3.5 \\
rs:ryu-vs-bd & api & verify & canonical/rfc-key-sorting & ryu-rs & json-canon-rs & json-canon-rs$^*$ & 1.06$\times$ & [1.039, 1.090] & 0.0002 & 0.0004 & 2.11 & 0.0000 & 4.0 \\
rs:dragonbox-vs-bd & api & verify & canonical/rfc-key-sorting & dragonbox-rs & json-canon-rs & json-canon-rs$^*$ & 1.08$\times$ & [1.038, 1.121] & 0.0024 & 0.0049 & 1.58 & 0.0000 & 4.0 \\
schubfach:go-vs-rs & api & verify & canonical/rfc-key-sorting & schubfach & schubfach-rs & schubfach-rs$^*$ & 2.15$\times$ & [2.101, 2.188] & 0.0002 & 0.0004 & 44.64 & 0.0000 & 4.2 \\
bd:go-vs-rs & api & verify & canonical/rfc-key-sorting & json-canon & json-canon-rs & json-canon-rs$^*$ & 2.19$\times$ & [2.153, 2.236] & 0.0002 & 0.0004 & 49.28 & 0.0000 & 4.0 \\
ryu:go-vs-rs & api & verify & canonical/rfc-key-sorting & ryu & ryu-rs & ryu-rs$^*$ & 2.06$\times$ & [2.028, 2.096] & 0.0002 & 0.0004 & 50.00 & 0.0000 & 3.5 \\
dragonbox:go-vs-rs & api & verify & canonical/rfc-key-sorting & dragonbox & dragonbox-rs & dragonbox-rs$^*$ & 2.04$\times$ & [1.967, 2.107] & 0.0002 & 0.0004 & 24.95 & 0.0001 & 7.1 \\
go:schubfach-vs-bd & api & verify & canonical/small & schubfach & json-canon & schubfach$^*$ & 1.35$\times$ & [1.344, 1.360] & 0.0002 & 0.0004 & -43.68 & 0.0000 & 1.0 \\
go:schubfach-vs-ryu & api & verify & canonical/small & schubfach & ryu & schubfach & 1.01$\times$ & [1.003, 1.016] & 0.0278 & 0.0504 & -1.08 & 0.0000 & 1.0 \\
go:schubfach-vs-dragonbox & api & verify & canonical/small & schubfach & dragonbox & schubfach & 1.00$\times$ & [0.998, 1.010] & 0.1806 & 0.2652 & -0.62 & 0.0000 & 1.0 \\
go:ryu-vs-dragonbox & api & verify & canonical/small & ryu & dragonbox & dragonbox & 1.00$\times$ & [0.999, 1.011] & 0.2092 & 0.3001 & 0.59 & 0.0000 & 0.8 \\
go:ryu-vs-bd & api & verify & canonical/small & ryu & json-canon & ryu$^*$ & 1.34$\times$ & [1.333, 1.348] & 0.0002 & 0.0004 & -44.13 & 0.0000 & 1.0 \\
go:dragonbox-vs-bd & api & verify & canonical/small & dragonbox & json-canon & dragonbox$^*$ & 1.35$\times$ & [1.340, 1.353] & 0.0002 & 0.0004 & -49.00 & 0.0000 & 0.8 \\
rs:schubfach-vs-bd & api & verify & canonical/small & schubfach-rs & json-canon-rs & schubfach-rs$^*$ & 1.73$\times$ & [1.654, 1.795] & 0.0002 & 0.0004 & -10.36 & 0.0000 & 6.4 \\
rs:schubfach-vs-ryu & api & verify & canonical/small & schubfach-rs & ryu-rs & schubfach-rs$^*$ & 1.08$\times$ & [1.033, 1.118] & 0.0046 & 0.0090 & -1.61 & 0.0000 & 6.4 \\
rs:schubfach-vs-dragonbox & api & verify & canonical/small & schubfach-rs & dragonbox-rs & schubfach-rs & 1.03$\times$ & [0.985, 1.069] & 0.2511 & 0.3506 & -0.53 & 0.0000 & 6.4 \\
rs:ryu-vs-dragonbox & api & verify & canonical/small & ryu-rs & dragonbox-rs & dragonbox-rs$^*$ & 1.05$\times$ & [1.008, 1.088] & 0.0264 & 0.0481 & 1.10 & 0.0000 & 5.9 \\
rs:ryu-vs-bd & api & verify & canonical/small & ryu-rs & json-canon-rs & ryu-rs$^*$ & 1.60$\times$ & [1.545, 1.668] & 0.0002 & 0.0004 & -9.48 & 0.0000 & 5.3 \\
rs:dragonbox-vs-bd & api & verify & canonical/small & dragonbox-rs & json-canon-rs & dragonbox-rs$^*$ & 1.68$\times$ & [1.611, 1.747] & 0.0002 & 0.0004 & -10.12 & 0.0000 & 5.9 \\
schubfach:go-vs-rs & api & verify & canonical/small & schubfach & schubfach-rs & schubfach-rs$^*$ & 1.84$\times$ & [1.780, 1.890] & 0.0002 & 0.0004 & 22.32 & 0.0000 & 6.4 \\
bd:go-vs-rs & api & verify & canonical/small & json-canon & json-canon-rs & json-canon-rs$^*$ & 1.44$\times$ & [1.399, 1.485] & 0.0002 & 0.0004 & 12.50 & 0.0001 & 6.2 \\
ryu:go-vs-rs & api & verify & canonical/small & ryu & ryu-rs & ryu-rs$^*$ & 1.72$\times$ & [1.681, 1.769] & 0.0002 & 0.0004 & 22.86 & 0.0000 & 5.3 \\
dragonbox:go-vs-rs & api & verify & canonical/small & dragonbox & dragonbox-rs & dragonbox-rs$^*$ & 1.80$\times$ & [1.746, 1.846] & 0.0002 & 0.0004 & 23.38 & 0.0000 & 5.9 \\
go:schubfach-vs-bd & api & verify & canonical/surrogate-pair & schubfach & json-canon & schubfach & 1.01$\times$ & [1.000, 1.009] & 0.0408 & 0.0715 & -0.93 & 0.0000 & 0.8 \\
go:schubfach-vs-ryu & api & verify & canonical/surrogate-pair & schubfach & ryu & schubfach & 1.00$\times$ & [0.995, 1.006] & 0.6767 & 0.7555 & -0.18 & 0.0000 & 0.8 \\
go:schubfach-vs-dragonbox & api & verify & canonical/surrogate-pair & schubfach & dragonbox & dragonbox & 1.00$\times$ & [0.996, 1.010] & 0.7842 & 0.8394 & 0.15 & 0.0000 & 1.1 \\
go:ryu-vs-dragonbox & api & verify & canonical/surrogate-pair & ryu & dragonbox & dragonbox & 1.00$\times$ & [0.996, 1.010] & 0.5503 & 0.6476 & 0.30 & 0.0000 & 1.1 \\
go:ryu-vs-bd & api & verify & canonical/surrogate-pair & ryu & json-canon & ryu & 1.00$\times$ & [0.999, 1.009] & 0.1556 & 0.2368 & -0.69 & 0.0000 & 0.9 \\
go:dragonbox-vs-bd & api & verify & canonical/surrogate-pair & dragonbox & json-canon & dragonbox & 1.01$\times$ & [1.002, 1.014] & 0.0404 & 0.0709 & -0.91 & 0.0000 & 1.1 \\
rs:schubfach-vs-bd & api & verify & canonical/surrogate-pair & schubfach-rs & json-canon-rs & json-canon-rs & 1.01$\times$ & [0.970, 1.062] & 0.6843 & 0.7613 & 0.20 & 0.0000 & 5.0 \\
rs:schubfach-vs-ryu & api & verify & canonical/surrogate-pair & schubfach-rs & ryu-rs & ryu-rs & 1.03$\times$ & [0.993, 1.081] & 0.2312 & 0.3273 & 0.58 & 0.0000 & 3.9 \\
rs:schubfach-vs-dragonbox & api & verify & canonical/surrogate-pair & schubfach-rs & dragonbox-rs & schubfach-rs & 1.08$\times$ & [1.016, 1.137] & 0.0306 & 0.0554 & -1.04 & 0.0000 & 8.0 \\
rs:ryu-vs-dragonbox & api & verify & canonical/surrogate-pair & ryu-rs & dragonbox-rs & ryu-rs$^*$ & 1.11$\times$ & [1.062, 1.165] & 0.0006 & 0.0013 & -1.72 & 0.0000 & 3.9 \\
rs:ryu-vs-bd & api & verify & canonical/surrogate-pair & ryu-rs & json-canon-rs & ryu-rs & 1.02$\times$ & [0.994, 1.056] & 0.2685 & 0.3701 & -0.53 & 0.0000 & 3.9 \\
rs:dragonbox-vs-bd & api & verify & canonical/surrogate-pair & dragonbox-rs & json-canon-rs & json-canon-rs$^*$ & 1.09$\times$ & [1.035, 1.146] & 0.0078 & 0.0151 & 1.35 & 0.0000 & 5.0 \\
schubfach:go-vs-rs & api & verify & canonical/surrogate-pair & schubfach & schubfach-rs & schubfach-rs$^*$ & 2.11$\times$ & [2.013, 2.177] & 0.0002 & 0.0004 & 23.97 & 0.0000 & 8.0 \\
bd:go-vs-rs & api & verify & canonical/surrogate-pair & json-canon & json-canon-rs & json-canon-rs$^*$ & 2.14$\times$ & [2.067, 2.178] & 0.0002 & 0.0004 & 39.70 & 0.0000 & 5.0 \\
ryu:go-vs-rs & api & verify & canonical/surrogate-pair & ryu & ryu-rs & ryu-rs$^*$ & 2.18$\times$ & [2.128, 2.211] & 0.0002 & 0.0004 & 47.86 & 0.0000 & 3.9 \\
dragonbox:go-vs-rs & api & verify & canonical/surrogate-pair & dragonbox & dragonbox-rs & dragonbox-rs$^*$ & 1.96$\times$ & [1.860, 2.035] & 0.0002 & 0.0004 & 17.59 & 0.0000 & 9.4 \\
go:schubfach-vs-bd & api & verify & canonical/unicode & schubfach & json-canon & json-canon & 1.00$\times$ & [0.999, 1.009] & 0.1526 & 0.2336 & 0.66 & 0.0000 & 0.6 \\
go:schubfach-vs-ryu & api & verify & canonical/unicode & schubfach & ryu & schubfach & 1.01$\times$ & [1.001, 1.016] & 0.0292 & 0.0529 & -1.06 & 0.0000 & 0.8 \\
go:schubfach-vs-dragonbox & api & verify & canonical/unicode & schubfach & dragonbox & schubfach & 1.01$\times$ & [1.001, 1.014] & 0.0390 & 0.0689 & -1.01 & 0.0000 & 0.8 \\
go:ryu-vs-dragonbox & api & verify & canonical/unicode & ryu & dragonbox & dragonbox & 1.00$\times$ & [0.993, 1.009] & 0.6873 & 0.7621 & 0.19 & 0.0000 & 1.1 \\
go:ryu-vs-bd & api & verify & canonical/unicode & ryu & json-canon & json-canon$^*$ & 1.01$\times$ & [1.005, 1.020] & 0.0034 & 0.0068 & 1.53 & 0.0000 & 0.6 \\
go:dragonbox-vs-bd & api & verify & canonical/unicode & dragonbox & json-canon & json-canon$^*$ & 1.01$\times$ & [1.006, 1.018] & 0.0020 & 0.0041 & 1.57 & 0.0000 & 0.6 \\
rs:schubfach-vs-bd & api & verify & canonical/unicode & schubfach-rs & json-canon-rs & schubfach-rs$^*$ & 1.02$\times$ & [1.010, 1.045] & 0.0098 & 0.0188 & -1.16 & 0.0000 & 1.9 \\
rs:schubfach-vs-ryu & api & verify & canonical/unicode & schubfach-rs & ryu-rs & ryu-rs & 1.00$\times$ & [0.990, 1.014] & 0.9304 & 0.9548 & 0.04 & 0.0000 & 1.7 \\
rs:schubfach-vs-dragonbox & api & verify & canonical/unicode & schubfach-rs & dragonbox-rs & schubfach-rs$^*$ & 1.08$\times$ & [1.066, 1.102] & 0.0002 & 0.0004 & -3.88 & 0.0000 & 1.9 \\
rs:ryu-vs-dragonbox & api & verify & canonical/unicode & ryu-rs & dragonbox-rs & ryu-rs$^*$ & 1.08$\times$ & [1.067, 1.103] & 0.0002 & 0.0004 & -3.99 & 0.0000 & 1.7 \\
rs:ryu-vs-bd & api & verify & canonical/unicode & ryu-rs & json-canon-rs & ryu-rs$^*$ & 1.03$\times$ & [1.011, 1.048] & 0.0082 & 0.0158 & -1.21 & 0.0000 & 1.7 \\
rs:dragonbox-vs-bd & api & verify & canonical/unicode & dragonbox-rs & json-canon-rs & json-canon-rs$^*$ & 1.06$\times$ & [1.033, 1.077] & 0.0002 & 0.0004 & 2.24 & 0.0000 & 3.2 \\
schubfach:go-vs-rs & api & verify & canonical/unicode & schubfach & schubfach-rs & schubfach-rs$^*$ & 1.03$\times$ & [1.023, 1.043] & 0.0004 & 0.0008 & 2.97 & 0.0000 & 1.9 \\
bd:go-vs-rs & api & verify & canonical/unicode & json-canon & json-canon-rs & json-canon-rs & 1.01$\times$ & [0.986, 1.018] & 0.5167 & 0.6179 & 0.33 & 0.0000 & 3.2 \\
ryu:go-vs-rs & api & verify & canonical/unicode & ryu & ryu-rs & ryu-rs$^*$ & 1.05$\times$ & [1.034, 1.056] & 0.0004 & 0.0008 & 3.52 & 0.0000 & 1.7 \\
dragonbox:go-vs-rs & api & verify & canonical/unicode & dragonbox & dragonbox-rs & dragonbox$^*$ & 1.04$\times$ & [1.024, 1.056] & 0.0002 & 0.0004 & -2.06 & 0.0000 & 1.1 \\
go:schubfach-vs-bd & api & verify & canonical/verify-whitespace & schubfach & json-canon & schubfach$^*$ & 1.71$\times$ & [1.698, 1.717] & 0.0004 & 0.0008 & -72.52 & 0.0000 & 0.7 \\
go:schubfach-vs-ryu & api & verify & canonical/verify-whitespace & schubfach & ryu & schubfach & 1.01$\times$ & [1.001, 1.011] & 0.0546 & 0.0934 & -0.91 & 0.0000 & 0.7 \\
go:schubfach-vs-dragonbox & api & verify & canonical/verify-whitespace & schubfach & dragonbox & schubfach$^*$ & 1.01$\times$ & [1.007, 1.019] & 0.0004 & 0.0008 & -1.92 & 0.0000 & 0.7 \\
go:ryu-vs-dragonbox & api & verify & canonical/verify-whitespace & ryu & dragonbox & ryu & 1.01$\times$ & [1.001, 1.013] & 0.0318 & 0.0573 & -1.03 & 0.0000 & 0.8 \\
go:ryu-vs-bd & api & verify & canonical/verify-whitespace & ryu & json-canon & ryu$^*$ & 1.70$\times$ & [1.687, 1.707] & 0.0002 & 0.0004 & -69.47 & 0.0000 & 0.8 \\
go:dragonbox-vs-bd & api & verify & canonical/verify-whitespace & dragonbox & json-canon & dragonbox$^*$ & 1.68$\times$ & [1.674, 1.696] & 0.0002 & 0.0004 & -65.86 & 0.0000 & 1.0 \\
rs:schubfach-vs-bd & api & verify & canonical/verify-whitespace & schubfach-rs & json-canon-rs & schubfach-rs$^*$ & 2.93$\times$ & [2.801, 3.004] & 0.0002 & 0.0004 & -33.38 & 0.0000 & 6.6 \\
rs:schubfach-vs-ryu & api & verify & canonical/verify-whitespace & schubfach-rs & ryu-rs & schubfach-rs & 1.01$\times$ & [0.965, 1.046] & 0.5223 & 0.6222 & -0.32 & 0.0000 & 6.6 \\
rs:schubfach-vs-dragonbox & api & verify & canonical/verify-whitespace & schubfach-rs & dragonbox-rs & schubfach-rs & 1.03$\times$ & [0.982, 1.060] & 0.0860 & 0.1414 & -0.82 & 0.0000 & 6.6 \\
rs:ryu-vs-dragonbox & api & verify & canonical/verify-whitespace & ryu-rs & dragonbox-rs & ryu-rs & 1.02$\times$ & [0.992, 1.044] & 0.1928 & 0.2803 & -0.61 & 0.0000 & 4.5 \\
rs:ryu-vs-bd & api & verify & canonical/verify-whitespace & ryu-rs & json-canon-rs & ryu-rs$^*$ & 2.89$\times$ & [2.818, 2.965] & 0.0002 & 0.0004 & -37.34 & 0.0000 & 4.5 \\
rs:dragonbox-vs-bd & api & verify & canonical/verify-whitespace & dragonbox-rs & json-canon-rs & dragonbox-rs$^*$ & 2.84$\times$ & [2.788, 2.902] & 0.0002 & 0.0004 & -39.48 & 0.0000 & 3.2 \\
schubfach:go-vs-rs & api & verify & canonical/verify-whitespace & schubfach & schubfach-rs & schubfach-rs$^*$ & 2.38$\times$ & [2.273, 2.434] & 0.0002 & 0.0004 & 35.97 & 0.0000 & 6.6 \\
bd:go-vs-rs & api & verify & canonical/verify-whitespace & json-canon & json-canon-rs & json-canon-rs$^*$ & 1.39$\times$ & [1.367, 1.404] & 0.0002 & 0.0004 & 23.08 & 0.0000 & 2.7 \\
ryu:go-vs-rs & api & verify & canonical/verify-whitespace & ryu & ryu-rs & ryu-rs$^*$ & 2.36$\times$ & [2.309, 2.413] & 0.0002 & 0.0004 & 48.72 & 0.0000 & 4.5 \\
dragonbox:go-vs-rs & api & verify & canonical/verify-whitespace & dragonbox & dragonbox-rs & dragonbox-rs$^*$ & 2.33$\times$ & [2.306, 2.388] & 0.0002 & 0.0004 & 59.94 & 0.0000 & 3.2 \\
go:schubfach-vs-bd & api & verify & noncanonical/control-escapes & schubfach & json-canon & schubfach$^*$ & 1.03$\times$ & [1.028, 1.042] & 0.0002 & 0.0004 & -3.96 & 0.0000 & 0.8 \\
go:schubfach-vs-ryu & api & verify & noncanonical/control-escapes & schubfach & ryu & schubfach$^*$ & 1.01$\times$ & [1.006, 1.021] & 0.0032 & 0.0064 & -1.62 & 0.0000 & 0.8 \\
go:schubfach-vs-dragonbox & api & verify & noncanonical/control-escapes & schubfach & dragonbox & schubfach$^*$ & 1.01$\times$ & [1.005, 1.015] & 0.0016 & 0.0033 & -1.89 & 0.0000 & 0.8 \\
go:ryu-vs-dragonbox & api & verify & noncanonical/control-escapes & ryu & dragonbox & dragonbox & 1.00$\times$ & [0.997, 1.010] & 0.3505 & 0.4550 & 0.43 & 0.0000 & 0.6 \\
go:ryu-vs-bd & api & verify & noncanonical/control-escapes & ryu & json-canon & ryu$^*$ & 1.02$\times$ & [1.012, 1.030] & 0.0004 & 0.0008 & -1.97 & 0.0000 & 1.3 \\
go:dragonbox-vs-bd & api & verify & noncanonical/control-escapes & dragonbox & json-canon & dragonbox$^*$ & 1.02$\times$ & [1.018, 1.031] & 0.0002 & 0.0004 & -2.95 & 0.0000 & 0.6 \\
go:schubfach-vs-bd & api & verify & noncanonical/escaped-key-order & schubfach & json-canon & schubfach$^*$ & 1.62$\times$ & [1.612, 1.627] & 0.0002 & 0.0004 & -84.30 & 0.0000 & 0.7 \\
go:schubfach-vs-ryu & api & verify & noncanonical/escaped-key-order & schubfach & ryu & schubfach$^*$ & 1.01$\times$ & [1.009, 1.020] & 0.0002 & 0.0004 & -2.36 & 0.0000 & 0.7 \\
go:schubfach-vs-dragonbox & api & verify & noncanonical/escaped-key-order & schubfach & dragonbox & schubfach$^*$ & 1.01$\times$ & [1.002, 1.018] & 0.0274 & 0.0498 & -1.14 & 0.0000 & 0.7 \\
go:ryu-vs-dragonbox & api & verify & noncanonical/escaped-key-order & ryu & dragonbox & dragonbox & 1.00$\times$ & [0.996, 1.013] & 0.4033 & 0.5107 & 0.36 & 0.0000 & 1.6 \\
go:ryu-vs-bd & api & verify & noncanonical/escaped-key-order & ryu & json-canon & ryu$^*$ & 1.60$\times$ & [1.589, 1.605] & 0.0002 & 0.0004 & -78.20 & 0.0000 & 0.8 \\
go:dragonbox-vs-bd & api & verify & noncanonical/escaped-key-order & dragonbox & json-canon & dragonbox$^*$ & 1.60$\times$ & [1.590, 1.617] & 0.0002 & 0.0004 & -56.17 & 0.0000 & 1.6 \\
go:schubfach-vs-bd & api & verify & noncanonical/large & schubfach & json-canon & schubfach$^*$ & 1.77$\times$ & [1.762, 1.774] & 0.0002 & 0.0004 & -141.67 & 0.0650 & 0.7 \\
go:schubfach-vs-ryu & api & verify & noncanonical/large & schubfach & ryu & schubfach & 1.01$\times$ & [1.000, 1.011] & 0.0734 & 0.1225 & -0.86 & 0.0650 & 0.7 \\
go:schubfach-vs-dragonbox & api & verify & noncanonical/large & schubfach & dragonbox & schubfach$^*$ & 1.01$\times$ & [1.005, 1.016] & 0.0032 & 0.0064 & -1.60 & 0.0650 & 0.7 \\
go:ryu-vs-dragonbox & api & verify & noncanonical/large & ryu & dragonbox & ryu & 1.01$\times$ & [0.998, 1.011] & 0.1454 & 0.2245 & -0.68 & 0.0878 & 0.9 \\
go:ryu-vs-bd & api & verify & noncanonical/large & ryu & json-canon & ryu$^*$ & 1.76$\times$ & [1.750, 1.766] & 0.0002 & 0.0004 & -117.91 & 0.0878 & 0.9 \\
go:dragonbox-vs-bd & api & verify & noncanonical/large & dragonbox & json-canon & dragonbox$^*$ & 1.75$\times$ & [1.742, 1.759] & 0.0002 & 0.0004 & -113.63 & 0.0919 & 1.0 \\
go:schubfach-vs-bd & api & verify & noncanonical/medium & schubfach & json-canon & schubfach$^*$ & 1.73$\times$ & [1.721, 1.739] & 0.0002 & 0.0004 & -115.83 & 0.0028 & 1.0 \\
go:schubfach-vs-ryu & api & verify & noncanonical/medium & schubfach & ryu & ryu & 1.00$\times$ & [0.995, 1.007] & 0.8732 & 0.9098 & 0.07 & 0.0021 & 0.8 \\
go:schubfach-vs-dragonbox & api & verify & noncanonical/medium & schubfach & dragonbox & dragonbox & 1.00$\times$ & [0.995, 1.007] & 0.9774 & 0.9820 & 0.01 & 0.0021 & 0.8 \\
go:ryu-vs-dragonbox & api & verify & noncanonical/medium & ryu & dragonbox & ryu & 1.00$\times$ & [0.994, 1.005] & 0.8852 & 0.9173 & -0.07 & 0.0021 & 0.8 \\
go:ryu-vs-bd & api & verify & noncanonical/medium & ryu & json-canon & ryu$^*$ & 1.73$\times$ & [1.725, 1.738] & 0.0002 & 0.0004 & -144.51 & 0.0021 & 0.8 \\
go:dragonbox-vs-bd & api & verify & noncanonical/medium & dragonbox & json-canon & dragonbox$^*$ & 1.73$\times$ & [1.726, 1.741] & 0.0002 & 0.0004 & -147.16 & 0.0021 & 0.8 \\
go:schubfach-vs-bd & api & verify & noncanonical/mixed-prod & schubfach & json-canon & schubfach$^*$ & 1.34$\times$ & [1.336, 1.352] & 0.0002 & 0.0004 & -41.03 & 0.0008 & 0.8 \\
go:schubfach-vs-ryu & api & verify & noncanonical/mixed-prod & schubfach & ryu & schubfach$^*$ & 1.01$\times$ & [1.005, 1.020] & 0.0070 & 0.0136 & -1.25 & 0.0008 & 0.8 \\
go:schubfach-vs-dragonbox & api & verify & noncanonical/mixed-prod & schubfach & dragonbox & schubfach$^*$ & 1.01$\times$ & [1.004, 1.019] & 0.0114 & 0.0217 & -1.30 & 0.0008 & 0.8 \\
go:ryu-vs-dragonbox & api & verify & noncanonical/mixed-prod & ryu & dragonbox & ryu & 1.00$\times$ & [0.990, 1.008] & 0.8900 & 0.9214 & -0.07 & 0.0014 & 1.3 \\
go:ryu-vs-bd & api & verify & noncanonical/mixed-prod & ryu & json-canon & ryu$^*$ & 1.33$\times$ & [1.317, 1.337] & 0.0002 & 0.0004 & -32.18 & 0.0014 & 1.3 \\
go:dragonbox-vs-bd & api & verify & noncanonical/mixed-prod & dragonbox & json-canon & dragonbox$^*$ & 1.33$\times$ & [1.318, 1.339] & 0.0004 & 0.0008 & -31.58 & 0.0015 & 1.4 \\
go:schubfach-vs-bd & api & verify & noncanonical/nested-mixed & schubfach & json-canon & schubfach$^*$ & 1.41$\times$ & [1.399, 1.412] & 0.0002 & 0.0004 & -59.97 & 0.0000 & 0.7 \\
go:schubfach-vs-ryu & api & verify & noncanonical/nested-mixed & schubfach & ryu & schubfach$^*$ & 1.01$\times$ & [1.005, 1.014] & 0.0028 & 0.0057 & -1.61 & 0.0000 & 0.7 \\
go:schubfach-vs-dragonbox & api & verify & noncanonical/nested-mixed & schubfach & dragonbox & schubfach$^*$ & 1.01$\times$ & [1.003, 1.011] & 0.0062 & 0.0121 & -1.40 & 0.0000 & 0.7 \\
go:ryu-vs-dragonbox & api & verify & noncanonical/nested-mixed & ryu & dragonbox & dragonbox & 1.00$\times$ & [0.999, 1.007] & 0.3251 & 0.4267 & 0.46 & 0.0000 & 0.5 \\
go:ryu-vs-bd & api & verify & noncanonical/nested-mixed & ryu & json-canon & ryu$^*$ & 1.39$\times$ & [1.386, 1.399] & 0.0002 & 0.0004 & -58.69 & 0.0000 & 0.7 \\
go:dragonbox-vs-bd & api & verify & noncanonical/nested-mixed & dragonbox & json-canon & dragonbox$^*$ & 1.40$\times$ & [1.391, 1.402] & 0.0002 & 0.0004 & -65.12 & 0.0000 & 0.5 \\
go:schubfach-vs-bd & api & verify & noncanonical/number-heavy & schubfach & json-canon & schubfach$^*$ & 1.56$\times$ & [1.542, 1.573] & 0.0004 & 0.0008 & -38.46 & 0.0002 & 1.7 \\
go:schubfach-vs-ryu & api & verify & noncanonical/number-heavy & schubfach & ryu & schubfach & 1.00$\times$ & [0.992, 1.010] & 0.5825 & 0.6790 & -0.26 & 0.0002 & 1.7 \\
go:schubfach-vs-dragonbox & api & verify & noncanonical/number-heavy & schubfach & dragonbox & schubfach & 1.00$\times$ & [0.991, 1.009] & 0.7149 & 0.7839 & -0.18 & 0.0002 & 1.7 \\
go:ryu-vs-dragonbox & api & verify & noncanonical/number-heavy & ryu & dragonbox & dragonbox & 1.00$\times$ & [0.996, 1.007] & 0.7822 & 0.8391 & 0.13 & 0.0001 & 0.8 \\
go:ryu-vs-bd & api & verify & noncanonical/number-heavy & ryu & json-canon & ryu$^*$ & 1.56$\times$ & [1.545, 1.566] & 0.0002 & 0.0004 & -47.80 & 0.0001 & 0.7 \\
go:dragonbox-vs-bd & api & verify & noncanonical/number-heavy & dragonbox & json-canon & dragonbox$^*$ & 1.56$\times$ & [1.547, 1.568] & 0.0002 & 0.0004 & -46.94 & 0.0001 & 0.8 \\
go:schubfach-vs-bd & api & verify & noncanonical/numeric-boundary & schubfach & json-canon & schubfach$^*$ & 1.50$\times$ & [1.485, 1.505] & 0.0002 & 0.0004 & -44.65 & 0.0001 & 0.8 \\
go:schubfach-vs-ryu & api & verify & noncanonical/numeric-boundary & schubfach & ryu & schubfach & 1.00$\times$ & [0.995, 1.007] & 0.6537 & 0.7394 & -0.21 & 0.0001 & 0.8 \\
go:schubfach-vs-dragonbox & api & verify & noncanonical/numeric-boundary & schubfach & dragonbox & schubfach & 1.00$\times$ & [0.993, 1.007] & 0.9622 & 0.9750 & -0.03 & 0.0001 & 0.8 \\
go:ryu-vs-dragonbox & api & verify & noncanonical/numeric-boundary & ryu & dragonbox & dragonbox & 1.00$\times$ & [0.994, 1.009] & 0.7706 & 0.8303 & 0.13 & 0.0001 & 1.4 \\
go:ryu-vs-bd & api & verify & noncanonical/numeric-boundary & ryu & json-canon & ryu$^*$ & 1.49$\times$ & [1.481, 1.503] & 0.0002 & 0.0004 & -42.91 & 0.0001 & 1.0 \\
go:dragonbox-vs-bd & api & verify & noncanonical/numeric-boundary & dragonbox & json-canon & dragonbox$^*$ & 1.49$\times$ & [1.483, 1.508] & 0.0002 & 0.0004 & -39.12 & 0.0001 & 1.4 \\
go:schubfach-vs-bd & api & verify & noncanonical/rfc-key-sorting & schubfach & json-canon & schubfach$^*$ & 1.02$\times$ & [1.017, 1.027] & 0.0002 & 0.0004 & -3.30 & 0.0000 & 0.4 \\
go:schubfach-vs-ryu & api & verify & noncanonical/rfc-key-sorting & schubfach & ryu & schubfach$^*$ & 1.02$\times$ & [1.011, 1.022] & 0.0002 & 0.0004 & -2.47 & 0.0000 & 0.4 \\
go:schubfach-vs-dragonbox & api & verify & noncanonical/rfc-key-sorting & schubfach & dragonbox & schubfach$^*$ & 1.02$\times$ & [1.015, 1.022] & 0.0002 & 0.0004 & -4.43 & 0.0000 & 0.4 \\
go:ryu-vs-dragonbox & api & verify & noncanonical/rfc-key-sorting & ryu & dragonbox & ryu & 1.00$\times$ & [0.997, 1.009] & 0.4681 & 0.5746 & -0.34 & 0.0000 & 1.1 \\
go:ryu-vs-bd & api & verify & noncanonical/rfc-key-sorting & ryu & json-canon & ryu & 1.00$\times$ & [0.998, 1.012] & 0.3341 & 0.4372 & -0.45 & 0.0000 & 1.1 \\
go:dragonbox-vs-bd & api & verify & noncanonical/rfc-key-sorting & dragonbox & json-canon & dragonbox & 1.00$\times$ & [0.997, 1.009] & 0.6753 & 0.7546 & -0.21 & 0.0000 & 0.6 \\
go:schubfach-vs-bd & api & verify & noncanonical/small & schubfach & json-canon & schubfach$^*$ & 1.35$\times$ & [1.344, 1.357] & 0.0002 & 0.0004 & -47.88 & 0.0000 & 0.8 \\
go:schubfach-vs-ryu & api & verify & noncanonical/small & schubfach & ryu & schubfach$^*$ & 1.02$\times$ & [1.012, 1.024] & 0.0002 & 0.0004 & -2.49 & 0.0000 & 0.8 \\
go:schubfach-vs-dragonbox & api & verify & noncanonical/small & schubfach & dragonbox & schubfach$^*$ & 1.02$\times$ & [1.014, 1.025] & 0.0002 & 0.0004 & -3.02 & 0.0000 & 0.8 \\
go:ryu-vs-dragonbox & api & verify & noncanonical/small & ryu & dragonbox & ryu & 1.00$\times$ & [0.996, 1.007] & 0.4573 & 0.5662 & -0.33 & 0.0000 & 0.9 \\
go:ryu-vs-bd & api & verify & noncanonical/small & ryu & json-canon & ryu$^*$ & 1.33$\times$ & [1.320, 1.334] & 0.0002 & 0.0004 & -43.72 & 0.0000 & 0.9 \\
go:dragonbox-vs-bd & api & verify & noncanonical/small & dragonbox & json-canon & dragonbox$^*$ & 1.32$\times$ & [1.318, 1.331] & 0.0002 & 0.0004 & -45.89 & 0.0000 & 0.8 \\
go:schubfach-vs-bd & api & verify & noncanonical/surrogate-pair & schubfach & json-canon & schubfach$^*$ & 1.01$\times$ & [1.004, 1.015] & 0.0066 & 0.0128 & -1.37 & 0.0000 & 0.8 \\
go:schubfach-vs-ryu & api & verify & noncanonical/surrogate-pair & schubfach & ryu & schubfach & 1.00$\times$ & [0.999, 1.010] & 0.1506 & 0.2311 & -0.69 & 0.0000 & 0.8 \\
go:schubfach-vs-dragonbox & api & verify & noncanonical/surrogate-pair & schubfach & dragonbox & schubfach & 1.01$\times$ & [1.001, 1.012] & 0.0478 & 0.0823 & -0.95 & 0.0000 & 0.8 \\
go:ryu-vs-dragonbox & api & verify & noncanonical/surrogate-pair & ryu & dragonbox & ryu & 1.00$\times$ & [0.996, 1.008] & 0.5705 & 0.6669 & -0.26 & 0.0000 & 0.9 \\
go:ryu-vs-bd & api & verify & noncanonical/surrogate-pair & ryu & json-canon & ryu & 1.00$\times$ & [0.999, 1.011] & 0.1634 & 0.2470 & -0.66 & 0.0000 & 0.9 \\
go:dragonbox-vs-bd & api & verify & noncanonical/surrogate-pair & dragonbox & json-canon & dragonbox & 1.00$\times$ & [0.997, 1.009] & 0.3995 & 0.5063 & -0.39 & 0.0000 & 0.9 \\
go:schubfach-vs-bd & api & verify & noncanonical/unicode & schubfach & json-canon & schubfach & 1.00$\times$ & [0.997, 1.006] & 0.5227 & 0.6222 & -0.30 & 0.0000 & 0.8 \\
go:schubfach-vs-ryu & api & verify & noncanonical/unicode & schubfach & ryu & schubfach$^*$ & 1.01$\times$ & [1.005, 1.014] & 0.0016 & 0.0033 & -1.78 & 0.0000 & 0.8 \\
go:schubfach-vs-dragonbox & api & verify & noncanonical/unicode & schubfach & dragonbox & schubfach$^*$ & 1.02$\times$ & [1.010, 1.020] & 0.0004 & 0.0008 & -2.43 & 0.0000 & 0.8 \\
go:ryu-vs-dragonbox & api & verify & noncanonical/unicode & ryu & dragonbox & ryu$^*$ & 1.01$\times$ & [1.001, 1.010] & 0.0260 & 0.0476 & -1.13 & 0.0000 & 0.5 \\
go:ryu-vs-bd & api & verify & noncanonical/unicode & ryu & json-canon & json-canon$^*$ & 1.01$\times$ & [1.004, 1.011] & 0.0004 & 0.0008 & 1.88 & 0.0000 & 0.5 \\
go:dragonbox-vs-bd & api & verify & noncanonical/unicode & dragonbox & json-canon & json-canon$^*$ & 1.01$\times$ & [1.009, 1.018] & 0.0002 & 0.0004 & 2.53 & 0.0000 & 0.5 \\
go:schubfach-vs-bd & api & verify & noncanonical/verify-whitespace & schubfach & json-canon & schubfach$^*$ & 1.72$\times$ & [1.706, 1.726] & 0.0002 & 0.0004 & -78.90 & 0.0000 & 1.0 \\
go:schubfach-vs-ryu & api & verify & noncanonical/verify-whitespace & schubfach & ryu & schubfach & 1.01$\times$ & [0.999, 1.011] & 0.1280 & 0.2000 & -0.72 & 0.0000 & 1.0 \\
go:schubfach-vs-dragonbox & api & verify & noncanonical/verify-whitespace & schubfach & dragonbox & schubfach$^*$ & 1.01$\times$ & [1.005, 1.018] & 0.0020 & 0.0041 & -1.60 & 0.0000 & 1.0 \\
go:ryu-vs-dragonbox & api & verify & noncanonical/verify-whitespace & ryu & dragonbox & ryu & 1.01$\times$ & [1.001, 1.012] & 0.0426 & 0.0742 & -1.00 & 0.0000 & 0.8 \\
go:ryu-vs-bd & api & verify & noncanonical/verify-whitespace & ryu & json-canon & ryu$^*$ & 1.71$\times$ & [1.699, 1.717] & 0.0002 & 0.0004 & -83.75 & 0.0000 & 0.8 \\
go:dragonbox-vs-bd & api & verify & noncanonical/verify-whitespace & dragonbox & json-canon & dragonbox$^*$ & 1.70$\times$ & [1.687, 1.706] & 0.0002 & 0.0004 & -80.14 & 0.0000 & 0.9 \\
\end{longtable}
\normalsize

\end{landscape}

\end{document}
